\chapter{GRB (O.P. Tandon): Complete Practice \ and  Objective Bank}
\label{chap:grb_all_questions}

\begin{problembox}
\textbf{Problem 1} \hfill \textbf{[GRB (O.P. Tandon) p.48 Q2]}
\label{q:ext-grb_all_questions-1}

[10;/textasciitilde();5r 
\underline{\hspace{1cm}} 0.77 + 0.0591 10 (95)3 x 10-5 
2 
g 
25x 2 
::::: 
(0.77 + 0.0226) 
0.7926 volt 
- ---------r.rnm$/cdot$Tt'!J>.;.' Tnri'/textasciitilde()'1'1i ,1 
.J un 
. th 
',/textasciitilde() , 
,'7 2+ 
equn num tS reacnea, fftlat-zs-
e-ratto-vJ <-vm,r;-ntratum OfLn--$/cdot$$/cdot$$/cdot$$/cdot$$/cdot$$/cdot$------.------- -----\underline{\hspace{1cm}} 
.... \underline{\hspace{1cm}}-
to Cu 2+ io/textasciitilde()ilibrium?\underline{\hspace{1cm}}'----------------$/cdot$-----jE/textasciitilde()$/circ$:::\underline{\hspace{1cm}}7:L\underline{\hspace{1cm}}:::::-="-'+\underline{\hspace{1cm}}$/cdot$.()..7\underline{\hspace{1cm}}W6-v.o1t\underline{\hspace{1cm}} 
.. u 
/textasciitilde(). /textasciitilde()':f/textasciitilde()' 
Hg/textasciitilde()+ IHg 
Solution: The reaction is, 
Zn +CuS04 /textasciitilde()Cu +ZnS04 
or 
Zn +Cu 2+ /textasciitilde()Cu +Zn2+ 
Eo.O.0591r [Zn.2:+-\underline{\hspace{1cm}}L 
::::: 
cell --2- og [Cu 2+] 
At equilibrium, E cell 
0 
EO 
0.0591 1, [Zn 2+] . 
cell ::::: -2- og [Cu 
] 
2+ 
2 
EO 
log [Zn 
]::::: 
X 
cell 
[Cu 2+] 
0.0591 
or 
(E/textasciitilde()II 
0.76+ 0.34 
1.lOvolt) 
::::: 2 X 1.10 ::::: 37.225 
0.0591 
(Zn 2+] 
1.679 X 1037 : 1 
[Cu 2+] 
Example 22. An excess' of liquid mercury is added to an 
acidified solution of 1.0 X 10-3 M Fe 3+. It is found that 5/% of 
Fe 3+ remains$/cdot$ at equilibrium at 25$/circ$ C Calculate EO 2+IH 
. 
Hgz 
g 
assuming that the only reaction that occurs is 
2Hg + 2Fe3+ /textasciitilde() 
Hg/textasciitilde()+ + 2Fe2+ 
(Given, EF 3+ 1 
2+::::: 0,77 volt). 
e 
Fe 
Solution: . 
2Hg 
At equilibrium, excess 
At equilibrium, Ecell 
+ 2Fe3+ /textasciitilde() 
Hg/textasciitilde()+ 
10-3 X 5. 10-3 x95 
100 
2x 100 
o 
(lIT 1995) 
+ 2Fe2+ 
10-3 x 95 
100 
c:;;/textasciitilde()/textasciitilde()E-xample 23. 
Prove that for two half reactions having 
potentials El and E2 which are combined to yield a third half 
reaction, having a potential E 3 , 
Solution: 
or 
or 
:::::,!I;/textasciitilde()l.+n'll!. . 
n3 
!!G3 
!!G1 + !!G2 
-n3FE3 ::::: - nl FE) - n2FE2 
n3E3 == n1 El + n2E2 
+ 
.:>Example 24. 
What is the standard potential of the Tl3+ ITI 
electrode? 
Tl3+ +2e- /textasciitilde()Tl+; EO:::::1.26volt 
Tl + + e - /textasciitilde() 
Tl; 
EO 
- 0.336volt 
Solution: 
T1 3+ + 2e- /textasciitilde() 
T1+; nFE o 
2x 1.26 x F::::: 2.52F 
T1+ + e- /textasciitilde() 
1'1; 
nFE o ::::: Ix (- 0.336) x F::::: - 0.336F . 
Adding 
2.52F - 0.336F 
2.184 
nF 
3 
::::: 0.728 volt 
Example 25. 
Calculate the minimum mass of NaOH 
required to be added in RHS to consume all the H + present in 
RHS of the cell of emf + 0.701 volt at 25$/circ$C before its use. Also 
. report the emf of the cell after addition of NaOH 
0.760 V 
Solution: 
The cell reaction is, 
Zn + 2H + /textasciitilde() 
Zn 2+ + H2
\end{problembox}

\begin{problembox}
\textbf{Problem 2} \hfill \textbf{[GRB (O.P. Tandon) p.51 Q1]}
\label{q:ext-grb_all_questions-2}

Electrochemistry is a branch ofpbysical chemistry 
which deals with the relationship between chemical energy and 
electrical energy and bow one can be converted into another. The 
subject is divided into two categories: (i) use of electrical energy 
to produce chemical changes (electrolysis ) and (ii) conversion of 
chemical energy iilto' electrical' energy (production of electricity 
by spontaneous redox reactions ).
\end{problembox}

\begin{problembox}
\textbf{Problem 3} \hfill \textbf{[GRB (O.P. Tandon) p.51 Q2]}
\label{q:ext-grb_all_questions-3}

Conductors are the substances whichallow'the passage 
of electric current through them. Those which do not allow the 
flow of electric$/cdot$ cwentthrough them are called insulators $/bullet$. 
Conductors are of two types:. 
.'. 
. 
(a) Metallic or electronic.cOJid!lttors .. are$/cdot$ those which 
remain. unchanged as cttrrent$/cdot$ fiowsthrough.them: These 
conductors transfer electric' curr:ept by transfer of electrons 
without transfer of matter.', Metals such . as 'CU, Ag, AI, Pt, etc., 
" 
.' 
-
. 
non- metals like carbon (graphite) and various aIloyshelongto 
this class.$/cdot$ 
. 
(b) Electrolytic conductors are the conductors like aqueous 
solutions of acids, bases and salts or substance in molten state. 
which allow the flow of electric current with chemical 
decomposition. These conductors are tennedas electrolytes. The 
substances whose aqueous solutions do not conduct electric 
current are called non-electrolytes. 
. 
To pass the current through an electrolytic conductor; two 
metallic rods or plates are required. These are termed electrodes. 
The electrode which is connected with the positive terminatof 
ba/textasciitilde()ery or through which electric' current enters the solution .is 
. termed anode and the electrode which is connected with negative 
terminal of battery or through which electric current leaves the 
solution is termed cathode. Actually the anode is the. electrode
\end{problembox}

\begin{problembox}
\textbf{Problem 4} \hfill \textbf{[GRB (O.P. Tandon) p.52 Q4]}
\label{q:ext-grb_all_questions-4}

Preferential discharge theory: 'If more than one type 
of coating an inferior metal with a superior metal by electrolysis. 
of ion is, attracted tbwardsaparticular electrode, then the ion
\end{problembox}

\begin{problembox}
\textbf{Problem 5} \hfill \textbf{[GRB (O.P. Tandon) p.52 Q6]}
\label{q:ext-grb_all_questions-5}

Arrhenius theory of electrolytic dissociation: It was 
discharged is the one which requires least energy. 
put forward by Arrhenius in 1884 to explain the properties of 
The, decreasing order' of the discharge potential or the 
electrolytic solutions. The main points of the theory are: 
increasmg ruder of deposition fOl the cations and anions is as 
(i) AD eleetfelyte, when dissolved in water, breaks up into 
foHows:' 
tWQ types Qf cllarged particles called iQns. The iQn carrying 
Cations': K + ,Na+, Ca2:", Mg2+, A13+, Zn2+, H +, Cu2+, Ag+, Au 3+ $/bullet$ 
positive charge is called cation and the ion carrying negative 
2 
' , 
'charge is called anion. 
Anions ,:S04-,NO;, OH",Cr, Br-, r. 
", 
" 
AB/textasciitilde()A+ (cation)+B- (anion) 
,S.FanidaY'slaws 6felectrolysis: These present the.. 
. . 
rellltionship,between the ,qmmtity of electric charge and the 
. (19 T/textasciitilde()e proces/textasciitilde() o/textasciitilde()spl/textasciitilde()ttmg of the m?lecules of the electrolyte 
'amount of the substance deposlted-artlie electrode.1hese were \underline{\hspace{1cm}},\underline{\hspace{1cm}}J11tO...lons.-1 and callecLu/textasciitilde()JlIsatio/textasciitilde().Jhe.../textasciitilde()acti/textasciitilde()D-of.thejotaJ number of 
given by Faraday in 1834. ' 
, 
" 
/textasciitilde()o/textasciitilde()ecu!es present m solut1On as 10ns IS known as degree of 
", (i) FirstJaw:,. ')Vhen an electric curre/textasciitilde()t is /textasciitilde()assed thr/textasciitilde()ugh an 
IOn:::tion/textasciitilde() 
, , 
electrolyte; 1hearnountof substancedeposlted Isproport1Onal to 
s de oted by(x . 
,the quantity ()f electric charge which flows' through the 
a = Number of molecules dissociated into ions 
electrQlyte.' 
Total number of molecules dissolved 
Matheniati/textasciitilde()ally, 
or 
WocQ 
W=ZQ=Zxlxt 
where, Zis a constant, known as electrochemical equivalent, I is ' 
the c/textasciitilde()nt in amperes and t is the time in seconds. When Q = 1 
coulomb or,oile ampere of current is passed for 1 seCbnd, W /textasciitilde() Z. 
'Zisthtrs-themassofthesubstance deposited 'by one" coulomb: 
(ii)Second' law: Wheri $/bullet$ the same quantity of current, is 
passed, through ,different electrolytes, the masse/textasciitilde() . of different 
substances deposited on the electrodes will be in the ratio of their 
equivalent masses. F.orexample,ifsame quantity of current is ' 
passeq through ,copper sulphate and silverttitrate solutions, then 
'M;/textasciitilde()ssofcoppei deposited Equi/textasciitilde()areritm/textasciitilde()s /textasciitilde()f copper 
'Massof silver deposited 
'Equiyalentmass of silver 
, ' One grrurtequivalent of /textasciitilde()y substance is deposited by passing 
same aniOUllt Of charge, i.e., 96500 coulomb (1 faraday). 
' 
. E:::dZ x 96500 
, E ' 
or:,. 
Z;::-'--
'96500 
Th 
h 
' db$/cdot$ 
nF, 
, e C arge carne 
y anIon = 
, 
coulomb 
, 
, 
6.02 X 1023 
When; n ='1,' 
F' 
96500 
The fundamental unit of charge = 
' 
;:: --/textasciitilde()-
, 
6.02 x 1023 
6.02 x 1023 
The degree of ionisation depends upon: (a) Nature of solute 
(b) Nature of solvent (c) Dilution and temperature. 
'. (iii) Ions present in solution constantly reunite to form neutral 
molecules; thus, there is a state of dynamic equilibrium between 
ionised and unionised molecules. 
AB/textasciitilde()A+ +B-
'[A+l[B-] 
., 
, 
, 
' 
;:: K (ionisation constant) 
[AB] 
, 
(iv) The ions are discharged always in equivalent aniounts 
when current is passed through electrolytic solution no matter 
what their relative speeds are. 
(v) The electrolytic solution is always neutral in nature as the 
total' charge on cations is equal to the total charge on anions. 
However, it is not necessary that, the number of cations and 
anions should be always equal. 
(vi) The properties of electrolytes in solution an: , the 
properties of ions present in the solution . 
(vii) The ions act as molecules for colligative properties. 
(viii) The conductivity of the solution depends onthe nature 
and number of ions. 
' 
' 
Evidences in favour of Arrhenius theory: (i) Ohm's law 
applicability: No part of the currerit is used in splitting up the 
molecules into ions. (ii) X-ray studies have shown that ions are 
present in solid electrolytes. The ionic compounds in molten state 
behave as good conductors. (iii) Ionic reactions are observed in 
the case of electrolytes. (iv) Constant value of heat of
\end{problembox}

\begin{problembox}
\textbf{Problem 6} \hfill \textbf{[GRB (O.P. Tandon) p.53 Q9]}
\label{q:ext-grb_all_questions-6}

Galvanic or Voltaic or Electrochemical cell: It is the
\end{problembox}

\begin{problembox}
\textbf{Problem 7} \hfill \textbf{[GRB (O.P. Tandon) p.53 Q7]}
\label{q:ext-grb_all_questions-7}

Conductivity is the property of the conductor which 
device in which chemical e/textasciitilde()ergy is converted into electrical 
facilitates the flow of electricity through it. It is equal to the 
energy. In this cell a redox reaction is carried out in an indirect 
reciprocal of resistance. 
manner and the decrease in free energy during the chemical 
C d' . 
I 
1 
process appears as electrical energy. 
. 
on uctlvlty = Resistance = R 
The common galvanic cell, dry cell imd the lead' storage 
It is expressed as $/Omega$ -lor mho. 
battery a,re the devices for converting chemical energy into 
electrical energy. Daniell cell is the typical example of galvanic 
Specific conductivity is defined as the reciprocal of specific 
resistance; i.e., it is the conductance of I cm3 of a conductor. It is 
cell. It consists of two half-cells, one containing zinc electrode 
dipping in solution of zinc sulphate (1 M) and another containing 
represented by the symbol K. 
copper electrode dipping in solution of copper sulphate (I M). 
Ki 
I x conductivity 
Zinc electrode acts as anode. Atthis electrode oxidation occurs. 
p 
a 
Zn :-7 Zn 2+ + 2e -. COReere1ectrode acts as cathode: At this 
h
'.:l$/cdot$ 
1... 
1 
d 
.:I 
. /%f'flfr-f.t:/textasciitilde()........!e;!!IQecldtJJroJ.!.dJl;eL.llr.e;;J.dlJ.u/textasciitilde()cclJtillOD1LlOJ1Ca;c.clJ
llll
TS;", .../C...Jlu..1 2.-'-+-:t:+/textasciitilde()2:.ee:....-....:=7.=$/cdot$ 
/textasciitilde()CI...llU\underline{\hspace{1cm}}. .../Wh'/Lue/textasciitilde()n:l-L.tbwe/textasciitilde()$/cdot$+t/w"uo--
-----Were, 
t 
,*lstaRce uetween 
 I ectr04lS anu 
tl ..... -area 0 
cross-section of electrode. 
half-cells are connected by a salt bridge and the metal electrodes' 
In the case of electrolytic solutions, specific conductance is 
are joined externally, the electrons from zinc electrode. (-ve 
the conductance of a solution of definite dilution enclosed in a 
electrode) move towards copper electrode (+ve electrode): i.e., 
cell having two electrodes, of unit. area separated by one 
current flows from cathode to anode. The cell reaction is: 
centimetre. The units ofspe.cifi<;;. con.ductance are .$/Omega$ -!. cm -I. 
Zn-+- /textasciitilde()l'tn.::.2+:!:.. ---li'<-ICPlu't-----=-----'----
The ratio of II a is called cell constant. 
Equivalent conductance: Conductivity of a solution 
centaining 19 equivalent of the electrolyte. It is denoted by 'N. If 
the concentration ef solution is C g equivalent per litre/textasciitilde() then 
A 
K 1000 
C 
or 
A=i(xV 
where, V is the volume in mL containing Ig equivalent of the 
electrolyte. The unit of equivalent conductance is $/Omega$'-I cm2 
equiv -'-I. 
Molar conductance: Conductance of a solution containing 
I g mole of an electrolyte. If V is volume in mL centaining 19 
mole, then 
-
Molar conductance, )1 /textasciitilde() K X V 
If C is the concentration ofthe solution in g mele per litre, then 
1000 
)1=KX 
C 
. Its unit is $/Omega$ -[ cm2 mol-I. 
. 
Molar conductance 
EqUIvalent conductance = 
. 
n 
Molecular 
where, 
n 
Equivalentmass
\end{problembox}

\begin{problembox}
\textbf{Problem 8} \hfill \textbf{[GRB (O.P. Tandon) p.53 Q8]}
\label{q:ext-grb_all_questions-8}

Kohlrausch's law: At infinite dilution, each ion makes 
a definite contribution tewards equivalent conductance of the 
electrolyte. The value of equivalent conductance at infinite 
dilution for any electrolyte is the sum ef conductances at infinite 
dilution of its constituent iens. Thus, 
A;' =Aa + Ac 
Aa and Ac are called the ionic conductances of anion and 
cation at infinite dilution respectively. It clm be used to determine 
degree of dissociation. 
S/textasciitilde()lt-bridge allows the flow of current by completing the 
circuit and lTIaintains . electrical neutrality .. It also prevents 
liquid-HqtiidjunctionpoiintiaJ.-
-" --
A galvanic cell is represented in the following manner: 
LHS 
Anode (-ve) 
Znl Zn2+ (1 M) 
Oxidation half-cell 
Zn /textasciitilde() Zn2+ + 2e-
Salt-bridge 
RHS 
Cathode (+ve) 
Cll/textasciitilde()+ (1 M) I Cll 
Reduction half-cell 
'. ell 2+ + 2e- = Cll 
,10.. Electrftdepotential/textasciitilde()- When--ametal--is--placoo--in-'a"-
solution of its ions, the metal either acquires a positive charge or 
negative charge with respect to solutien. On account of this a 
definite potential is developed betwe.en. the metal and the 
solutien. This potential difference is termed electrode potential. 
The magnitude of petential depends on the nature of el(;lctrode, 
concentration of ions and temperature. It is a ,measure ef the 
tendency of the metal to lose or gain electrons or a measure of the 
relative tendency to undergo. oxidation or reductio.n. Depending 
cn. the nature of the metal electrode, the electrode potential is of 
two types: 
(i) Oxidation potential: When electrode is negatively 
charged with respect to solution, i,e., oxidation occurs and the' 
electrode acts as anode 
M :-7M n+ + ne-
(ii) Reduction potential: When electrode is positively 
charged with respect to solution, i.e., reductien occurs and the, 
electrode acts as cathode, 
M"+ + ne- :-7 M
\end{problembox}

\begin{problembox}
\textbf{Problem 9} \hfill \textbf{[GRB (O.P. Tandon) p.54 Q12]}
\label{q:ext-grb_all_questions-9}

Emf of a cell: The difference in potentials of two 
half-cells is known as the emf of the celL 
E/textasciitilde()ell = Red. pot. cathode Red. pot. anode (E/textasciitilde()ght 
E/textasciitilde()ft ) 
= Red. pot. cathode + Oxid. pot. anode 
Oxid. potential anode' Oxid. potential cathode
\end{problembox}

\begin{problembox}
\textbf{Problem 10} \hfill \textbf{[GRB (O.P. Tandon) p.54 Q13]}
\label{q:ext-grb_all_questions-10}

Some more reference electrodes: Since, a hydrogen 
electrode is difficult to prepare and maintain, it is usually 
replaced by other reference electrodes which are known as 
secondary reference electrodes. These are convenient to handle 
and prepared easily. Calomel electrode and silver-silver chloride 
are used as reference electrodes. The potentials of these 
electrodes are fixed on hydrogen scale but their values depend on 
strong electrolyte concentration.
\end{problembox}

\begin{problembox}
\textbf{Problem 11} \hfill \textbf{[GRB (O.P. Tandon) p.54 Q14]}
\label{q:ext-grb_all_questions-11}

Reversible and irreversible cells: A cell is said to be 
reversible if the following two conditions are fulfilled: 
(i) The chemical reaction of the cell stops when an exactly 
equal external emf is applied. 
(ii) The chemical reaction of the' cell is reversed and the 
current flows in opposite direction when the external emf is 
slightly higher than that of the celL 
' 
Any other cell which does not obey the above two conditions 
is termed as irreversible.' Daniell cell is reversible but 
Zn I H2804 .1 Ag cell is irreversible in nature. 
. 
.
\end{problembox}

\begin{problembox}
\textbf{Problem 12} \hfill \textbf{[GRB (O.P. Tandon) p.54 Q15]}
\label{q:ext-grb_all_questions-12}

Prediction for occurrence of a redox reaction: Any 
redox reaction would occur spontaneously if the free energy 
change (LiG) is negative. 
LiGo 
nFEo where, 'n' is the number of electrons involved, 
FisthevalueoffaradayandE" is the cellemfLlGo can be 
negative if EO is positive. 
When EO is positive, the cell reaction will be spontaneous and 
serves as a SOQrGe of electfical MeFg  ..$/bullet$ 
'. \underline{\hspace{1cm}}l/textasciitilde() ... Nel;'usLeqltatien: The pottmtial of aa electrode 
changes with the change in concentration of ions in solution in 
. contact with the electrode. Increase in concentration of cations 
results in an increase of reduction potential of an electrode. 
or 
Consider 
M"+ +ne-4M 
pnEo ' /textasciitilde()7;3rj3Rl$/cdot$10-$/cdot$- (reducedlNni] .: 
E red 
red 
. 
nF 
g [oxidised.forrn] 
2.303RT$/cdot$ 0.0591 
F 
and the concentration of metal is unity. 
'Eo 
0:0591 1-' ['M 11+] 
red + 
og 
n 
Similarly, consider the electrode where oxidation occurs. 
M -'7M n+ + ne-
E 
EO 
0.0591 log [M 11+ ] 
oxid. 
oxid. 
n 
Cell potential depends on the potential of anode. and cathode.\underline{\hspace{1cm}} 
ECell 
EAnode + ECathode 
Oxid. pot. of anode /textasciitilde()Red. pot. of cathode 
E o 
0.05911 
[A 
d" 
. 
]' 
= oxid; 
og 
no IC Ion cone. 
n 
E o 
0.05911 
[C hod' '. 
] 
+ 
red - --. og 
at 
Ie IOn conc. 
n 
$/circ$ . 
$/circ$ 
0.0591 
[Anodic ion conc. ] 
= Eoxid, + Ered -
log 
. . 
$/circ$ 
Eeen 
n 
[Cathodic IOn conc.] 
0.0591 
[Products] 
log----
n 
[Reactants] 
E o 
0.05911 
cell - -- og ---":':':::" 
n 
[ion]RHS
\end{problembox}

\begin{problembox}
\textbf{Problem 13} \hfill \textbf{[GRB (O.P. Tandon) p.54 Q17]}
\label{q:ext-grb_all_questions-13}

Electrochemical series: When the electrodes (metals 
and non-metals) in contact with their ions are arranged on the 
basis of values of their standard reduction potentials or standard 
oxidation potentials, the resulting series 
is 
called the 
electrochemical or activity series of the elements. 
I
\end{problembox}

\begin{problembox}
\textbf{Problem 14} \hfill \textbf{[GRB (O.P. Tandon) p.55 Q18]}
\label{q:ext-grb_all_questions-14}

Relation between emfand equilibrium constant: 
AGO = - 2.303RT log K 
AGo= nFE:eIl 
So, 
nFE:eIl = 2.303RTlogK 
EO 
- 2.303RT 1 K 
cell -
nF 
og 
0.0591 log K(T= 298 K) 
n' 
,
\end{problembox}

\begin{problembox}
\textbf{Problem 15} \hfill \textbf{[GRB (O.P. Tandon) p.55 Q19]}
\label{q:ext-grb_all_questions-15}

Primary cells are those whichl;lave definite . life and 
become dead ovtlr a titne. These cannot be recharged; Examples 
are dry cell, mercury ce)l, etc. , .' 
' ' 
" , 
' 
Secondary cells 'are the cells which can be recharged by . 
passing direct current through them. Lead storage battery and 
nickel-cadmium cell, are examples of secondary cells. Lead 
storage battery consists of spongy lead as anode, a grid of lead 
packed with Pb02 ,as cathode and an aqueous /textasciitilde()olution of H2S04 
(38/% by mass or 20/% H2S04 of speciftc 'gravity 2.15) as 
electrol e.' 
. 
Pb +Pb02 + 4H + + 2S0/textasciitilde()- -7 
2PbS04 +2H:!O 
Voltage 2.0 volt. The concentration 
of H2S04 decreases during 
. 
discharging. 
2PbS04 + 2H20 -7 Pb + Pb02 
4H+ + 2S0/textasciitilde()-' 
. Nameof/textasciitilde() 
Anooe 
CathOde 
Electrolyte 
uSed 
Dry cell 
Zn 
Graphite 
Mn02 + C 
(around 
cathode) 
NH4CI + ZnCl2 ' 
(around anode) 
Mercury cell 
Zn 
Graphite 
HgO+ KOH 
(moist) 
Lead 
storage 
Pb 
Pb02 
H2SO4 
(38/% by 
battery 
mass) 
Lithimn battery 
Li 
Metal sulphide LiCI04 in organic 
solvent 
Ni-Cadcell 
Cd 
NiO 
KOH solution 
Fuel cell 
Porous carbon Porous carbon 
Cone. 
NaOH 
or 
(H2-0 2) 
(with catalyst) (with catalyst) 
KOH (aq.).
\end{problembox}

\begin{problembox}
\textbf{Problem 16} \hfill \textbf{[GRB (O.P. Tandon) p.55 Q20]}
\label{q:ext-grb_all_questions-16}

Fuel cell: 
The cells which convert chemical energy of 
fuels directly mto electrical energy are known as fuel cells. One 
of sucll cells is mel cell (H2 -02) in which electficalenergy IS 
generated by the use ofH2 and O2, In fuel cell NaOH is used as 
an electrolyte and hydrogen gas is diffused at the anode made of 
porous carbon. At the cathode oxygen is diffused through a 
porous carbon. The following chemical reactions occur: 
2H+ 20H- --t 2H20+2e-
(ii) At cathode, O2 + 2H20 + 4e - --t 40H-
Net cell reaction 
2Hz + O2 --t 2H2 0 + energy 
Fuel cells .are quite efficient, and free from pollution. These 
are used in spacecrafts.
\end{problembox}

\begin{problembox}
\textbf{Problem 17} \hfill \textbf{[GRB (O.P. Tandon) p.55 Q21]}
\label{q:ext-grb_all_questions-17}

Concentration cells: If two plates of the sallie metal 
are dipped separately into two solutions of the same electrolyte 
and are cOI11":ected with a salt bridge, the whole arrangement is 
found to acts as galvanic cell. Such ce)ls are known as 
concentration cells. These are of two types: '. 
(i) Electrode concentration cells 
,. 
,. , 
.Pt, H2 -<pressurePtJI H +IH2 (pressu.rep;),.Pt 
". 
. 
. 
' 
. 
. If PI >P2 oxidation occUrs at LHS /textasciitilde()lecirode and reduction 
occurs at RHS electrode; 
. 
E 
=$/circ$.0591 10 (P1l at 2SOC 
' .$/bullet$ 
cell 
2 
g ( 
) ". 
, 
P2.. . 
(ii) Electrolyte concentration cells 
M I M n+ (CdIlMn+ (C2 ) I,M 
ZnlZnn+ (CdIlZnH+ (C2 )IZn 
C2 is higher than C1. . 
E 
- 0.0591 1 
C2 t 25$/circ$C 
cell --- og-a 
n 
C1
\end{problembox}

\begin{problembox}
\textbf{Problem 18} \hfill \textbf{[GRB (O.P. Tandon) p.56 Q1]}
\label{q:ext-grb_all_questions-18}

Match the List-I with List-II and List-III: 
List-I 
List-II 
List-ID 
(Quantity) 
(Symbol) 
(Unit) 
(a) Conductivity 
(P)Am 
(u) mhocm:-I 
(b) Cell constant . 
(q) K 
(v) cm-1 
(c) Molar conductance (r) Ae 
(w) $/Omega$ -,-1 cm2 mol-l 
(d) Equivalent 
(s) 1I A 
(x) $/Omega$-I cm2 eq-l 
conductance
\end{problembox}

\begin{problembox}
\textbf{Problem 19} \hfill \textbf{[GRB (O.P. Tandon) p.56 Q2]}
\label{q:ext-grb_all_questions-19}

Match the salts in the Column-I with their Use in Column-ll: 
 Glumo-I 
(b) Agar-agar 
(c) 0.1 N KCl 
(d) Quinhydrone 
Column II 
(q) Calomel.electrQde 
(r) Used in ice cream 
(s) Redox electrode 
$/cdot$3. MatclnneColuriiiFlwitIiColumn::I1:$/cdot$ 
Column-I 
Column-II 
(Relation) 
(Term) 
(a) (l (Degree of ionisation) 
(P ).,,+ I A 
(b) t+ (Transport number) 
(c) Fraction ofa mole 
undergoing ,ionization 
(d) NaCl 
(q) NmIA/textasciitilde() 
(r) U+ IU+ + U-
(s) N lAo 
e 
e 
(t) Am = Ae
\end{problembox}

\begin{problembox}
\textbf{Problem 20} \hfill \textbf{[GRB (O.P. Tandon) p.56 Q4]}
\label{q:ext-grb_all_questions-20}

Match the physical quantities in the List-I with their units in 
List-II: 
List'-I 
List-II 
(a) Resistance 
(P) $/Omega$ 
(b) Resistivity 
(q) volt amp-I 
(c) Conductivity 
(r) $/Omega$m 
(d) Specific conductance 
(s) $/Omega$-I m-I
\end{problembox}

\begin{problembox}
\textbf{Problem 21} \hfill \textbf{[GRB (O.P. Tandon) p.56 Q5]}
\label{q:ext-grb_all_questions-21}

Match the Cobnnn-I with Column-ll: 
Column-I . 
Column-II 
(a) Specific conductance, K 
(P) A';" I A"", 
(b) Molar conductance, Am 
(q) Decreases with dilution 
(c) Resistance of electrolyte (r) Increases with dilution 
solutionR 
(d) Degree of ionization of 
weak: electrolyte, (l 
(s) Increases with increase in 
the distance between 
parallel plates
\end{problembox}

\begin{problembox}
\textbf{Problem 22} \hfill \textbf{[GRB (O.P. Tandon) p.56 Q6]}
\label{q:ext-grb_all_questions-22}

Match Column-I with Column-II: 
Column-I 
(Combination ofhalf-ceU reactions) 
Column-II 
(potential of 
overaU 
process) 
(a) 60Er + Br-/textasciitilde() BrOi + 3H20+6e-
(P) EO= 0.56 V 
EO =-0.61 V 
20Er + Br-/textasciitilde() BrO- + H20 + 2e-
 0 -0.76 V 
+4e-/textasciitilde() S+ 3H20 
EO =0.45 V 
SO/textasciitilde()- + 4W + 2e-/textasciitilde() H2S03 + H20 
.  O=0.17V 
EO=? 
(c) CIOi +6W +6e-""'-: cr + 3H20 
 0 
+1.45 V 
CI2+2e-/textasciitilde() 2Cr 
EO == +1.36 V 
 0 ? 
(q) E"= - 0.535V 
(r)  0= +0.36 V 
(d) Znta:V -+ 2e--7 Zn(s)  0 = -0.76 V (s) EO =+1.47 V 
Ag+ + e- --7 Ag(s)  0 = +0.80 V 
Zn(s) + 2Ag+/textasciitilde() Zn2+ (aq)+ 2Ag(s) 
 O=?
\end{problembox}

\begin{problembox}
\textbf{Problem 23} \hfill \textbf{[GRB (O.P. Tandon) p.56 Q7]}
\label{q:ext-grb_all_questions-23}

Match the Colqrnn-I with Column-ll: 
Column-I 
. (a) Concentration cell 
(b) Edison cell 
(c) Mercury cell 
(d) Dry cell 
Column-II 
(P) Fe is oxidised by Ni20 3 
(q) Zinc anode 
(r) HgO cathode 
(s)  0 =0
\end{problembox}

\begin{problembox}
\textbf{Problem 24} \hfill \textbf{[GRB (O.P. Tandon) p.56 Q8]}
\label{q:ext-grb_all_questions-24}

Match the Column-I with Column-II: 
Column-I (Electrode) 
(a) Calomel 
(b) Glass 
(c) Hydrogen 
(d) Quinhydrone 
Column-II (Type) 
(P) Reference 
(q) Redox 
(r) Membrane 
(s) Gas 
I
\end{problembox}

\begin{problembox}
\textbf{Problem 25} \hfill \textbf{[GRB (O.P. Tandon) p.57 Q9]}
\label{q:ext-grb_all_questions-25}

Match the Column-I with Column-II: 
Column-I 
(a) Kohlrausch law 
(b) Am 
(c) K (kappa) 
(d) Ci 
Column-II 
(P) Am / A/textasciitilde()II 
(q) J:. x 1 
R A 
(r) A:IICa3 (P04 h 
= 3).,$/circ$CaH + 2)"oPO!-
(s) K x 1000 
M
\end{problembox}

\begin{problembox}
\textbf{Problem 26} \hfill \textbf{[GRB (O.P. Tandon) p.57 Q10]}
\label{q:ext-grb_all_questions-26}

Match the Column-I with Column-II: 
Column-I 
(a) EO 
0 
(b) E= 0 
(c) LlG=O 
(d) 1 Faraday 
Column-ll 
(p) Cell is discharged 
(q) Q K 
(r) 96500 coulomb 
(8) 1 mol electrons 
(t) Concentration cell 
[,4/textasciitilde(). -----------------1
\end{problembox}

\begin{problembox}
\textbf{Problem 27} \hfill \textbf{[GRB (O.P. Tandon) p.57 Q1]}
\label{q:ext-grb_all_questions-27}

(a-q, u), (b--s, v), (c-p, w), (d/textasciitilde()r, x)
\end{problembox}

\begin{problembox}
\textbf{Problem 28} \hfill \textbf{[GRB (O.P. Tandon) p.57 Q2]}
\label{q:ext-grb_all_questions-28}

(a-q), (b--p, r), (c-p), (d-s)
\end{problembox}

\begin{problembox}
\textbf{Problem 29} \hfill \textbf{[GRB (O.P. Tandon) p.57 Q3]}
\label{q:ext-grb_all_questions-29}

(a-q, s), (b/textasciitilde()p, r), (c-q, s), (d-t) 
.4. 
(a-p, q), (b--r), (c/textasciitilde()s), (d-s)
\end{problembox}

\begin{problembox}
\textbf{Problem 30} \hfill \textbf{[GRB (O.P. Tandon) p.57 Q5]}
\label{q:ext-grb_all_questions-30}

(a-q), (b--r), (c-q, s), Cd-p, r)
\end{problembox}

\begin{problembox}
\textbf{Problem 31} \hfill \textbf{[GRB (O.P. Tandon) p.57 Q6]}
\label{q:ext-grb_all_questions-31}

(a-q), (b--$/cdot$r), (c-s), (d-p)
\end{problembox}

\begin{problembox}
\textbf{Problem 32} \hfill \textbf{[GRB (O.P. Tandon) p.57 Q7]}
\label{q:ext-grb_all_questions-32}

(a--s), (b--p), (c-q, r), (d-$/cdot$.q)
\end{problembox}

\begin{problembox}
\textbf{Problem 33} \hfill \textbf{[GRB (O.P. Tandon) p.57 Q8]}
\label{q:ext-grb_all_questions-33}

(a/textasciitilde()p), (b-r), (c-s), (d-q)
\end{problembox}

\begin{problembox}
\textbf{Problem 34} \hfill \textbf{[GRB (O.P. Tandon) p.57 Q9]}
\label{q:ext-grb_all_questions-34}

(a-r), (b-s), (c-q), (d-p)
\end{problembox}

\begin{problembox}
\textbf{Problem 35} \hfill \textbf{[GRB (O.P. Tandon) p.57 Q10]}
\label{q:ext-grb_all_questions-35}

(a-t), (b--p, q), (c-p, q), (d-r, s)
\end{problembox}

\begin{problembox}
\textbf{Problem 36} \hfill \textbf{[GRB (O.P. Tandon) p.58 Q1]}
\label{q:ext-grb_all_questions-36}

How many coulombs are required for the following 
reductions? 
(i) I mole of Ag + ions to Ag 
(ii) I mole of Cu 2+ ions to Cu 
(iii) 1 mole ofMn04 ions to Mn/textasciitilde()-
[Ans. 
(i) 96500 coulomb (ii) 193000 coulomb (iii) 96500 
coulomb] 
.
\end{problembox}

\begin{problembox}
\textbf{Problem 37} \hfill \textbf{[GRB (O.P. Tandon) p.58 Q2]}
\label{q:ext-grb_all_questions-37}

How many faradays are needed to reduce 2 gram-mole of 
Cu 2+ to Cu metal? 
. 
[Ans. 
4 faraday]
\end{problembox}

\begin{problembox}
\textbf{Problem 38} \hfill \textbf{[GRB (O.P. Tandon) p.58 Q3]}
\label{q:ext-grb_all_questions-38}

How many faradays are released when 12.7 g of copper metal 
is changed into copper ions? 
. 
n. A lead storage battery is discharged in which the following 
electrode reaction takes place: 
Pb(s) + PbOis) + 2HzS04(l) /textasciitilde() 
2PbS04 + 2H20 
The original volume in electrolyte is 1 litre. During discharge 
the concentration of H2S04 changes from 34.6/% by weight 
(density 1.261 g/mL) to 27/% by weight. How many faradays . 
have left the' anodic half-cell? Water produced during 
electrolysis is used up. 
. 
[Ans. 
1.254 faraday]
\end{problembox}

\begin{problembox}
\textbf{Problem 39} \hfill \textbf{[GRB (O.P. Tandon) p.58 Q12]}
\label{q:ext-grb_all_questions-39}

The chemical reaction given below: 
C12(g) + S02(g) + 2H20(l) /textasciitilde() 
2Cqaq.) + 3H+(aq.) 
-----------1[/textasciitilde()H/textasciitilde()inwt/textasciitilde()$/cdot$/textasciitilde()/textasciitilde()C/textasciitilde()H/textasciitilde()--/textasciitilde())/textasciitilde()C/textasciitilde()U/textasciitilde()2+/textasciitilde()/textasciitilde()I/textasciitilde()2/textasciitilde()e/textasciitilde()-/textasciitilde()]----------/textasciitilde()--------------pp/textasciitilde()ro/textasciitilde()c/textasciitilde()eeed/textasciitilde()s/textasciitilde()r/textasciitilde()e/textasciitilde()addUih '/textasciitilde()in/textasciitilde()aaG)/textasciitilde()/textasciitilde()e/textasciitilde()O/textasciitilde()Y/textasciitilde()8/textasciitilde()/textasciitilde()'/textasciitilde()io/textasciitilde()nh.----------/textasciitilde()--------
1 mole 
2 mole 
J 
'1 
$/bullet$$/bullet$ ' 
63.5 
2F 
(i) Give the hair-cell reactions. 
.
\end{problembox}

\begin{problembox}
\textbf{Problem 40} \hfill \textbf{[GRB (O.P. Tandon) p.58 Q4]}
\label{q:ext-grb_all_questions-40}

Calculate the number of coulombs required to deposit 40.5 g 
(ii) Ifa fully charged cell initially held I mole ofCl2, for how 
Al when the electrode reaction is: 
many days could it sustain a current of 0.05 ampere, 
Al3+ + 3e- ---7 Al 
assuming the cell becomes non-operative when 90/% of 
[Ans.,-4T34-x.)05,eoulomb1
, 
"'''''' 
initial CI 2 has been used up? 
'J 
'(Ans.'40.z"cfitysr---- ",,-,---- --.------.----.
\end{problembox}

\begin{problembox}
\textbf{Problem 41} \hfill \textbf{[GRB (O.P. Tandon) p.58 Q5]}
\label{q:ext-grb_all_questions-41}

Calculate the number of coulombs required to deposit 50 g of 
silver at cathode from silver nitrate solution.
\end{problembox}

\begin{problembox}
\textbf{Problem 42} \hfill \textbf{[GRB (O.P. Tandon) p.58 Q13]}
\label{q:ext-grb_all_questions-42}

A current of 80 microampere is passed through a solution of 
(Atomic mass of silver = 108) 
AgN03 for 32 minutes using platinum electrodes. A uniform 
[Ans. 
44675.9 coulomb] 
single atom thick layer is deposited covering 86/% of the
\end{problembox}

\begin{problembox}
\textbf{Problem 43} \hfill \textbf{[GRB (O.P. Tandon) p.58 Q6]}
\label{q:ext-grb_all_questions-43}

How much electric charge is required to produce 20.0 g of 
cathode surface. If total surface area of cathode is 601.7 cm
2
, 
calcium from molten CaCI
2? ' 
calculate the area covered by one Ag atom. 
[Ans. 
96500 coulomb] 
[Ans. 
5.4 x 10-
16 cm
2
]
\end{problembox}

\begin{problembox}
\textbf{Problem 44} \hfill \textbf{[GRB (O.P. Tandon) p.58 Q7]}
\label{q:ext-grb_all_questions-44}

If 3 faradays of electricity are passed through an iron (II)
\end{problembox}

\begin{problembox}
\textbf{Problem 45} \hfill \textbf{[GRB (O.P. Tandon) p.58 Q14]}
\label{q:ext-grb_all_questions-45}

Anthracene can be electrolytically oxidised to anthraquinone. 
bromide solution, how many grams Of iron metal will be 
$/circ$ 
deposited? (At. mass of iron = 56) 
II 
[Ans. 
84 g]
\end{problembox}

\begin{problembox}
\textbf{Problem 46} \hfill \textbf{[GRB (O.P. Tandon) p.58 Q8]}
\label{q:ext-grb_all_questions-46}

A certain quantity of electricity deposits 0.54 g of Ag from 
silver nitrate solution. What volume of hydrogen will be 
liberated by the same quantity of electricity at 27$/circ$C and 
750 mm ofHg pressure? 
Mass of 
[Hint: 
Mass of Ag 
1 x 0.54 
Mass of H2 
0.005 g 
108 
/textasciitilde()olumeof H2 
0.005 x 0.0821 x 300 x 760 
2 x750 
0.06234 litre] 
(V =/textasciitilde(). RT) 
M 
P
\end{problembox}

\begin{problembox}
\textbf{Problem 47} \hfill \textbf{[GRB (O.P. Tandon) p.58 Q9]}
\label{q:ext-grb_all_questions-47}

In the electrolysis of an aqueous cupric bromide solution, how 
many grams of bromine are fmmed on passing a current of 1 
ampcrc for 16 minutes and 5 seconds? Write the anode and the 
cathode half reactions. 
[Ans. 
0.80 g; cathode reaction Cu 2+ + 2e- ---7 Cu; 
anode reaction 2Br- ---7 Br2 + 2e-]
\end{problembox}

\begin{problembox}
\textbf{Problem 48} \hfill \textbf{[GRB (O.P. Tandon) p.58 Q10]}
\label{q:ext-grb_all_questions-48}

A current of 1.5 amperes is passed through a$/cdot$solution of a salt 
of a bivalent metal for 30 minutes. Increase in mass of cathode 
is 0:8898 g. Find the atomic mass of the metal. 
[f/ 115. 
63,6] 
3[0] 
C14H1o /textasciitilde() 
Anthracene 
II $/circ$ 
Anthraquinone 
What mass of anthraquinone can be produced by the passage 
of 1 ampere current for an hour at 100/% efficiency? 
[Ans. 
1.2932 g]
\end{problembox}

\begin{problembox}
\textbf{Problem 49} \hfill \textbf{[GRB (O.P. Tandon) p.58 Q15]}
\label{q:ext-grb_all_questions-49}

Dal lake has water 8.2 x 1012 litre approximately. A power 
reactor produces electricity at the rate of 1.5 x 106 coulomb 
per second at an appropriate voltage. How much time would it 
take to electro lyse the lake? 
[Ans. 
1.9 million year]
\end{problembox}

\begin{problembox}
\textbf{Problem 50} \hfill \textbf{[GRB (O.P. Tandon) p.58 Q16]}
\label{q:ext-grb_all_questions-50}

How many grams of H2 and 02 are produced during the 
electrolysis of water by a 1.30 ampere current for 5 hours? 
What volumes of dry gases are produced at NTP? 
[Ans. 
0.242 g Hz; 1.94 g 02; 2.72 litre H2; L36litre 0, ]
\end{problembox}

\begin{problembox}
\textbf{Problem 51} \hfill \textbf{[GRB (O.P. Tandon) p.59 Q17]}
\label{q:ext-grb_all_questions-51}

How many coulombs of electricity would be required to 
reduce the iron in 36.0 g of potassium hexacyano ferrate(III), 
K3Fe(CN)6' to metallic iron?
\end{problembox}

\begin{problembox}
\textbf{Problem 52} \hfill \textbf{[GRB (O.P. Tandon) p.59 Q26]}
\label{q:ext-grb_all_questions-52}

Three electrolytic cells A ,B and C containing electrolytes zinc 
sulphate, silver nitrate and copper sulphate respectively, were 
connected in series. A steady current of 1.50 ampere was 
passed through them until 1.45 g of silver was deposited at the 
[Ans. 
3.17 x 104 coulomb] 
[Hint: Iron in the complex is Fe3+ . 
cathode of cell B. How long did the current flow? What 
K
3Fe(CN)6 + 
3e-
----73K+ + 6CN- + Fe] 
masses of copper and zinc were deposited? 
1 mole 
3 x 96500 C 
(At. mass: 
Cu 
63.5, Zn = 65.4, Ag = 107.8)
\end{problembox}

\begin{problembox}
\textbf{Problem 53} \hfill \textbf{[GRB (O.P. Tandon) p.59 Q18]}
\label{q:ext-grb_all_questions-53}

What current stj:ength in ampere will be required to liberate 
[Ans. 865 second; 0.427 g Cu; .0.44 g Zn] 
5 g of iodine from potassium iodide solution in 30 minutes?
\end{problembox}

\begin{problembox}
\textbf{Problem 54} \hfill \textbf{[GRB (O.P. Tandon) p.59 Q27]}
\label{q:ext-grb_all_questions-54}

An electric current passing for 6 minutes through a dilute 
[Ans. 2.11 ampere] 
$/bullet$ 
H2S04 solution gave 40 mL of the electrolytic gas (H2 + 02)
\end{problembox}

\begin{problembox}
\textbf{Problem 55} \hfill \textbf{[GRB (O.P. Tandon) p.59 Q19]}
\label{q:ext-grb_all_questions-55}

The mass of copper deposited from a solution of copper 
measured at NTP. What was the average value of current? 
sulphate by a uniform current of 0.25 ampere flowing for one 
[Ans. 0.638 ampere] 
hour is 0.295 g. Find the equivalent mass of copper. (1 faraday 
[Hint: 96500 coulomb give 11.2 litre of H2 and 5.6 litre of 
= 96500coulomb) 
. 
$/circ$2, i.e .$/bullet$ 16.8litreofelectro1yticgasatNTP.]$/cdot$ 
[Ans. 31.75]
\end{problembox}

\begin{problembox}
\textbf{Problem 56} \hfill \textbf{[GRB (O.P. Tandon) p.59 Q28]}
\label{q:ext-grb_all_questions-56}

The same quantity of electricity that liberated 2.158 g silver
\end{problembox}

\begin{problembox}
\textbf{Problem 57} \hfill \textbf{[GRB (O.P. Tandon) p.59 Q20]}
\label{q:ext-grb_all_questions-57}

The solution of a salt of a metal of atomic mass' 112 was 
was passed through a solution of a gold salt and 1.314 g of 
.-.- -----eIeCtiQIYsecnor1"5'rrlliiliteswrth-acu.rreirtOmampere:-The/textasciitilde()--$/cdot$-7---gotd-wardep'o-sited:11reeqntvalent mass ofSi1ver iSTOT9-. -... --
mass of the metal deposited was O. 788 g. Fmd the valency of 
Calettlate-the-eqtivalent-mass-of-gold:-What is the-oxidatioll 
the metal. 
state of gold in this gold salt? (At. mass of gold = 197) 
[Ans. 2] 
[AlI/textasciitilde()" Eq. mass 
Oxidation state 
3]
\end{problembox}

\begin{problembox}
\textbf{Problem 58} \hfill \textbf{[GRB (O.P. Tandon) p.59 Q21]}
\label{q:ext-grb_all_questions-58}

If a monovalent metal ion carries 1.6 x 10-19 coulomb of/textasciitilde()lj. Calculate the volume of Cl2 at NTP produced during 
electricity, what is the amount of electricity carried by one 
electrolysis of MgCl2 which produces 6.50 g Mg. 
gram molecular-mass of-the metal ions? -
CAt.lllasJ'. QfMg=.=.).$/pm$J).\underline{\hspace{1cm}} \underline{\hspace{1cm}} \underline{\hspace{1cm}} 
[Ans. 9.6336 x 104 cOulomb] 
[/)w;. 5.991itre]
\end{problembox}

\begin{problembox}
\textbf{Problem 59} \hfill \textbf{[GRB (O.P. Tandon) p.59 Q22]}
\label{q:ext-grb_all_questions-59}

Calculate approximately how much current is necessary to 
:m. How long does it take to deposit 100 g of Al from an 
produce oxygen gas at the rate of1 mL per second? 
electrolytic cell containing Al20 3 using a current of 125 
[Ans. 
11.23 ampere] 
ampere? 
[Hint: 
I g-equlvalent of oxygen or 22400 mL is liberated by 
[All/textasciitilde(). 8577.8 second] 
96500 coulomb.] 
4 
'.11, 109 fairly concentrated solution of CuS04 .is electrolysed
\end{problembox}

\begin{problembox}
\textbf{Problem 60} \hfill \textbf{[GRB (O.P. Tandon) p.59 Q23]}
\label{q:ext-grb_all_questions-60}

0.5 faraday of electricity was required to deposit all the copper 
using 1.01 faraday of electricity. Calculate the mass of the 
in 500 mL of a copper sulphate solution. What is the normality 
resulting solution. 
of the copper sulphate solution? 
. 
[AII,< 9.6025 g] 
[Ans .. Normality = IN] 
J), The density of copper is 8.94 g mL-1.Find out the number of
\end{problembox}

\begin{problembox}
\textbf{Problem 61} \hfill \textbf{[GRB (O.P. Tandon) p.59 Q24]}
\label{q:ext-grb_all_questions-61}

Copper is deposited at an electrode according to the reaction 
coulombs needed to plate an area lOx 10 cm
1 to a' thickness 
2+ ;: 2e-
.... 
. 
. .' 
of 10-2 cm using CuS04 solution as electrolyte. (At. mass of 
25/textasciitilde() 
Cu 
/textasciitilde() 
Cu. When a current of 1:34 ampere was passed 
Cu 
63.6) 
through a solutionof copper sulphate for 10 hours, 15.885 g of 
. 
[Alii;. 27129.2 coulomb.] 
copper was deposited. 
. ' 
(i) How many coulombs of electricity were passedthfough 
.11, What mass orAg (At. mass 108) could be plated on a spoon 
the solution? 
. . 
from electrolysis ofAgN03 solution by one ampere current for 
10 n:UIJ.utes? 
. 
(Dh;<nb:uU990) 
(ii) How many moles of copper were deposited? 
. 
[.'.lb 0.6715 g] 
(iii) How many coulombs are required for the deposition of 1 
'! If a current of 0.3 ampere is drawn from a Daniell cell for I 
mole of copper? 
hour, what would be the change in mass of electrodes? (At. 
(iv) If the charge on an electron is 1.6 x 10-19 coulomb 
mass ofCu = 63.5 and Zn = 65.31) 
calculate the value of Avogadro's number. .' 
, 
[Ans; (i) 482.40 coulomb (ii) 0.25 mol (iii) 2 x 96500 coulomb 
[",/!,;, 0.356g Cu deposited, 0.366 g ofZn dissolves] 
An element A (At. mass 112) and an element B (At. mass 27) 
(iv) 6.03 x 1023] 
form chlorides. Solutions of these chlorides are electrolysed 
Calculate the mass and volume at NTP of hydrogen and 
separately and it is found that when the same quantity of 
chlorine that will be formed by passing 10,000 coulomb of 
electricity is passed, 5.6 g of A was deposited while in the 
charge through an aqueous solution of potassium chloride. 
other cell only 0.9 g of B was deposited. What is the valency 
The cell reaction is: 
of A if the valency of B is 3? 
(Dhanbad 1991) 
2KCl + 2HOH /textasciitilde() 
2KOH + 02 + H2 
[ 
2] 
[Ans. 0.1036 g Hz, volume 1.063 litre; 3.678 g C12, volume" How many coulombs must be applied to a cell for the 
1.1603 litre] 
electrolytic production of 245 g NaCI04 from NaCI03$/bullet$ The 
anode efficiency for the desired reaction is 60/%. . 
["'Ji. 6.43 X 105 coulomb]
\end{problembox}

\begin{problembox}
\textbf{Problem 62} \hfill \textbf{[GRB (O.P. Tandon) p.60 Q37]}
\label{q:ext-grb_all_questions-62}

Calculate the mass of Hg/textasciitilde()CI2 which can be prepared by the 
reduction of mercury(II) ions in the presence of chloride ions . 
by the passage of5.0 a/textasciitilde()pere current for 3.0 hours. 
[ADS. 
132.34 g] 
[Hint: 
2Hg2+ + 2Cl-' + 2e- /textasciitilde() 
Hg2C12 ] 
2 F 
1 mole
\end{problembox}

\begin{problembox}
\textbf{Problem 63} \hfill \textbf{[GRB (O.P. Tandon) p.60 Q38]}
\label{q:ext-grb_all_questions-63}

How long will it take for a uniform current of 6.0 ampere to 
deposit 78,0 g gold from a solution' of AuCt:j? What mass of 
chlorine gas will be formed simultaneously at the anode in the 
electrolytic cell? 
[ADS. 
19104 second; chlorine liberated = 42.17 g] 
[HiDt: 
AuCe;: + 3e- /textasciitilde() 
Au 
+ 4Cl- (Cathode) 
3F 
I mole 
2Cl- /textasciitilde() 
C12 
+ 2e- (Anode)] 
I mole 
Z F
\end{problembox}

\begin{problembox}
\textbf{Problem 64} \hfill \textbf{[GRB (O.P. Tandon) p.60 Q43]}
\label{q:ext-grb_all_questions-64}

The resistance of 0.01 N solution at 25
DC is 200 $/Omega$. Cell 
constant of the conductivity cell is unity. Calculate the 
equivalent$/cdot$ conductance of the solution. 
(ADS. 
500 $/Omega$-I cm2 eq-I]
\end{problembox}

\begin{problembox}
\textbf{Problem 65} \hfill \textbf{[GRB (O.P. Tandon) p.60 Q44]}
\label{q:ext-grb_all_questions-65}

A conductivity cell was filled with 0.01 M solution of KCl 
which was known to have a specific conductivity of 0.1413 
-uhm-I m- I at 298 K. Its measured resistance at 298 K was 
94.3 $/Omega$. When the cell was filled with 0.02 M AgN03 
solution, its resistance was 50.3 $/Omega$. Calculate (i) cell 
constant and (ii) the specific conductance ofAgN03 solution. 
[ADs. 
13.32m-l ; 2.648 x 10-1 $/Omega$- I m-I ]
\end{problembox}

\begin{problembox}
\textbf{Problem 66} \hfill \textbf{[GRB (O.P. Tandon) p.60 Q45]}
\label{q:ext-grb_all_questions-66}

A conductance cell was calibrated by filling it with a 0.02 M 
solution of potassium chloride. (specific 
conductance 
= 0.2768 $/Omega$-l m-I) and measuring the resistance at 298 K 
which was found to be 457.3 $/Omega$. The cell was then filled
\end{problembox}

\begin{problembox}
\textbf{Problem 67} \hfill \textbf{[GRB (O.P. Tandon) p.60 Q39]}
\label{q:ext-grb_all_questions-67}

How long will it take 5 ampere of current to def)osit 2 g of 
with a calcium chloride solution containing 0.555 g of CaCl2 
... 
/textasciitilde()/textasciitilde()/textasciitilde()/textasciitilde()/textasciitilde()/textasciitilde()/textasciitilde()/textasciitilde()/textasciitilde()/textasciitilde()/textasciitilde()/textasciitilde()\underline{\hspace{1cm}}/textasciitilde() 
\underline{\hspace{1cm}} /textasciitilde() \underline{\hspace{1cm}} 
/textasciitilde()--:-\underline{\hspace{1cm}} /textasciitilde()/textasciitilde() 
\underline{\hspace{1cm}} /textasciitilde() \underline{\hspace{1cm}} /textasciitilde() \underline{\hspace{1cm}} 
/textasciitilde()$/cdot$/textasciitilde()C/textasciitilde()oEPp/textasciitilde()e/textasciitilde()rc.!fr/textasciitilde()o/textasciitilde()m/textasciitilde()a/textasciitilde()s/textasciitilde()o/textasciitilde()lu,:,t/textasciitilde()io/textasciitilde()n/textasciitilde()o/textasciitilde()f/textasciitilde()c/textasciitilde()o:jp/textasciitilde()p/textasciitilde()e/textasciitilde()r/textasciitilde()su,:,l/textasciitilde()p/textasciitilde()h/textasciitilde()at/textasciitilde()e/textasciitilde()?/textasciitilde()/textasciitilde()/textasciitilde()\underline{\hspace{1cm}}-,--\underline{\hspace{1cm}} ----.!per litre. The measured resistance was 1050 $/Omega$. Calculate 
(Given, CE of copper = 32, F = 96500 coulomb) 
the molar conductlVlty of Cli02 solutIOn. 
(Ranchi 1996) 
[ADS. 
0.0241 mho m2 mol-I] 
[ADS. 
20.1 minute] 
[Hint: -Determine cell constant with KCl data. its value IS
\end{problembox}

\begin{problembox}
\textbf{Problem 68} \hfill \textbf{[GRB (O.P. Tandon) p.60 Q40]}
\label{q:ext-grb_all_questions-68}

The amount of lactic acid, HC3/textasciitilde()03' produced in a sample of 
126.6 m -I . 
muscle tissue was analysed by reaction with hydroxide ion. 
1 
. --Hydroxide--ioti -was'--ptbdUc-ed- in-the--sample--mixrure by 
Sp.\underline{\hspace{1cm}} condo of CaC!2 soln. = Cell cons!. x Resistance 
---\underline{\hspace{1cm}} ... \underline{\hspace{1cm}} .. 
electrolysis. The cathode reaction was, 
I 
I 
1 
= 126.6 x -- = 0.1206 $/Omega$- rri-
, 
2Hz0(l) + 2e- ----7 Hz(g) + 20W(aq.) 
Hydroxide ion rea<;;ts with lactic acid as soon as it is produced. 
The end point 
0 of the reaction is detected with an acid-base 
indicator. It required 115 seconds for a current of 15.6 rnA to 
reach end point. How many grams of lactic acid (a monoprotic 
acid) were present in the sample? 
[ADS. 
1.674 x 10-3 g] 
[Hint: No. of moles oflactic acid = No. of moles ofOH - used 
,;, No. of faraday used in 
electrolysis 
N 
b 
f
Co 
$/circ$d$/cdot$ 
d$/cdot$ I x t 
15.6 X 10-3 x 115 
urn er 0 lara ay use = -- = 
. 
. 
96500 
96500 
= 1.86 x 10-5 
Mass of lactic acid = Number of moles x Molecular mass 
= 1.86 x 10-5 x 90 
= 1.674 X 10-3 g]
\end{problembox}

\begin{problembox}
\textbf{Problem 69} \hfill \textbf{[GRB (O.P. Tandon) p.60 Q41]}
\label{q:ext-grb_all_questions-69}

In whatdirection, can the reaction: 
2NaCI + Fez (S04)3 /textasciitilde() 
2FeS04 + Clz + Na ZS04 
proceed spontaneously? 
(Given:-E;e3+/Fe2+ = O.77V; 
E/textasciitilde()I-/CI2 = 1.36V) 
(BCECE 2005) 
[ADS. 
Forward]
\end{problembox}

\begin{problembox}
\textbf{Problem 70} \hfill \textbf{[GRB (O.P. Tandon) p.60 Q42]}
\label{q:ext-grb_all_questions-70}

How many faradays of electricity will be required to 
completely electro lyse one mole of molten Al20 3 to produce 
Al metal and 0z gas? 
IJEE (West Bengal) 200S1 
[ADS. 
3 F] 
Sp. condo 
Molar conductance = ....!....--
Conc. 
1050 
Cone. = 0.555 = 0.005 mol dm-3 = 5 mol m-3 ] 
111.0
\end{problembox}

\begin{problembox}
\textbf{Problem 71} \hfill \textbf{[GRB (O.P. Tandon) p.60 Q46]}
\label{q:ext-grb_all_questions-71}

The molar conductivities at infinite dilution of Kel, 
KN03 and AgN03 at 298 K are 0.01499 mho mZ mol-I, 
0.01250 mho mZ mol-J and 0.01334 mho mZ mo]-J 
respectively. What is the molar conductivity of AgCl at 
il)finite dilution at this temperature? 
[ADS. 
0.01383 mho m2 mol-I] 
47, The electrodes in a conductivity cell have area 1.2 x 10-4 m2 
and they are fixed 3 x 10-3 m apart. A solution containing 200 
g equivalent. of the electrolyte per m3 of the solution has a 
resistance of 60 $/Omega$ at 298 K. Calculate the equivalent 
conductivity of the solution. 
[ADS. 
2.09 x lO-z mho m2 eq-I]
\end{problembox}

\begin{problembox}
\textbf{Problem 72} \hfill \textbf{[GRB (O.P. Tandon) p.60 Q48]}
\label{q:ext-grb_all_questions-72}

The molar conductivities of NH4+ ion and Cl- ion are 73.5 
mho cmz mol-I and 76.2 mho cm2 mo]-I respectively. The 
specific 
conductivity 
of 
0.1 
M 
NH4CI 
is 
. 1.288 X lO-z mho cm-I. What is the dissociation constant of 
NH4Cl? 
[Ans. 
0.86]
\end{problembox}

\begin{problembox}
\textbf{Problem 73} \hfill \textbf{[GRB (O.P. Tandon) p.60 Q49]}
\label{q:ext-grb_all_questions-73}

The specific conductivity of a saturated solution of silver 
chloride at 18 DC is 1.24 X 10-6 mho after subtracting that of 
water. Ionic conductances at infinite dilution of Ag + and Cl-
ions at this temperature are 53.8 and 65.3 respectively. 
Calculate the solubility of silver chloride in gram per litre. 
[ADS. 
1.494 x 10-3 g per litre] 
I
\end{problembox}

\begin{problembox}
\textbf{Problem 74} \hfill \textbf{[GRB (O.P. Tandon) p.61 Q50]}
\label{q:ext-grb_all_questions-74}

Given: 
A/textasciitilde() LYz Mg2+] = 53.06 $/Omega$-I cm2 mol-I 
A/textasciitilde() U'i SO/textasciitilde()-] = 80 $/Omega$-I cm2 mol-1 
A/textasciitilde() LX AI3+] = 63 $/Omega(-1)$ cm2 mol-1 
Calculate the values of A:, [AI2(S04h] and A/textasciitilde() [MgS04]. 
[Ans. 858 $/Omega$-I cm2 mol-I, 266.12 $/Omega$-I cm2 mol-I] 
(i) Iron reduces Zn 2+ ions (ii) Iron reduces Ni2+ ions. 
Given: EO 2+ 
= - 0.76 volt; EO 2+ 
= - 0.44 volt and 
Zn 
IZn 
Fe 
IFe 
E/textasciitilde()i 2+ INi = - 0.25 volt 
[Ans. Iron reducesNi 2+ ions as the E;ell in this case has positive 
value.] 
57/textasciitilde() Can a solution of 1 M CuS04 be stored in a vessel made of 
nickel metal? 
Given: 
EO 2+ 
= - 0.25 volt 
(Ni 
;Ni)
\end{problembox}

\begin{problembox}
\textbf{Problem 75} \hfill \textbf{[GRB (O.P. Tandon) p.61 Q51]}
\label{q:ext-grb_all_questions-75}

Hydrofluoric acid is a weak acid. At 25$/circ$C; the molar 
conductivity of 0.002 M HF is 176.2 $/Omega$-I cm2 mol-I. If its 
A';' = 405.2 $/Omega$-I cm2 mol-I, 
calculate 
its 
degree 
of 
. 
EOc 2+ 
C = 0.34 volt 
dissociation 
and 
equilibrium 
constant 
at 
the 
given 
(u 
; u) 
concentration. 
Ans. No; The reaction takes place between CuS04 and nickel 
[Ans .. a= 0.435, K = 6.7 x 10-4] 
s E/textasciitilde()II is positive.]
\end{problembox}

\begin{problembox}
\textbf{Problem 76} \hfill \textbf{[GRB (O.P. Tandon) p.61 Q52]}
\label{q:ext-grb_all_questions-76}

For mercury at O$/circ$C specific conductance,
\end{problembox}

\begin{problembox}
\textbf{Problem 77} \hfill \textbf{[GRB (O.P. Tandon) p.61 Q58]}
\label{q:ext-grb_all_questions-77}

An . .ver whether, under standard conditions, the following 
reactions are possible or not: 
--------------------------------------------------------------------------------------------------------------------/textasciitilde():,--/textasciitilde()--------------/textasciitilde()-*K/textasciitilde()+.1.Q/textasciitilde()/textasciitilde()/textasciitilde()/textasciitilde()m/textasciitilde()./textasciitilde()-------------------------(/textasciitilde()I/textasciitilde())/textasciitilde()Wrr$/cdot$ 
1/textasciitilde()II/textasciitilde()c/textasciitilde()o/textasciitilde()p=p/textasciitilde()cr/textasciitilde()r/textasciitilde()e/textasciitilde()du/textasciitilde()c/textasciitilde()e/textasciitilde()A'g/textasciitilde()+/textasciitilde()/textasciitilde()t/textasciitilde()o/textasciitilde()A'g/textasciitilde()?/textasciitilde()Glven: E/textasciitilde()g+/Ag 
--------------------------------/textasciitilde()----/textasciitilde()------------------/textasciitilde()--------------------------/textasciitilde()----/textasciitilde()----------------/textasciitilde()----/textasciitilde()$/cdot$,--------------,(a) If the resistance of a eell contammg mercury IS O.243T66 
= 0.799 volt; EO 2+ 
= - 0.337 volt 
. 
Cu 
ICu 
$/Omega$, what is the cell constant of the cell? 
(ii) Will Fe3+ be reduced to Fe2+ by Sn 2+ ion? Given: 
(b) If the same cell is filled with KCI solution at O$/circ$C, the 
Fe3+ I Fe2+ = 0.771 volt; 
resistance of the cell is 3.966 x 104 $/Omega$. What is the 
conductivity ofKCI solution? 
Sn 2+ I Sn4+ = - 0.250 volt 
\underline{\hspace{1cm}}J(;) If the 
averag/textasciitilde()\underline{\hspace{1cm}}/textasciitilde()ss-s/textasciitilde()ctio!lllLJ!/textasciitilde()iL9L.JhtL\underline{\hspace{1cm}}/textasciitilde()1l/textasciitilde()j/textasciitilde()L 
---(iii)-Would-yoU-Use-a-silver-spoon-to-stir-a-solution....gf--. 
0.9643 mm2, what is the effective distance between the 
Cu(N03h? 
electrodes? 
[Ans. 
(i) Yes, E;ell Cu I C;u2+ II Ag+ I Ag is positive (+ 0.462 
[Ans. (a) 2:58476 x 105 S/m; (b) 6.517 S/m; (c) 0.2492 m] 
volt)
\end{problembox}

\begin{problembox}
\textbf{Problem 78} \hfill \textbf{[GRB (O.P. Tandon) p.61 Q53]}
\label{q:ext-grb_all_questions-78}

The mobility' of the NH4+ ions is 7.623 x 10-8 m2 I Vs. 
(ii) Yes, 
E;ell 
Pt I Sn4+, Sn2+ II Fe3+, Fe2+ I Pt is positive 
(+0.621 volt) ." 
Calculate: 
.:; 
(iii) Yes, the reaction does not occur.] 
(a) the molar conductivity ofNH4+ ion;
\end{problembox}

\begin{problembox}
\textbf{Problem 79} \hfill \textbf{[GRB (O.P. Tandon) p.61 Q59]}
\label{q:ext-grb_all_questions-79}

The electrode potentials of two half reactions are as follows: 
(b) velocity of the ion if 15 volt are applied across the 
electrodes 25 cm apart. 
[Ans. 
(a) 73.55 x 10- 4 S m 2/mol; (b)4.574I-lrnls] 
54; . Construct electric cells for the following reactions: 
(a) Fe + Cu 2+ -----7 Cu + Fe2+ 
(b) Cu + 2Ag + -----7 2Ag + Cu 2+ 
(c) 6Fe2+ + Cr20/textasciitilde()- + 14H + -----7 6Fe3+ + 2Cr3+ + 7H20 
(d) Cd + 12 -----7 Cd2+ + 2r 
(e) 2Fe3+ + 2Cl- -----7 2Fe2+ + CI2 
[Ans. (a) Fel Fe2+ II Cu2+ I Cu 
(b) Cu I Cu2+ II Ag+ lAg 
(c) Pt I Fe2+, Fe3+ II H +, Cr20/textasciitilde()-, Cr3+ I Pt 
(d) Cd I Cd2+ II r, 12 I Pt 
(e) Pt I C12, Cl- II Fe3+, Fe2+1 Pt ] 
5/.. Predict 
whether 
the 
following 
reaction 
will 
occur 
spontaneously or not: 
C02+ + Sn -----7 Co + Sn 2+ 
EO 
2+ = 0.277 voli and 
EO 
2+ = 0.136 volt 
Col Co 
t 
Sn/Sn 
[Ans. 
No; E;ell is negative = - 0.141 volt]
\end{problembox}

\begin{problembox}
\textbf{Problem 80} \hfill \textbf{[GRB (O.P. Tandon) p.61 Q56]}
\label{q:ext-grb_all_questions-80}

An iron wire is immersed in a solution containing ZnS04 and 
NiS04.Predict giving reasons which of the following 
reactions is likely to proceed? 
Fe3+ + e - -----7 Fe2+ .,' 
EO = + 0.76 volt 
Ce4+ + e- /textasciitilde() 
Ce3+ ; 
EO = + 1.60 volt 
Giving reason, d/textasciitilde()scribe iCe3+ can be oxidised by Fe3+. 
[Hint: Ce3+ + Fe3+ /textasciitilde() 
Ce4+ + Fe2+ 
Ce3+ /textasciitilde() 
ee4+ + e- (oxidation); EO = -l.60 volt 
Fe3+ + e- /textasciitilde() 
Fe2+ 
(reduction); EO = + 0.76 volt 
Adding both, EO ofthe reaction = - 0.84 volt 
The oxidation of Ce3+ by Fe3+ is not possible.]
\end{problembox}

\begin{problembox}
\textbf{Problem 81} \hfill \textbf{[GRB (O.P. Tandon) p.61 Q60]}
\label{q:ext-grb_all_questions-81}

Calculate the standard reduction potential of Ni 2+ I Ni 
electrode when the emf of the cell 
Ni I Ni2+ (1 M) II Cu 2+ (1 M) I Cu 
is 0.59 volt and Eco 2+ 
= + 0.34 volt 
. 
u 
ICu 
[Ans. 
- 0.25 volt]
\end{problembox}

\begin{problembox}
\textbf{Problem 82} \hfill \textbf{[GRB (O.P. Tandon) p.61 Q61]}
\label{q:ext-grb_all_questions-82}

Calculate the half-ceil potential at 298 K for the reaction, 
Zn 2+ + 2e - -----7 Zn 
if[Zn 2+] = 0.1 M and EO = - 0.76 volt 
[Ans. 
-:- 0.789 volt]
\end{problembox}

\begin{problembox}
\textbf{Problem 83} \hfill \textbf{[GRB (O.P. Tandon) p.61 Q62]}
\label{q:ext-grb_all_questions-83}

A galvanic cell consists of a metallic zinc plate iinmersed in 
0.1 M Zn(N03h solution. Calculate the emf of the cell at 
25$/circ$C. Write the chemical equations for the electrode reactions 
and represent the cell. (Given: EO 2+ 
= - 0.76 volt and 
Zn 
IZn
\end{problembox}

\begin{problembox}
\textbf{Problem 84} \hfill \textbf{[GRB (O.P. Tandon) p.62 Q70]}
\label{q:ext-grb_all_questions-84}

Calculate the emf of the cell 
Zn I Zn 2+ II Pb2+ I Pb 
0.1 M 
0.02 M 
Cr I Cr3+ (0.1 M) II Fe2+ (0.01 M) I Fe 
 :ell = 0.63 volt 
Eeell = 0.609 volt] 
(Given: EO 3+ 
= -'0.75 volt; 
Cr 
ICr 
IFe = - 0.45 volt)
\end{problembox}

\begin{problembox}
\textbf{Problem 85} \hfill \textbf{[GRB (O.P. Tandon) p.62 Q63]}
\label{q:ext-grb_all_questions-85}

Calculate the cell potential of the cell at 25$/circ$C having the cell 
[Alis. 
0.261 volt] 
reaction:
\end{problembox}

\begin{problembox}
\textbf{Problem 86} \hfill \textbf{[GRB (O.P. Tandon) p.62 Q71]}
\label{q:ext-grb_all_questions-86}

Calculate the equilibrium constant for the reaction, 
2Ag+ + Cu 
2Ag + Cu 2+ 
(Given that, 
[Ag+] 
1.0 x 10-3 M, 
Zn(s) + Ag20(s) + H20(l) 
2Ag(s) + Zn 2+(aq.) 
[Cu 2+]=1.0.x1O-4M, E:olI = 0.46 volt) 
[Ans. 
0040 volt] 
when E;ell = 1.11 at 298 K.
\end{problembox}

\begin{problembox}
\textbf{Problem 87} \hfill \textbf{[GRB (O.P. Tandon) p.62 Q64]}
\label{q:ext-grb_all_questions-87}

Calculate the E and 
of the cell 
[Ans. Kc "'" 30499 X 1037] 
Ni I Ni2+ II.Cu2+ I Cu
\end{problembox}

\begin{problembox}
\textbf{Problem 88} \hfill \textbf{[GRB (O.P. Tandon) p.62 Q72]}
\label{q:ext-grb_all_questions-88}

The standard reauction potential ofAg+ lAg electrode is 0.80 
from the followirig half-cell reactions: 
volt. 
Calculate 
the 
standard 
electrode' potential 
of 
Ni2+ + 2e- ----7 Ni; 
EO 
- 0.25 volt 
Cl- I AgCII Ag 
at 
25'oC. 
Giyen 
solubility 
product, 
------fCC'ffUZ$/pm$+-2/textasciitilde() ) Cu, 
EO 
+ 0.34 volt 
Ksp-(AgCI) 
1.2 x 10-10$/bullet$ 
------t(ftGi/'en:1Ni:2+] \underline{\hspace{1cm}} 1 M andreu.2+J/textasciitilde()i11O'F-/textasciitilde()3 
"*lv..!,(t", )'1:-. --c--/textasciitilde()------fl:[Hnl"'nt .. 
: /textasciitilde()Ap'p1y--:-----------------
. [Ans. E:'l1 = 0.59 volt; Eeell 
0.5015 volt] 
E:ell = 0.0591 x log [Ag+ HCn 0.05911<,>g Ksp(AgCI)
\end{problembox}

\begin{problembox}
\textbf{Problem 89} \hfill \textbf{[GRB (O.P. Tandon) p.62 Q65]}
\label{q:ext-grb_all_questions-89}

Use EO values to calculate t:..Go for the reaction 
E:ell = - 0.576 volt 
+ 
Fe3+ +Ag 
EO 4.-=-0.80-volt ,and  0"3+\underline{\hspace{1cm}}2+::;:\underline{\hspace{1cm}}0.72-volt 
Ag lAg 
PtIFe 
,Fe 
"-
[Ans. 
2895 J] 
[Hint: E:ell 
0.03 volt, t:..Go = -nFE:ell ]
\end{problembox}

\begin{problembox}
\textbf{Problem 90} \hfill \textbf{[GRB (O.P. Tandon) p.62 Q66]}
\label{q:ext-grb_all_questions-90}

A piece of zinc metal is dipped in a 0.1 M solution of zinc salt. 
The salt is dissociated to the extent of 20/%. Calculate the 
electrode potential of Zn 2+ I Zn. 
(Given: E;n2+ /Zn 
- 0.76 volt.) 
[Ans. 
/textasciitilde() 0.8102 volt] 
[Hint: 
[Zn2+] 
0.1 x 0.2 
0.02 M 
(Reduction) 
E 
-
Zn2+ IZn -
0.0591 I 
[Z 2+.].] 
+-\underline{\hspace{1cm}}$/cdot$oglO n 
IZn 
2 
.
\end{problembox}

\begin{problembox}
\textbf{Problem 91} \hfill \textbf{[GRB (O.P. Tandon) p.62 Q67]}
\label{q:ext-grb_all_questions-91}

Estimate the concentration limit beyond which the half-cell 
potential of Cu 2+ I Cu will be zero. 
[Ans. 
[Hint: 
= 0.34 volt 
68 .. Calculate equilibrium constant for the following reaction: 
-Zn + CuS04 ----7 ZnS04 + Cu 
'0 
EZnIZn2+ 
0.765 volt; 
[Ans.
\end{problembox}

\begin{problembox}
\textbf{Problem 92} \hfill \textbf{[GRB (O.P. Tandon) p.62 Q1]}
\label{q:ext-grb_all_questions-92}

862 x 1037] 
69., For the cell reaction, 
EO 2+ C = 0.347 volt 
Cu 
I u 
Ni I Ni2+ II Ag + lAg 
. 
, 
Calculate the equilibrium constant at 25$/circ$C. How much 
maximum work would be obtained by the operation of this 
cell? 
. 
[Ans. 
3.98 x 1035 , Max. work 
202650 J] 
t'. 
E:eu = Oxid. pot. anode + Red. pot. cathode 
-Red.pot./textasciitilde()athode[Cl=IAgCI-IAg]=-:0..5-76."7(--O.80)--- ----
0.5.16 + 0.80 = 0.224 volt] 
Derive Nemst equation for the cell. 
Ni(s) I Ni2+ (aq. 0.1 M) II Ag+(aq. 0.1 M) I Ag(s) 
and also find its cell potential. Given: 
EO + 
0.80 volt and ENo" 2+ /N1. = -,0.25 volt 
Ag lAg 
[Hint: 'Cell reaction, 
Ni + 2Ag+ 
2Ag + Ni2+; 
E:ell = 1.05 volt 
0.059 i I 
[Ni 2+ ) 
Nemst equation, Eeoll 
E:olI 
og --
2 
[Ag+]2 
(1.05 
0.0295) volt 
= 1.02 volt]
\end{problembox}

\begin{problembox}
\textbf{Problem 93} \hfill \textbf{[GRB (O.P. Tandon) p.62 Q74]}
\label{q:ext-grb_all_questions-93}

Determine the equilibrium constant of the following reaction 
at 298 K: 
2Fe3+ + Sn 2+ ----7 2Fe2+ + Sn 4+ 
(Given: EsO 4+ s 2+= 0.15 volt, EO 3+ 
2+ = 0.771 volt) 
n 
I n 
Fe 
IFe 
[Ans. 
K 
1.0 x 1021] 
[Hint: 
Calculate E:.II $/bullet$ The value is 0.621 volt. 
$/circ$ 
0.0591 
Apply Eeen = 
log K]
\end{problembox}

\begin{problembox}
\textbf{Problem 94} \hfill \textbf{[GRB (O.P. Tandon) p.62 Q75]}
\label{q:ext-grb_all_questions-94}

The zinc-silver oxide cell is used in hearing aids and electric 
watches. From the following data: 
+ 2e-
Zn; 
EO = - 0.76 volt 
Ag20 + H20 + 2e - ----7 2Ag + 20fr; 
EO = 0.334 volt 
Answer the following: 
(a) Construct the cell and what will be its emf? 
(b) What is its cell reaction and what will be t:..Go value of this 
cell reaction? 
I
\end{problembox}

\begin{problembox}
\textbf{Problem 95} \hfill \textbf{[GRB (O.P. Tandon) p.63 Q76]}
\label{q:ext-grb_all_questions-95}

An excess of Hg was added to 10-3 M acidified solution of 
Fe3+ ions. It was found that only 4.6/% of the ions remained as 
Fe3+ at equilibrium at 25$/circ$C. Calculate EO for 2Hg I 
at 
25$/circ$C for 
2Hg + 2Fe3+ /textasciitilde() 
Hg/textasciitilde()+ + 2Fe2+ 
(Ans. 
- 0.7912 volt] 
[Hint: 
By passage of 105 
coulomb charge, 
zinc 
ion 
concehtration increases by 0.518 mole while copper ion 
concentration decreases by 0.518 mole. 
h 
"+ ' 
2+ 
T us, [Zn- ]= 1.518Mand[Cu 
] 0.482M 
(Now apply Nemst equation)]
\end{problembox}

\begin{problembox}
\textbf{Problem 96} \hfill \textbf{[GRB (O.P. Tandon) p.63 Q81]}
\label{q:ext-grb_all_questions-96}

Calculate the standard potential for the reaction, 
Hg2Cl2 + Cl2 /textasciitilde() 
2Hg2+ + 4Cl-
(Given: 
Hg2Cl2 + 2e- /textasciitilde() 
2Hg + 2Cr; 
Hgi+ /textasciitilde() 
2Hg2+ + 2e -; 
2Hg /textasciitilde() 
Hg;+ + 2e-; 
Cl2 + 2e- /textasciitilde() 
2Cl-; 
Ans. 
- 0.08 volt] 
LHint: 
First determine EO for 
EO 
0.270 volt 
EO 
0.92 volt 
EO 
0.79 volt 
EO 
1.36 volt) 
---c----/textasciitilde()AgJ..AgL(()j15f-MJ(l\underline{\hspace{1cm}}j,WW5)M/textasciitilde()g-- 
is 0.788 volt Calculate the solubility product of AgI. 
--------uIt/textasciitilde()c/textasciitilde()olilln/textasciitilde()e/textasciitilde()/textasciitilde()/textasciitilde()/textasciitilde()--------
JAns. 1.10 x 10-16]
\end{problembox}

\begin{problembox}
\textbf{Problem 97} \hfill \textbf{[GRB (O.P. Tandon) p.63 Q78]}
\label{q:ext-grb_all_questions-97}

At equimolar concentrations of Fe2+ and Fe3+, what must 
[Ag+] be so that the voltage of the galvanic cell made from 
Kg+-I]/gandP-e3+1'Fe2+el-ect1'ooe/textasciitilde()requals zero ?Th:eTeacfioffls 
Fe2++ Ag+ /textasciitilde() 
Fe3+ + Ag. Detennin.e 
t/textasciitilde()e, /textasciitilde()equilibrium 
. 
. 
. 
-0' 
c/textasciitilde()/textasciitilde()tant at 25"Cfor thereaction. (Given: E.Ag + lAg 
0.799 volt 
and E;e3+ IFe2 + = 0.771 volt ) 
[Hint: 
= -0,771 + 0.799 = 0.028 volt 
At/textasciitilde()quilibrium, Eeell = 0 
0- EO \underline{\hspace{1cm}} 0.0591 1 
-
cell 
I' og 
]lAg+ 1 
$/circ$ 
1 
= EceU - 0.0591 log --
JAg+]= 0.34 
10 K = nEo 
g 
0.0591 
K =3.0] 
[Ag+]
\end{problembox}

\begin{problembox}
\textbf{Problem 98} \hfill \textbf{[GRB (O.P. Tandon) p.63 Q79]}
\label{q:ext-grb_all_questions-98}

Using Nemst equation for the cell reaction, 
Pb + Sn 2+ /textasciitilde() 
Pb2+ + Sn 
i 
; 
i 
. 
[Pb2+] 
, 
Calculate the ratIO --2- for which Eccll 
[Sn +] 
O. 
(Given: E;b/Pb2+ ='0.13 volt and E:n2+/sn ::: 
0.14 volt) 
, 
[Pb2+] 
(Ans. 
--2 -
= 0.5] 
. 
[Sn +] 
I
\end{problembox}

\begin{problembox}
\textbf{Problem 99} \hfill \textbf{[GRB (O.P. Tandon) p.63 Q80]}
\label{q:ext-grb_all_questions-99}

, Determine the potential of a Daniell cell, initially containing 
1.00 L each of 1.0 M copper(II) ion and 1.0 M zinc(II) ion, 
after passage of 105 coulomb charge. The E" of the Daniell 
cell is 1.1 0 volt. 
[Ans. 
1.09 volt] 
This electrode is now coupled with elz + 2e- ----7 2Cl-
electrode. ]
\end{problembox}

\begin{problembox}
\textbf{Problem 100} \hfill \textbf{[GRB (O.P. Tandon) p.63 Q82]}
\label{q:ext-grb_all_questions-100}

Given: 
EO 
0.15 volt 
. \underline{\hspace{1cm}} Cu/textasciitilde()\underline{\hspace{1cm}}+\underline{\hspace{1cm}}e 
\underline{\hspace{1cm}} u""LCU;\underline{\hspace{1cm}}. \underline{\hspace{1cm}} 
... \underline{\hspace{1cm}}.\underline{\hspace{1cm}}E/textasciitilde() .... ::;,\underline{\hspace{1cm}}05 volL. ""n"/textasciitilde()/textasciitilde() \underline{\hspace{1cm}} 
n 
Calculate potential for Cll 2+ + 2e- /textasciitilde() 
Cu. 
[Ans. 
0.325 volt]
\end{problembox}

\begin{problembox}
\textbf{Problem 101} \hfill \textbf{[GRB (O.P. Tandon) p.63 Q83]}
\label{q:ext-grb_all_questions-101}

What is the standard electrode potential for the electrode 
Mn04/Mn02 in solution?
\end{problembox}

\begin{problembox}
\textbf{Problem 102} \hfill \textbf{[GRB (O.P. Tandon) p.63 Q84]}
\label{q:ext-grb_all_questions-102}

(Given: EO 0-
"+ = 1.51 volt, 
Mn 4/Mn-
[Ans. 
1.7 volt] 
[Hint: 
1.23 volt) 
MnO;j + 8H+ + 5e- ----7 Mn 2+ + AH20; EO 
LSI volt 
Mn02 + 4H+ + 2e- ----7 Mn2+ + 
EO 
1.23 volt 
MnO;j + 4H+ +3e- ----7 Mn02 + 2H20; EO =? 
EO '" 5 x 1.51- 2 x 1.23 = 7.55 - 2.46 
5.09 
1.70 volt] 
333 
What ratio of Pb2+ to Sn 2+ concentration is needed to reverse 
the following cell reaction? 
Sn(s) + Pb2+(aq.) /textasciitilde() 
Sn2+(/textasciitilde()q.) + Pb(s) 
EO 
= - 0.136 volt and 
Sn 2+ ISn 
0.126 volt 
[Ans.
\end{problembox}

\begin{problembox}
\textbf{Problem 103} \hfill \textbf{[GRB (O.P. Tandon) p.63 Q85]}
\label{q:ext-grb_all_questions-103}

for the cell Mg I Mg2+ II Ag+ I Ag, calculate the equilibrium 
constant at 25$/circ$C and the maximum work that can be obtained 
during the operation of the cell. 
(Given: EO 
M?+ = 2.37 volt and 
0.80 volt) 
Mg/l g-
. 
[Ans. 
K= 2.86 x 10107; Wm"" = 6.118?< 105 1].
\end{problembox}

\begin{problembox}
\textbf{Problem 104} \hfill \textbf{[GRB (O.P. Tandon) p.64 Q86]}
\label{q:ext-grb_all_questions-104}

Detem1ine the potential for the cell: 
Pt I Fez. Fc3+ II Cr20/textasciitilde()-, Cr3+, H + I Pt 
in which [Fe l '] and [Fe3+ 1 are O.5M and O.75Mrespectively 
and 
rCr20/textasciitilde()-l[Crl+] ancl [H+] are 2M, 4M and 1M 
respectively. Given: 
fe 3+ + e- /textasciitilde() 
Fe"+ ; 
EO = 0.770 volt 
14H + + 6e- + Cr20/textasciitilde()- = 2Cr3+ + 7H20; 
EO = 1.35 volt 
[Ans. 
0.56 volt]
\end{problembox}

\begin{problembox}
\textbf{Problem 105} \hfill \textbf{[GRB (O.P. Tandon) p.64 Q87]}
\label{q:ext-grb_all_questions-105}

For the measurement of the solubility. product of AgCl the 
following cell is constructed: 
Ag I AgClIl KCl (0.1 Af) II AgN03 (0.1 M) lAg 
Thc emf of the cell is 0.45 volt. In the cell, KCI is dissociated 
to the extent of 83/% and AgNO, is dissociated to the extent of 
86/%: Calculate the solubility p/textasciitilde()oduct of AgCl at 298 K. 
[Ans. 
1.735 x IO-I/textasciitilde() 
. 
[Hint: . Eeell = E;ell - 0.0591 log ..:.[A---=::g,...+::..cJA",no""d"-.e ] 
[Ag + lcathode
\end{problembox}

\begin{problembox}
\textbf{Problem 106} \hfill \textbf{[GRB (O.P. Tandon) p.64 Q88]}
\label{q:ext-grb_all_questions-106}

Excess of AgCl is added to 0.1 M solution of KBr at 298 K. 
Calculate the equilibrium concentrations ofBr- and Cl- ions. 
ECol- 1 
. 
= 0.222 volt 
EO \underline{\hspace{1cm}} 
I 
= 0.095 volt 
AgCI! Ag 
. 
Br I AgBr Ag 
[Ans. 
[CI-] = 0.09929 mol L-I $/bullet$ [Br-] = 0.00071 mol L-I ] 
[Hint: 
AgCI + 
Br-
/textasciitilde() 
AgBr + CI-
At equilibrium. 
(0.1 - x) 
x 
Apply 
1 
K 
E;cll 
ou =--
b 
0.0591 
or 
$/circ$ 
I 
x 
Eeell ] 
og (O.I-x) = 0.0591
\end{problembox}

\begin{problembox}
\textbf{Problem 107} \hfill \textbf{[GRB (O.P. Tandon) p.64 Q89]}
\label{q:ext-grb_all_questions-107}

The emf of the cell Cu I CuS04 II CuS'04 I Cu is - 0.03 Vat 
a = 0.1 
x 
25$/circ$C. Calculate the activity of copper sulphate solution in the 
right hand side electrode. 
[Ans. 
0.009616 mol L- I] 
[Hint: 
Apply E = (E/textasciitilde()u - E/textasciitilde()u) -
0.0/textasciitilde()9lI0g /textasciitilde()/textasciitilde()I]
\end{problembox}

\begin{problembox}
\textbf{Problem 108} \hfill \textbf{[GRB (O.P. Tandon) p.64 Q90]}
\label{q:ext-grb_all_questions-108}

The cell, 
Pt I H2 (l atm) II H + (pH = x) I 
normal calomel electrode, has an enif of 0.67 volt as 25$/circ$C. 
Calculate the pH of the solution. The oxidatiOli potential of 
calomel electrode on the hydrogen scale is - 0.28 volt. 
[Ans. 
pH = 6.61]
\end{problembox}

\begin{problembox}
\textbf{Problem 109} \hfill \textbf{[GRB (O.P. Tandon) p.64 Q91]}
\label{q:ext-grb_all_questions-109}

The emf of the cell, 
Ag I AgI, Kl (0.05 AI) II AgN03 (0.05 M).I Ag 
is 0.7R8 volt. Calculate the solubility product of AgI. 
[Ans. 
1.101 x 10-16]
\end{problembox}

\begin{problembox}
\textbf{Problem 110} \hfill \textbf{[GRB (O.P. Tandon) p.64 Q92]}
\label{q:ext-grb_all_questions-110}

The standard reduction potential for Cu 2+ / Cu is + 0.34 V. 
Calculate the reduction potential at pH = 14 for the above 
couple. Ksp ofCu(OH)2 is 1.0 X 10- 19 . 
[Ans. 
- 0.22 V] 
(TIT 1996)
\end{problembox}

\begin{problembox}
\textbf{Problem 111} \hfill \textbf{[GRB (O.P. Tandon) p.64 Q93]}
\label{q:ext-grb_all_questions-111}

Calculate the potential ofa cell in which H electrode is 
immersed in a solution of pH 3.5 and in a solution of pH 10.7. 
[Ans. 
0.4255 volt]
\end{problembox}

\begin{problembox}
\textbf{Problem 112} \hfill \textbf{[GRB (O.P. Tandon) p.64 Q94]}
\label{q:ext-grb_all_questions-112}

Calculate the emf of the cell, 
Cll (I atm) I NaCl (aq.) I Clz (0.1 atm) 
[Ans. 
- 0.0295 volt]
\end{problembox}

\begin{problembox}
\textbf{Problem 113} \hfill \textbf{[GRB (O.P. Tandon) p.64 Q95]}
\label{q:ext-grb_all_questions-113}

The standard free energy change for the reaction, 
H2(g) + 2AgCI(s) /textasciitilde() 
2Ag(s) + 2H +(aq.) + 2Cl-(aq.) 
is - 10.26 keal mol- l at 25$/circ$C. A cell using the above reaction 
is operated at 25$/circ$C under PH 2 = 1 atm, [H +] and [CluJ = 0.1. 
Calculate the emf of the cell. 
[Ans. 
0.340 volt]
\end{problembox}

\begin{problembox}
\textbf{Problem 114} \hfill \textbf{[GRB (O.P. Tandon) p.64 Q96]}
\label{q:ext-grb_all_questions-114}

For the reaction, Fe3+ + 3e- /textasciitilde() 
Fe. EO is - 0.036 volt and 
the standard electrode potential for Fe3+ + e- /textasciitilde() 
Fe2+ is 
-uO/textasciitilde()771volr.CalcurareffieEo/textasciitilde()rorFeL+-:j:Le , . '-Fe"".------
[Ans. 
- 0.4395 volt]
\end{problembox}

\begin{problembox}
\textbf{Problem 115} \hfill \textbf{[GRB (O.P. Tandon) p.64 Q97]}
\label{q:ext-grb_all_questions-115}

The standard reduction potential at 25$/circ$C of the reaction 
2H20 + 2e- /textasciitilde() 
Hl + 20W is - 0.8277 va.lt. Calculate the 
equilibrium constant forthe reaction, 
2H20 /textasciitilde() 
H30+ + OW at 25$/circ$C 
[Ans. 
'" 1 014 ] 
[Hint: 
H20 + e- --'-7,YiH2 + OW (Cathode); EO = - 0.8277 volt 
H20 + ,YiH2 -7 HP+ + e-
(Anode);  0 = 0 
EO for the cell = - 0.8277 volt 
Apply now 
EO = 0.0591 [log K ] , . 
(n = 1)] 
. n
\end{problembox}

\begin{problembox}
\textbf{Problem 116} \hfill \textbf{[GRB (O.P. Tandon) p.64 Q98]}
\label{q:ext-grb_all_questions-116}

Calculate the emf of the following cell: 
. Pt(Hz I atm) I CH3CH2COOH (0.15 M) II 0.01 M 
NH40HI H2 (1 atm) Pt 
K" for CH3CH]COOH = 1.4 x 10-5 
K b for NH40H = 1.8 x 10-5 
[Ans. 
- 0.4603 volt] 
[Hint: [W Jin CH3CH2COOH = /textasciitilde()C x Ka = /textasciitilde()0.15 x 1.4 xIO-5 
= 1.449 X 10-3 
[OW] in NH40H = /textasciitilde()C XKb = /textasciitilde()0.01 x1.8 x 10-5 = 0.4242x.iQ-J 
+ . 
10-14 
-II 
[H ]m NH40H = 
. 
, = 2.3573 x 10 
, 
0.4242 xl 0-" 
[H+ lRHs 
Eeell = 0.0591 log -+--] 
[H luiS
\end{problembox}

\begin{problembox}
\textbf{Problem 117} \hfill \textbf{[GRB (O.P. Tandon) p.64 Q99]}
\label{q:ext-grb_all_questions-117}

Calculate equilibrium constantJor 
12 + 1 /textasciitilde() 
I3 
at 298 K from the following information: 
I2(aq.) + 2e- /textasciitilde() 
21-; 
EO = 0.6197 volt 
I; + 2e- /textasciitilde() 
31-; 
[Ans. 
706.9] 
fRint: 
12 + 2e- --721-; 
EO = 0.5355 volt 
31 - --7 13- + 2e - ; EO = - 0.5355 volt 
EO = 0.6197 - 0.5355 
= 0.0842 volt
\end{problembox}

\begin{problembox}
\textbf{Problem 118} \hfill \textbf{[GRB (O.P. Tandon) p.65 Q100]}
\label{q:ext-grb_all_questions-118}

101. 
K 
antilo' [ nEo J = antilo [2 x 0.0842J 
706.9] 
g 0.0591 
g 
0.0591 
A lead storage battery has initially 200 g of lead and 200 g of 
Pb02 , plus excess H2S04, How long could this cell deliver a 
current of 10 amp, without recharging, if it was possible to 
operate it so that the reaction goes to completion? 
[Ans. 
4.48 hour] 
Zn(s) + 2AgCl(s) /textasciitilde() 
ZnC12 (0.555 M) + 2Ag(s) 
Eooc 
1.015 volt (dEl = - 4.02 X 10-4 volt per degree. 
dT IP 
FindI:!.G,I:!.S.
\end{problembox}

\begin{problembox}
\textbf{Problem 119} \hfill \textbf{[GRB (O.P. Tandon) p.65 Q105]}
\label{q:ext-grb_all_questions-119}

What will be the value of' A for /textasciitilde() 0.001 M aqueous NH3 
solution? 
Kb = 1.6 X 10-5 and Ao = 2.38 X 10-2 $/Omega(-1)$ m2 morl 
[Ans. 
2.998 x ro-3 $/Omega$-I cm2 morl ] 
[Hint: 
K b = Ca2 
0.126 
1.6 x 10-5 = 0.001 x a 2 
a = 0.126 
(A"", = molar conductance 
at concentration 'c' and 
A:. :::: molar conductance at 
infinite dilution) 
\underline{\hspace{1cm}} ......... [/textasciitilde()I:!.G .. = .. =.195\underline{\hspace{1cm}}895\underline{\hspace{1cm}}kJ/textasciitilde()M-=-=18.55-cal] ........... -/textasciitilde()---.... ---.. . 
A"", :::: 2.998 ;(10::::3 $/Omega$-I cm2 morl 1 
[Hint: 
We know that, 
Mf=nF[T(/textasciitilde():)p -E] 
2x96500 ['-273 ><4.02 x ro-4 
1.015] 
217075.98 joule = - 217.075 kJ 
I:!.G 
nFE == - 2 x 96500 x 1.015 =-195S95 J 
'., ".:::: -195.895 kJ 
'"'. 
ti..G =Mf 
TliS 
...:.195.895 
217.075 - 273 x /)"S 
I::.S 
- 0.07758 kJ = -77.58 J = -18.55 cal]
\end{problembox}

\begin{problembox}
\textbf{Problem 120} \hfill \textbf{[GRB (O.P. Tandon) p.65 Q102]}
\label{q:ext-grb_all_questions-120}

Calculate the equilibrium constant for the reaction, 
2Fe3+ + 3r /textasciitilde() 
2Fe2+ + 13
\end{problembox}

\begin{problembox}
\textbf{Problem 121} \hfill \textbf{[GRB (O.P. Tandon) p.65 Q103]}
\label{q:ext-grb_all_questions-121}

The standard reduction potentials in acidic conditions are 0.77 
and 0.54 V respectively for Fe3+ I Fe2+ and 13/r couples. 
(liT 19(8) 
[Ans. 
6.26 x 107] 
Electrolysis ofa solution of MnS04 in aqueous sulphuric acid 
is a method for the preparation of Mn02 as the per reaction, 
Mn 2+(aq.) + 2HzO ----t MnOz + 2/textasciitilde()(aq.) + H2 
Passing a current of 27 A for 24 hours gives one kg of Mn0z. 
What is tlje value of current efficiency? Write the reactions 
taking pla'ce at the cathode and at the anode. lilT (May) 19971 
[Hint: 
Apply 
Ex ix t 
w=---
96500 
1000 = 87 x i x 24 x 60 x 60 
2 x 96500 
i :::: 25.6 ampere 
Current efficiency:::: 25.6 x roo = 94.8/% 
27.0 
Reactions: 
Anode: 
Catbode: 
Mn2+ /textasciitilde() 
Mn4+ + 2e-
2H+ + 2e- /textasciitilde()H2]
\end{problembox}

\begin{problembox}
\textbf{Problem 122} \hfill \textbf{[GRB (O.P. Tandon) p.65 Q104]}
\label{q:ext-grb_all_questions-122}

Calculate the number of kWh of electricity is necessary .to 
produce 1 metric ton (1000 kg) of aluminium by Hall process 
in a cel/ operating at 15 V 
[Ans. 
4.47 x 104 kWh]
\end{problembox}

\begin{problembox}
\textbf{Problem 123} \hfill \textbf{[GRB (O.P. Tandon) p.65 Q106]}
\label{q:ext-grb_all_questions-123}

107. 
A weak monobasic acid is 5/% dissociated in 0.01 mol dm-3 
solution. The limiting molar conductivity at infinite dilution is 
4.00 x 10-2 $/Omega$-I m2 mol-I. Calculate the conductivity of a 
0.05 moldriC" solution of'fhe"aci<:[ 
[Ans. 
8.92 x 10-4 $/Omega$-l cm2 morl ] 
[Hint: 
Dissociation constant of a/textasciitilde()id Ka = Ca2 
Ka 
Ca2 
2.5 x 10-5 
0.05 x a 2 
a= 0.0223 
We know that, a 
0.0223 
= 0.01 x (0.05)2 
= 2.5 X 10-5 
A"", 
8.92 X 10-4 $/Omega$-I cm2 mol-I] 
1.05 g of a lead ore containing impurity of Ag was dissolved in 
quantityofHN03 and the volume was made 350 mL. A Ag 
electrode was dipped in the solution and EceU of, 
Pt(H2) I H+(1 M) II Ag+ lAg 
was 0.503 V at 298 K. Calculate /% of lead in the ore. 
EO 
. 
0.80V 
Ag+ I Ag 
[Ans. 
0.0339/%] 
[Hint: 
Anode! H2(g) /textasciitilde() 
H+ (1 M) + e-
2 
Cathode Ag+ (x) + e- /textasciitilde() 
Ag(s) 
Ag(s) + H+ (1 M) 
(n = 1) 
EO = 0.80 - 0 = 0.80 volt 
Q= 
E = EO 
0.0591 1oglO Q 
n 
0.0591 1 
(1) 
0.503 = 0.80 
oglO -
1 
x
\end{problembox}

\begin{problembox}
\textbf{Problem 124} \hfill \textbf{[GRB (O.P. Tandon) p.66 Q108]}
\label{q:ext-grb_all_questions-124}

Calculate EO of the following half-cell reaction at 298 K: 
Ag(NH3); + e- /textasciitilde() 
Ag + 2NH3 
+e- /textasciitilde()Ag; 
Ag(NH3); /textasciitilde()A!t + 2NH3 ; 
[ADS. 
0.373 volt] 
==0.80 V 
K := 6 X 10-8 
Set-1: Questions with single correct answer
\end{problembox}

\begin{problembox}
\textbf{Problem 125} \hfill \textbf{[GRB (O.P. Tandon) p.66 Q1]}
\label{q:ext-grb_all_questions-125}

An electrolyte is a substance which: 
(a) conducts electricity 
(b) decomposes on heating 
(c) is acidic in nature 
(d) when dissolved in water, dissociates into ions
\end{problembox}

\begin{problembox}
\textbf{Problem 126} \hfill \textbf{[GRB (O.P. Tandon) p.66 Q2]}
\label{q:ext-grb_all_questions-126}

The theory of ionisation was presented by: 
(a) Faraday 
(b) Arrhenius (c) Ostwald 
(d) Rutherford
\end{problembox}

\begin{problembox}
\textbf{Problem 127} \hfill \textbf{[GRB (O.P. Tandon) p.66 Q3]}
\label{q:ext-grb_all_questions-127}

Dissociation of an electrolyte in water into ncgative and 
positive ions is called: 
(a) ionisation 
(b) electrolysis 
( c) decomposition 
(d) hydrolysis
\end{problembox}

\begin{problembox}
\textbf{Problem 128} \hfill \textbf{[GRB (O.P. Tandon) p.66 Q4]}
\label{q:ext-grb_all_questions-128}

Degree of ionisation is equal to: 
(a) total number of moles of the electrolyte present in solution 
(b) total number of moles of the electrolyte dissociated into 
ions 
(c) number of moles dissociated Itotal number of moles 
dissolved 
(d) total number of moles dissolved/number of moles dissociated
\end{problembox}

\begin{problembox}
\textbf{Problem 129} \hfill \textbf{[GRB (O.P. Tandon) p.66 Q5]}
\label{q:ext-grb_all_questions-129}

Conductivity of aqueous solution of an electrolyte depends un: 
(a) molecular mass of the electrolyte 
(b) boiling point of solvent 
(c) degree of ionisation 
(d) volume of the solvent
\end{problembox}

\begin{problembox}
\textbf{Problem 130} \hfill \textbf{[GRB (O.P. Tandon) p.66 Q6]}
\label{q:ext-grb_all_questions-130}

Degree ()f ionisation does not depend on: 
(a) nature of the solvent 
(b) nature of the electrolyte 
(c) dilution 
(d) molecular mass of the electrolyte
\end{problembox}

\begin{problembox}
\textbf{Problem 131} \hfill \textbf{[GRB (O.P. Tandon) p.66 Q7]}
\label{q:ext-grb_all_questions-131}

Substances which give good conducting aqueous solution are 
called: 
(a) weak electrolytes 
(b) strong electrolytes 
(c) non-electrolytes 
(d) catalysts
\end{problembox}

\begin{problembox}
\textbf{Problem 132} \hfill \textbf{[GRB (O.P. Tandon) p.66 Q8]}
\label{q:ext-grb_all_questions-132}

The number of ions given by one molecule ofK4Fe(CN)6 after 
complete dissociation is: 
llIint: 
Anode 
EO 
0.80 volt 
Q = [Ag + HNH312 
[Ag(NH3hf 
Ag(s) + Ag+ + 2NH3 
6 X 10-8 
(n 
I) 
E/textasciitilde()ell 
(v - 0,80) 
At equilibrium, E == 0 
.. 
E == EO \underline{\hspace{1cm}} 0.OS91 1og Q 
(a) 5 
11 
o (v:'" 0.80) - 0.OS91 1og (6 x 10-8 ) 
1 
(b) 11 
. (c) 2 
(d) 10
\end{problembox}

\begin{problembox}
\textbf{Problem 133} \hfill \textbf{[GRB (O.P. Tandon) p.66 Q9]}
\label{q:ext-grb_all_questions-133}

The amount of electricity required to produce one mole of 
copper from copper sulphate solution will be .
\end{problembox}

\begin{problembox}
\textbf{Problem 134} \hfill \textbf{[GRB (O.P. Tandon) p.66 Q10]}
\label{q:ext-grb_all_questions-134}

(a) 1 faraday 
( c) 2 faraday 
lHint: Cu 2+ + 2e-
(VITEEE 2008) 
(b) 2.33 faraday 
(d) 1.33 faraday 
Cu (s) 
1 mole copper requires 2 mole electron, i.e .. 2 faraday charge.] 
The process in which chemical change occurs on passing 
electricity is termed: 
(a) ionisation 
(b) neutralisation 
(c) electrolysis 
(d) hydrolysis
\end{problembox}

\begin{problembox}
\textbf{Problem 135} \hfill \textbf{[GRB (O.P. Tandon) p.66 Q11]}
\label{q:ext-grb_all_questions-135}

Which of the following condition is correct for operation of 
electrolytic cell? 
(a) b.G = 0, E , 0 
(c) b.G? 0, E < 0 
(b) ,b.G < 0, E > 0 
(d) b.G > 0, E > 0
\end{problembox}

\begin{problembox}
\textbf{Problem 136} \hfill \textbf{[GRB (O.P. Tandon) p.66 Q12]}
\label{q:ext-grb_all_questions-136}

Which one is the correct equation that represents the first law 
of electrolysis? 
(a) mZ = et 
(b) m 
eZt 
(c) me 
Zt 
(d) e 
mZt 
13.. 
When one coulomb' of electricity is passed through an 
electrolytic solution, the mass deposited on the electrode is 
equal to: 
(a) equivalent weight 
(b) molecular weight 
(c) electrochemical equivalent (d) one gram
\end{problembox}

\begin{problembox}
\textbf{Problem 137} \hfill \textbf{[GRB (O.P. Tandon) p.66 Q14]}
\label{q:ext-grb_all_questions-137}

One faraday is equal to: 
(a) 9650 coulomb 
(c) 19640 coulomb 
(b) 10,000 coulomb 
(d) 96500 coulomb
\end{problembox}

\begin{problembox}
\textbf{Problem 138} \hfill \textbf{[GRB (O.P. Tandon) p.66 Q15]}
\label{q:ext-grb_all_questions-138}

,When one faraday of electric current is passed, the mass 
deposited is equal to: 
(a) one gram equivalent 
(b) one gram mole 
( c) electrochemical equivalent (d) half gram equivalent
\end{problembox}

\begin{problembox}
\textbf{Problem 139} \hfill \textbf{[GRB (O.P. Tandon) p.66 Q16]}
\label{q:ext-grb_all_questions-139}

On passing one faraday of electricity through a dilute solution 
of an acid, the volume of hydrogen obtained at NTP is: 
(a) 22400 mL 
(b) 1120 mL 
(c) 2240 mL 
(d) 11200 mL 
I
\end{problembox}

\begin{problembox}
\textbf{Problem 140} \hfill \textbf{[GRB (O.P. Tandon) p.67 Q17]}
\label{q:ext-grb_all_questions-140}

w g of copper is deposited in a copper voltan1.eter when an 
electric current of 2 ampere is passed for 2 hours. If one 
ampere of electric current is passed for 4 hours in the same 
voltameter, copper deposited will be: 
(a) w 
(b) wl2 
(c) wl4 
(d) 2w
\end{problembox}

\begin{problembox}
\textbf{Problem 141} \hfill \textbf{[GRB (O.P. Tandon) p.67 Q18]}
\label{q:ext-grb_all_questions-141}

Copper sulphate soluticn is electrolysed between two 
platinum electrodes. A current is passed until 1.6 g of oxygen 
is liberated at anode. Thl/textasciitilde() amount of copper deposited at the 
cathode during the same period is: 
(a) 6.36 g 
(b) 63.6 g 
(c) 12.7 g 
(d) 3.2 g
\end{problembox}

\begin{problembox}
\textbf{Problem 142} \hfill \textbf{[GRB (O.P. Tandon) p.67 Q19]}
\label{q:ext-grb_all_questions-142}

When electricity is passed through a solution of AIC!3' 13.5 g
\end{problembox}

\begin{problembox}
\textbf{Problem 143} \hfill \textbf{[GRB (O.P. Tandon) p.67 Q28]}
\label{q:ext-grb_all_questions-143}

The extent of ionisation increases: 
(a) with increase in concentration of the electrolyte 
(b) on decreasing temperature of solution 
(c) on addition of excess of water 
(d) on stirring the solution vigorously
\end{problembox}

\begin{problembox}
\textbf{Problem 144} \hfill \textbf{[GRB (O.P. Tandon) p.67 Q29]}
\label{q:ext-grb_all_questions-144}

Molten sodium chloride conducts electricity due to the 
presence of: 
(a) free electrons 
(b) free ions 
(c) free molecules$/cdot$ 
(d) free atoms ofNa and Cl
\end{problembox}

\begin{problembox}
\textbf{Problem 145} \hfill \textbf{[GRB (O.P. Tandon) p.67 Q30]}
\label{q:ext-grb_all_questions-145}

When NaCi is dissolved in water, the sodium ion is: 
(a) oxidised 
(b) reduced 
(c) hydrolysed 
(d) hydrated 
of Al is discharged. The amount of charge passed is:
\end{problembox}

\begin{problembox}
\textbf{Problem 146} \hfill \textbf{[GRB (O.P. Tandon) p.67 Q31]}
\label{q:ext-grb_all_questions-146}

A solution of sodium sulphate in water is electrolysed using 
/textasciitilde() 
.... --' -.......:...(a)1.5 .. r-1o)---o/textasciitilde()-r--(c)-t:U/textasciitilde()/textasciitilde()(G)---r.OF-- ...... -
m------ptatinum-etectrod'es-:-The proauctsafCathooe and' anode are 
20.' When the same electrie current is passed through the solution 
of different electrolytes in series, the amounts of elements 
deposited on the electrodes are in the ratio oftheir: 
(a) atomic numbers 
(b) atomic masses 
. \underline{\hspace{1cm}} . 
\underline{\hspace{1cm}} /textasciitilde()\underline{\hspace{1cm}}/textasciitilde()\underline{\hspace{1cm}}M \underline{\hspace{1cm}} $/cdot$ \underline{\hspace{1cm}} $/cdot$/textasciitilde() \underline{\hspace{1cm}} \underline{\hspace{1cm}} . \underline{\hspace{1cm}} . \underline{\hspace{1cm}} ... \underline{\hspace{1cm}} 
.. /textasciitilde() \underline{\hspace{1cm}} /textasciitilde()\underline{\hspace{1cm}},\underline{\hspace{1cm}}. 
\underline{\hspace{1cm}} .. 
(c) specific gravities 
(d) equivalent masses
\end{problembox}

\begin{problembox}
\textbf{Problem 147} \hfill \textbf{[GRB (O.P. Tandon) p.67 Q21]}
\label{q:ext-grb_all_questions-147}

Faraday'S laws of electrolysis are related to: \underline{\hspace{1cm}} 
(a) atomic number of the cation 
(b) atomic number of the anion 
(c), equivalent mass of the electrolyte 
(d) speed of the cation
\end{problembox}

\begin{problembox}
\textbf{Problem 148} \hfill \textbf{[GRB (O.P. Tandon) p.67 Q22]}
\label{q:ext-grb_all_questions-148}

The specific conductance of a 0.01 Iv! solution of KCl is 
0.0014 $/Omega$'" I cm- I at 25$/circ$C. Its equivalent conductance is: 
(a) 14 
(b) 140' 
(c) 1.4 
Cd) 0.14
\end{problembox}

\begin{problembox}
\textbf{Problem 149} \hfill \textbf{[GRB (O.P. Tandon) p.67 Q23]}
\label{q:ext-grb_all_questions-149}

The equivalent conductivity of 0.1 N CH3COOH at 25$/circ$C is 80 
and at infinite dilution 400 $/Omega$-I. The degree of dissociation 
of CH3COOH is: 
(a) I 
(b) 0.2 
(c) 0.1 
(d) 0.5
\end{problembox}

\begin{problembox}
\textbf{Problem 150} \hfill \textbf{[GRB (O.P. Tandon) p.67 Q24]}
\label{q:ext-grb_all_questions-150}

The number of electrons involved when one faraday of 
electricity is passed through an electrolytic solution is: 
(a) 96500 
(b) 8 x 106 
(c) 12 X 1016 (d) 6 x 1023
\end{problembox}

\begin{problembox}
\textbf{Problem 151} \hfill \textbf{[GRB (O.P. Tandon) p.67 Q25]}
\label{q:ext-grb_all_questions-151}

One faraday of charge was passed through the electrolytic 
cells placed in series containing solutions of Ag+, Ni2+ and 
Cr3+ respectively. The amount of Ag (At. mass lOS), Ni (At. 
mass 59) and Cr (At. mass 52) deposited will be: 
Ag 
Ni 
Cr 
(a) 10Sg 
29.5g 
17.5g 
(b) 108 g 
59.0 g 
52.0 g 
(c) 108g 
IOS.0g 
108.0g 
(d) 108g 
117.5g 
166.0g
\end{problembox}

\begin{problembox}
\textbf{Problem 152} \hfill \textbf{[GRB (O.P. Tandon) p.67 Q26]}
\label{q:ext-grb_all_questions-152}

One faraday of electricity will liberate one gram mole of the 
metal from the solution of: 
(a) BaCl2 
(b) CuS04 
(c) AICl3 
(d) NaCl
\end{problembox}

\begin{problembox}
\textbf{Problem 153} \hfill \textbf{[GRB (O.P. Tandon) p.67 Q27]}
\label{q:ext-grb_all_questions-153}

Strong electrolytes are those which: 
(a) dissolve readily in water 
(b) conduct electricity 
(I;) dissociate into ions even at high concentration 
(d) dissociate into ions at high dilution 
respectively: 
(a) H2,02 
(c) 02' Na 
(b) 02' H2 
(d) 02' S02
\end{problembox}

\begin{problembox}
\textbf{Problem 154} \hfill \textbf{[GRB (O.P. Tandon) p.67 Q32]}
\label{q:ext-grb_all_questions-154}

The electric conduction of a salt solution in waterdependseOIl 
$/cdot$ .. the:-
/textasciitilde()'-"'-""-'/textasciitilde()"-'c"-."-"-. - ..... \underline{\hspace{1cm}}-\underline{\hspace{1cm}} .. \underline{\hspace{1cm}} ... 
(a) shape of molecules 
(b) size of its molecules 
(c) size of solvent molecules (d) extent of its ionisation
\end{problembox}

\begin{problembox}
\textbf{Problem 155} \hfill \textbf{[GRB (O.P. Tandon) p.67 Q33]}
\label{q:ext-grb_all_questions-155}

In electroplating, the article to be electroplated - /textasciitilde()rves as: 
(a) cathode 
(b) t.-" .. ".WIYlt: 
(c) anode 
Cd) coilducf6r
\end{problembox}

\begin{problembox}
\textbf{Problem 156} \hfill \textbf{[GRB (O.P. Tandon) p.67 Q34]}
\label{q:ext-grb_all_questions-156}

The amount of electricity that can deposit lOS g of silver from 
silver nitrate solution is: 
(AFMC 1993) 
(a) I ampere 
(b) I coulomb 
(c) I faraday 
(d) 2 ampere
\end{problembox}

\begin{problembox}
\textbf{Problem 157} \hfill \textbf{[GRB (O.P. Tandon) p.67 Q35]}
\label{q:ext-grb_all_questions-157}

A certain current liberated 0.504 g of hydrogen in 2 hours. 
How many grams of copper can be liberated by the same 
current flowing for the same time in CuS04 solution? 
Ca) 12.7g 
.(b) 15.9g 
(c) 31.Sg 
(d) 63.5g
\end{problembox}

\begin{problembox}
\textbf{Problem 158} \hfill \textbf{[GRB (O.P. Tandon) p.67 Q36]}
\label{q:ext-grb_all_questions-158}

If the specific resistance of a solution of concentration C g 
equivalent litre -I is R, then its equivalent conductance is: 
(a) 100R' 
(b) RC 
(c) 1000 
(d) 
C 
C 
1000 
RC 
1000R
\end{problembox}

\begin{problembox}
\textbf{Problem 159} \hfill \textbf{[GRB (O.P. Tandon) p.67 Q37]}
\label{q:ext-grb_all_questions-159}

If the specific conductance and conductance of a solution are 
same, then its cell constant is equal to: 
(a) I 
(b) 0 
(c) 10 
(d) 100
\end{problembox}

\begin{problembox}
\textbf{Problem 160} \hfill \textbf{[GRB (O.P. Tandon) p.67 Q38]}
\label{q:ext-grb_all_questions-160}

On increasing the dilution, the specific conductance: 
(Jiwaji 1990) 
(a) increases 
(b) decreases 
(c) remains constant 
(d) none of these
\end{problembox}

\begin{problembox}
\textbf{Problem 161} \hfill \textbf{[GRB (O.P. Tandon) p.67 Q39]}
\label{q:ext-grb_all_questions-161}

The distance between two electrodes of a cell is 2.5 cm and 
area of each electrode is 5 cm2 $/bullet$ The cell constant is: 
(Jabalpur 1990) 
(a) 2 
(b) 12.5 
(c) 7.5 
(d) 0.5
\end{problembox}

\begin{problembox}
\textbf{Problem 162} \hfill \textbf{[GRB (O.P. Tandon) p.67 Q40]}
\label{q:ext-grb_all_questions-162}

At 25$/circ$C, the molar conductances at infinite dilution for the 
strong electrolytes NaOH, NaCI and BaCI2 are 24S x 10-4, 
126 X 10-4 and 2S0x 10-4 Sm2 mol-i respectively. t//textasciitilde()nBa(OHh 
in S m2 marl: 
(EAMCET 20(9) 
(a) 52.4 x 10-4 
(b) 524 x 10-4 
(c) 402 x 10-4 
(d) 262x ]0-4 
[Hint : A/textasciitilde(),Ba(OHh = A:IIBaCI2 + 2AomNaOH 
2A/textasciitilde(),NaCI 
= 2S0x/textasciitilde()0-4 + 2x 248x 10-4 - 2 x 126x 10-4 
= 524 X 1O-4S m2 morl]
\end{problembox}

\begin{problembox}
\textbf{Problem 163} \hfill \textbf{[GRB (O.P. Tandon) p.68 Q41]}
\label{q:ext-grb_all_questions-163}

The electrochemical cell stops working after some time 
because: 
(VITEEE 2008) 
(a) electrode potential of both the electrodes becomes zero 
(b) electrode potential of both the electrodes becomes equal 
(c) one of the electrodes is eaten away 
(d) the cell reaction gets reversed
\end{problembox}

\begin{problembox}
\textbf{Problem 164} \hfill \textbf{[GRB (O.P. Tandon) p.68 Q42]}
\label{q:ext-grb_all_questions-164}

Which reaction will take place at cathode when fused calcium 
chloride is electrolysed? 
(a) Ca2+ + 2e- /textasciitilde() 
Ca 
(b) Cl- /textasciitilde() 
CI + e-
\end{problembox}

\begin{problembox}
\textbf{Problem 165} \hfill \textbf{[GRB (O.P. Tandon) p.68 Q43]}
\label{q:ext-grb_all_questions-165}

In electrolysis oxidation takes place at: 
(a) both the electrodes(J:  cathode 
(c) anode 
(d) in the solution
\end{problembox}

\begin{problembox}
\textbf{Problem 166} \hfill \textbf{[GRB (O.P. Tandon) p.68 Q44]}
\label{q:ext-grb_all_questions-166}

The equation representing the process by which standard 
reduction potential of zinc can be defined is: 
(a) Zn 2+(s) + 2e- /textasciitilde() 
Zn 
(5)zii(gr/textasciitilde()$/cdot$ zii2'f(g) + 
(c) Zn 2+(g) + 2e- /textasciitilde() 
Zn 
(d) Zn 2+(aq.) + 2e- /textasciitilde() 
Zn(s)
\end{problembox}

\begin{problembox}
\textbf{Problem 167} \hfill \textbf{[GRB (O.P. Tandon) p.68 Q45]}
\label{q:ext-grb_all_questions-167}

The measured potential for, 
Mg2+ + 2e- /textasciitilde() 
Mg(s) 
does not depend upon: 
(a) raising the temperature 
(b) increasing the concentration of Mg2+ ions 
(c) making the magnesium plate bigger 
(d) purity of magnesium plate
\end{problembox}

\begin{problembox}
\textbf{Problem 168} \hfill \textbf{[GRB (O.P. Tandon) p.68 Q46]}
\label{q:ext-grb_all_questions-168}

All cells do not contain: 
(a) an anode 
(b) a cathode 
(c) ions 
(d) a porous partition
\end{problembox}

\begin{problembox}
\textbf{Problem 169} \hfill \textbf{[GRB (O.P. Tandon) p.68 Q47]}
\label{q:ext-grb_all_questions-169}

When lead accumulator is charged, it is: 
(a) an eleCtrolytic cell 
(b) a galvanic cell 
(c) a Daniell cell 
(d) none of these
\end{problembox}

\begin{problembox}
\textbf{Problem 170} \hfill \textbf{[GRB (O.P. Tandon) p.68 Q48]}
\label{q:ext-grb_all_questions-170}

The standard electrode potentials for the reactions, 
Ag+(aq.) + e- ---7 Ag(s) 
Sn 2+(aq.) + 2e- ---7 Sn(s) 
at 25$/circ$C are 0.80 volt and - 0.14 volt respectively. The emf of 
the cell, 
Sn I Sn 2+(1 M) II Ag+(1 M) I Ag 
is: 
(a) 0.66 volt (b) 0.80 volt (c) 1.08 volt (d) 0.94 volt
\end{problembox}

\begin{problembox}
\textbf{Problem 171} \hfill \textbf{[GRB (O.P. Tandon) p.68 Q49]}
\label{q:ext-grb_all_questions-171}

The cathodic reaction in electrolysis of dilute H2S04 with 
platinum electrode is: 
(a) oxidation 
(b) reduction 
(c) both oxidation and reduction 
(d) neutralization
\end{problembox}

\begin{problembox}
\textbf{Problem 172} \hfill \textbf{[GRB (O.P. Tandon) p.68 Q50]}
\label{q:ext-grb_all_questions-172}

Strongest reducing agent is: 
(a) K 
(b) Mg 
(c) Al 
(d) I
\end{problembox}

\begin{problembox}
\textbf{Problem 173} \hfill \textbf{[GRB (O.P. Tandon) p.68 Q51]}
\label{q:ext-grb_all_questions-173}

The metal oxide which decomposes on heating, is: 
(a) ZnO 
(b) HgO 
(c) A120 3 
(d) CuO
\end{problembox}

\begin{problembox}
\textbf{Problem 174} \hfill \textbf{[GRB (O.P. Tandon) p.68 Q52]}
\label{q:ext-grb_all_questions-174}

Thereaction,/textasciitilde()H2(g) + AgCI(s) = H +(aq.) + Cnaq.) + Ag(s) 
occurs in the galvanic cell: 
(AFMC 2009) 
(a) 
I AgCl(s) I KCI (soln.) II AgN03 (soln.) lAg 
(b) Pt I H2(g) I HCl (soln.) II AgN03 (soln.) I Ag 
(c) 'Pt I H2(g) I HCl (soln.) II AgCl(s) I Ag 
(d) Pt I H2(g) I KCI (soln.) II AgCI(s) I Ag
\end{problembox}

\begin{problembox}
\textbf{Problem 175} \hfill \textbf{[GRB (O.P. Tandon) p.68 Q53]}
\label{q:ext-grb_all_questions-175}

The standard oxidation potentials, EO ,. for the half reactions 
are as, 
Zn /textasciitilde() 
Zn 2+ + 2e-; 
EO = + 0.76 volt 
Pe /textasciitilde() 
Pe2+ + 2e-; 
EO = + 0.41 volt 
lheemfofthecen/textasciitilde()/textasciitilde()+-+-Zn /textasciitilde() 
Zn2+-:j:-Pe IS: 
ICET (Karnataka) 2009] 
(a) + 0.35 volt 
(b) - 0.35 volt 
(c) +1.17 volt 
(d) -1.17 volt
\end{problembox}

\begin{problembox}
\textbf{Problem 176} \hfill \textbf{[GRB (O.P. Tandon) p.68 Q54]}
\label{q:ext-grb_all_questions-176}

The. standard/textasciitilde()reduction-potentials/textasciitilde()at.25':CJorthe.fonowm/textasciitilde()--
half reactions are given against each; 
. 
. . 
55:
\end{problembox}

\begin{problembox}
\textbf{Problem 177} \hfill \textbf{[GRB (O.P. Tandon) p.68 Q58]}
\label{q:ext-grb_all_questions-177}

59. 
Zn 2+(aq.) + 2e- /textasciitilde() 
Zn(s); 
Cr3+(aq.)+ 3e- /textasciitilde() 
Cr(s); 
2W + 2e-
H/textasciitilde()(g); 
Fe3+ + e- /textasciitilde() 
Fe2+; 
Which is the strongest reducing agent? 
(a) Zn 
(b) Cr 
(c) H2(g) 
(d) Pe2+(aq.) 
Hydrogen gas will not reduce heated: 
- 0.762 
- 0.740 
0.00 
0.77 
(VITEEE 2007) 
(a) cupric oxide 
(b) ferric oxide 
(c) stannic oxide 
(d) aluminium oxide 
A solution containing one mole per litre of each Cu(N03h; 
AgN03; Hg2(N03h and Mg(N03)z is being electrolysed by 
using inert electrodes. The values of standard electrode 
potentials (reduction potentials) are Ag/Ag+ 
0.80 volt, 
2HglHi+ = 0.79 volt, Cu/Cu 2+= "" 0.24 volt, Mg/Mg/textasciitilde()+ 
= - 2.37 volt. With increasing voltage, the sequence of 
deposition of metals on the cathode will be: 
[PMT (Kerala) 2004] 
(a) Ag, Hg, Cu 
(b) Cu, Hg, Ag 
(c) Ag, Hg, Cu, Mg 
(d) Mg, Cu, Hg, Ag 
Four colourless salt solutions are placed in separate test tubes 
and a strip of copper is dipped in each. Which solution finally 
turns blue? 
(a) Pb(N03h (b) AgN03 
(c) Zn(N03h (d) Cd(N03h 
Red hot carbon will remove oxygen from the oxide XO and YO 
but not from Zo. Y will remove oxygen from Xo. Use this 
evidence to deduce the order of acti'(ity of the three metals X , 
Y and Z putting the most active first: 
(a) XYZ 
(b) ZIT 
(c) YXZ 
(d) ZXY 
Which of the following metals does not give the following 
reaction? 
M + Water ---7 Oxide or hydroxide + H2 
(a) Iron 
(b) S.odium 
(c) Mercury (d) Magnesium
\end{problembox}

\begin{problembox}
\textbf{Problem 178} \hfill \textbf{[GRB (O.P. Tandon) p.69 Q60]}
\label{q:ext-grb_all_questions-178}

Which is the best reducing agent? 
(a) F 
(b) Cl-
(c) Br-
(d) 1-
\end{problembox}

\begin{problembox}
\textbf{Problem 179} \hfill \textbf{[GRB (O.P. Tandon) p.69 Q61]}
\label{q:ext-grb_all_questions-179}

If a spoon of copper metal is placed in a solution of ferrous 
sulphate: 
(a) Cu will precipitate out 
(b) iron will precipitate 
(c) Cu and Fe will precipitate 
(d) no reaction wiil take place
\end{problembox}

\begin{problembox}
\textbf{Problem 180} \hfill \textbf{[GRB (O.P. Tandon) p.69 Q62]}
\label{q:ext-grb_all_questions-180}

Among Na, Hg, S, Pt and graphite, which can be used as 
electrodes in electrolytic cells having aqueous solutions? 
(a) Hg and Pt 
(b) Hg, Pt and graphite 
and S
\end{problembox}

\begin{problembox}
\textbf{Problem 181} \hfill \textbf{[GRB (O.P. Tandon) p.69 Q63]}
\label{q:ext-grb_all_questions-181}

The most reactive metal among the following is: 
(a) AI 
(b) Ni 
(c) Pb 
Cd) Cu
\end{problembox}

\begin{problembox}
\textbf{Problem 182} \hfill \textbf{[GRB (O.P. Tandon) p.69 Q64]}
\label{q:ext-grb_all_questions-182}

Which of the following metals is most readily corroded in 
moist air? 
/textasciitilde()\underline{\hspace{1cm}}.- ---ta)-C-opper- '---/textasciitilde()--/textasciitilde()------Eb)--Iron --------------
(c) Silver 
(d) Nickel
\end{problembox}

\begin{problembox}
\textbf{Problem 183} \hfill \textbf{[GRB (O.P. Tandon) p.69 Q65]}
\label{q:ext-grb_all_questions-183}

Which one will liberate Br2 from KBr? 
(a) HI 
(b) 12 
(c) Clz 
(d) SOz
\end{problembox}

\begin{problembox}
\textbf{Problem 184} \hfill \textbf{[GRB (O.P. Tandon) p.69 Q66]}
\label{q:ext-grb_all_questions-184}

Which cine of the following is not the correct representation? 
(a) 
Red. pot. of cathode + Oxid. pot. of anode 
(b) E/textasciitilde()ell :::: Red. pot. of cathode - Oxid. pot. of anode 
(c) E/textasciitilde()ell :::: Red. pot. of cathode 
Red. pot. of anode 
(d) E/textasciitilde()ell 
Oxid. pot. of cathode + Oxid. pot. of anode
\end{problembox}

\begin{problembox}
\textbf{Problem 185} \hfill \textbf{[GRB (O.P. Tandon) p.69 Q67]}
\label{q:ext-grb_all_questions-185}

Which of the following represents the potential of silver wire 
dipped into 0.1 M AgN03 solution at 25$/circ$C? 
(a) 
(b) (E;oo. + 0.059) 
(c) (E/textasciitilde()xid. - 0.059) 
(d) 
0.059)
\end{problembox}

\begin{problembox}
\textbf{Problem 186} \hfill \textbf{[GRB (O.P. Tandon) p.69 Q68]}
\label{q:ext-grb_all_questions-186}

If the solution ofthe CUS04 in which copper rod is immersed 
is diluted to lO'times, the electrode potential: 
(a) increases by 0.030 volt 
(b) decreases by 0.030 volt 
(c) increases by 0.059 volt 
(d) decreases by 0.0059 volt
\end{problembox}

\begin{problembox}
\textbf{Problem 187} \hfill \textbf{[GRB (O.P. Tandon) p.69 Q69]}
\label{q:ext-grb_all_questions-187}

A solution of Cu(II) sulphate is reacted with KCI and Kl. In 
. which case will the Cu 2+ be reduced to Cu +? 
(a) In both the cases 
(b) When reacted with KCI 
(c) When reacted with KI 
(d) In both the case/textasciitilde() but in presence of H +
\end{problembox}

\begin{problembox}
\textbf{Problem 188} \hfill \textbf{[GRB (O.P. Tandon) p.69 Q70]}
\label{q:ext-grb_all_questions-188}

From the electrochemical series, it can be concluded that: 
(a) Znz+ will liberateH2 from I MHCl 
. 
/textasciitilde()b) Ag metal reacts spontaneously with Zn2+ 
(c) Zn metal will liberate H2 from I MHCl 
(d) Ag metal will liberate H2 from I M HCI
\end{problembox}

\begin{problembox}
\textbf{Problem 189} \hfill \textbf{[GRB (O.P. Tandon) p.69 Q71]}
\label{q:ext-grb_all_questions-189}

The potential of a hydrogen electrode at pH =: I is: 
. (a) 0.059 volt 
(b) 0 volt 
(c) 
0.059 volt 
(d) 0.59 volt
\end{problembox}

\begin{problembox}
\textbf{Problem 190} \hfill \textbf{[GRB (O.P. Tandon) p.69 Q72]}
\label{q:ext-grb_all_questions-190}

. Which is not true for a standard hydrogen electrode? 
(a) The hydrogen ion concentration is I M 
(b) Temperature is 25$/circ$C 
(c) Pressure of hydrogen is I atmosphere 
(d) It contains a metallic conductor which does not absorb 
. hydrogen
\end{problembox}

\begin{problembox}
\textbf{Problem 191} \hfill \textbf{[GRB (O.P. Tandon) p.69 Q73]}
\label{q:ext-grb_all_questions-191}

The electrode potential becomes equal to standard electrode 
potential when reactants and products ratio is: 
(a) equal to I 
(b) greater than I 
(c) less than I 
(d) none of these
\end{problembox}

\begin{problembox}
\textbf{Problem 192} \hfill \textbf{[GRB (O.P. Tandon) p.69 Q74]}
\label{q:ext-grb_all_questions-192}

For the half-cell reaction, 
Au 3+ + 3e-
Au 
the villue of n used in Nernst equation is: 
(a) 3 
(b) 2 
(c) I 
(d) 3 x 96500 
When a piece of sodium metal 
in water, a ''''''''''IU 
takes place to yield hydrogen because: 
(a) sodium loses electrons 
(b) sodium acts as an oxidising agent 
(c).\underline{\hspace{1cm}}wate.r .loses electron/textasciitilde()c 
(d) water acts as a reducing agent
\end{problembox}

\begin{problembox}
\textbf{Problem 193} \hfill \textbf{[GRB (O.P. Tandon) p.69 Q76]}
\label{q:ext-grb_all_questions-193}

Which one is the wrong statement about$/cdot$ electrochemical 
series? 
(a) Active metals have negative reduction potentials 
(b) Active non/textasciitilde()metals have positive reduction potentials 
(c) Metals above hydrogen liberate hydrogen from acids 
(d) Metals below hydrogen are strong reducing agents
\end{problembox}

\begin{problembox}
\textbf{Problem 194} \hfill \textbf{[GRB (O.P. Tandon) p.69 Q77]}
\label{q:ext-grb_all_questions-194}

The reduction potential values are given below: 
3+' 
2 
Al 1 Al =: -1.67 volt, 
Mg + IMg =: 
2.34 volt, 
$/cdot$Cu 2+ICu :::: + 0.34 volt, 
12/21 
+ 0.53 volt. 
Which one is the best reducing agent? 
(a) Al 
(b) Mg 
(c) Cu 
.(d) 12
\end{problembox}

\begin{problembox}
\textbf{Problem 195} \hfill \textbf{[GRB (O.P. Tandon) p.69 Q78]}
\label{q:ext-grb_all_questions-195}

From the values given in question No. 77, which one is the 
best oxidising agent? 
(a) Al 
(b) Mg 
(c) 12 
(d) Cu
\end{problembox}

\begin{problembox}
\textbf{Problem 196} \hfill \textbf{[GRB (O.P. Tandon) p.69 Q79]}
\label{q:ext-grb_all_questions-196}

When iron is rusted, it is: 
(a) reduced 
(b) oxidised 
(c) evaporated 
(d) decomposed
\end{problembox}

\begin{problembox}
\textbf{Problem 197} \hfill \textbf{[GRB (O.P. Tandon) p.69 Q80]}
\label{q:ext-grb_all_questions-197}

Galvanization of iron denotes coating with: 
(a). Cu 
.(b) Sn 
(c) Zn 
(d) Al
\end{problembox}

\begin{problembox}
\textbf{Problem 198} \hfill \textbf{[GRB (O.P. Tandon) p.69 Q81]}
\label{q:ext-grb_all_questions-198}

The standard electrode potentials 9f four elements A, B, C and 
Dare -3.05, 1.66, ---0.40 and 0.80 volt. The highest chemical 
activity will be shown by: 
(a) A 
(b) B 
(e) C 
(d) D
\end{problembox}

\begin{problembox}
\textbf{Problem 199} \hfill \textbf{[GRB (O.P. Tandon) p.69 Q82]}
\label{q:ext-grb_all_questions-199}

Which of the following methods does not liberate hydrogen? 
(a) Zn + H2S04 (dit.) 
(b) Mg + H2S04 (dit.) 
(c) Cu + H2S04 (dB.) 
(d) Zn + HCl (dil.)
\end{problembox}

\begin{problembox}
\textbf{Problem 200} \hfill \textbf{[GRB (O.P. Tandon) p.69 Q83]}
\label{q:ext-grb_all_questions-200}

A depolarizer used in dry cell is: 
(a) ammonium chloride 
(b) manganese dioxide 
(c) potassium oxide 
(d) sodium phosphate 
84 .. The oxide which can be reduced by hydrogen is: 
(a) Na20$/cdot$ 
(b) CaO 
(c) K20 
(d) CuO
\end{problembox}

\begin{problembox}
\textbf{Problem 201} \hfill \textbf{[GRB (O.P. Tandon) p.69 Q85]}
\label{q:ext-grb_all_questions-201}

The reference electrode is made from which of the following? 
(a) ZnCI2 
(b) CuS04 
(c) Hg2Cl2 
(d) HgCl2
\end{problembox}

\begin{problembox}
\textbf{Problem 202} \hfill \textbf{[GRB (O.P. Tandon) p.70 Q86]}
\label{q:ext-grb_all_questions-202}

Given, standard electrode potentials; 
Fe3+ + 3e- /textasciitilde() 
Fe; 
EO = - 0.036 volt 
Fe2+ + 2e - /textasciitilde() 
Fe; 
EO = - 0.440 volt 
The standard electrode potential EO for Fe3+ + e- /textasciitilde() 
Fe2+ 
is: 
IA/'fU (Medical) 2007; AIEEE 2009) 
(a) 
0.476 volt 
(b) 
0.404 volt 
(c) 0.440 volt 
(d) 
0.772 volt 
[Hint: Apply !!"G = - nFE 
Fe2+ + 2e- /textasciitilde() 
Fe !!"G = 
2 x F x (- 0.440 V) = O.8SF ... (i) 
+ 3e- /textasciitilde() 
Fe !!"G = 
3 x F x (-0.036) = O.lOSF 
... (ii) 
Subtracting eq. (i) from eq. (ii), 
. /textasciitilde()/textasciitilde()---"Fe?"+-e- /textasciitilde() 
Fer+;/textasciitilde()'KG;;;O.lOSF"": O.SSOF--
O.772F 
E" for the reaction = - !!"G = c-:- 0.772F) = 0.772 volt] 
nF 
1 xF 
81 .. .K;. Caand-Li-metals-maybe$/cdot$arrangedin-thedecreasing$/cdot$order$/cdot$ 
oftheir standard electrode potentials as: 
. 
(a) K, Ca and Li 
(b) Li, K and Ca 
(c) Li, Ca and K 
(d) Ca, Li and K
\end{problembox}

\begin{problembox}
\textbf{Problem 203} \hfill \textbf{[GRB (O.P. Tandon) p.70 Q88]}
\label{q:ext-grb_all_questions-203}

In a galvanic cell energy changes occur like: . 
(a) chemical energy /textasciitilde() 
electrical energy 
(b) electrical energy /textasciitilde() 
chemical energy 
(c) chemical energy /textasciitilde() 
internal energy 
(d) internal energy /textasciitilde() 
electrical energy
\end{problembox}

\begin{problembox}
\textbf{Problem 204} \hfill \textbf{[GRB (O.P. Tandon) p.70 Q89]}
\label{q:ext-grb_all_questions-204}

The reaction is spontaneous if the cell potential is: 
(a) positive 
(b) negative 
(c) zero 
(d) infinite
\end{problembox}

\begin{problembox}
\textbf{Problem 205} \hfill \textbf{[GRB (O.P. Tandon) p.70 Q90]}
\label{q:ext-grb_all_questions-205}

Is the reaction, 2Al + 3Fe2+ 
2Al3+ + 2Fe possible? 
(a) No, because standard oxidation potential of Al < Fe 
(b) Yes, because standard oxidation potential of Al > Fe 
(c) Cannot be predicted 
(d) Yes, because aluminium is a strong oxidising agent
\end{problembox}

\begin{problembox}
\textbf{Problem 206} \hfill \textbf{[GRB (O.P. Tandon) p.70 Q91]}
\label{q:ext-grb_all_questions-206}

In an experimental set-up for the measurement of emf of a 
half-cell using a reference electrode and a salt bridge, when 
the salt bridge is removed, the voltage: 
(a) does not change 
(b) increases to maximum 
(c) decreases half the value (d) drops to zero
\end{problembox}

\begin{problembox}
\textbf{Problem 207} \hfill \textbf{[GRB (O.P. Tandon) p.70 Q92]}
\label{q:ext-grb_all_questions-207}

More electronegative elements have: 
(a) negative reduction potential 
(b) tendency to lose elections 
(c) pos/textasciitilde()tiye reduction potential 
(d) positive oxidation potential
\end{problembox}

\begin{problembox}
\textbf{Problem 208} \hfill \textbf{[GRB (O.P. Tandon) p.70 Q93]}
\label{q:ext-grb_all_questions-208}

Which metal will be depos*d in the galvanic cell? . 
Cu I Cu 2+ II 
I Ag 
(a) Cu 
. (b) Ag 
(c) Both 
.' (d) None of these
\end{problembox}

\begin{problembox}
\textbf{Problem 209} \hfill \textbf{[GRB (O.P. Tandon) p.70 Q94]}
\label{q:ext-grb_all_questions-209}

In the cell Zn I /textasciitilde()n 2+ II CU 2+ I Cu, the negative electrode is: 
(a) eu . 
(b) Cu 2+ 
(c) Zn 
(d) Zn 2+
\end{problembox}

\begin{problembox}
\textbf{Problem 210} \hfill \textbf{[GRB (O.P. Tandon) p.70 Q95]}
\label{q:ext-grb_all_questions-210}

Which of the following gains electrons more easily? 
(a) Na+ 
(b) Zn 2+ 
(c) A13+ 
(d) H+
\end{problembox}

\begin{problembox}
\textbf{Problem 211} \hfill \textbf{[GRB (O.P. Tandon) p.70 Q96]}
\label{q:ext-grb_all_questions-211}

The potential of a hydrogen electrode at pH lOis: 
(a) 0.51 volt 
(b) 0 volt 
(c) - 0.591 volt 
(d) 0.059 volt
\end{problembox}

\begin{problembox}
\textbf{Problem 212} \hfill \textbf{[GRB (O.P. Tandon) p.70 Q97]}
\label{q:ext-grb_all_questions-212}

The emf ofthe cell, . 
Ni I Ni2+ (1.0 M) II Ag+ (1.0 M) lAg 
0.25 volt, EO for Ag+/Ag = 0.80 volt) 
is given by: 
(a) - 0.25 + 0.80 == 0.55 volt 
(b) - 0.25 - (+ 0.80) 
1.05 volt 
(c) 0 + 0.80 - (- 0.25) 
+ 1.05 volt 
(d) .- 0.80 -
(/textasciitilde() 0.25) 
0.55 volt 
..... /textasciitilde()98 $/bullet$. :]::he.reacti()nf0f-the.celI, 
Zn I Zn 2+ (1.0M) II Cd2+ (1.0 M) I Cd 
is: 
(a) Cd /textasciitilde() 
Cd2+ + 2e-
.(b)Zn2:':/textasciitilde()Zn ... .2e= 
(c) Cd+ Zn2+ /textasciitilde() 
Cd2+ + Zn 
(d) Zn + Cd2+ /textasciitilde() 
Cd + Zn 2+
\end{problembox}

\begin{problembox}
\textbf{Problem 213} \hfill \textbf{[GRB (O.P. Tandon) p.70 Q99]}
\label{q:ext-grb_all_questions-213}

The position of some metals in the electrochemical series in 
decreasing electropositive character is givenas Mg > Al > Zn. 
> Cu > Ag. What will happen if a copper spoon is used to stir a 
solution of aluminium 'nitrate? 
(a) The spoon will get coated with aluminium 
(b) An alloy of aluminium and copper is formed 
(c) The solution becomes blue 
(d) There is no reaction
\end{problembox}

\begin{problembox}
\textbf{Problem 214} \hfill \textbf{[GRB (O.P. Tandon) p.70 Q100]}
\label{q:ext-grb_all_questions-214}

The half-cell reaction is the one that: 
. 
(a) takes place atone electro'de 
(b) consumes half a unit of electricity 
(c) . involves half a mole of electrolyte 
(d) goes halfway to completion 
Which Nernst equation is true to fmd out the potential of 
non-standard electrochemical cell from the following? . 
Fe(s) 
(xM)11 I-(aq$/cdot$)1 12 (s)(Pt) 
ICE:]:: (Gujarat) 2008J 
. 
0 
0.592 
2+]2 
(a) Eeel! = Eeel! -
loglo (Fe ](i 
n 
. 
(b) E 
= E" - 0.0592 I 
[F 2+] [1-]2 
cell 
eel! 
og 10 
e 
. 
n 
$/circ$ 
0.0592 I 
[2+ 
(C) E ell = EeeH -
oglo Fe ](1] 
e 
n 
o 
0.0592 I 
[Fe2+](if 
(d) Eeel! == Ecell 
nF 
oglo 
[Fe](I2]
\end{problembox}

\begin{problembox}
\textbf{Problem 215} \hfill \textbf{[GRB (O.P. Tandon) p.70 Q102]}
\label{q:ext-grb_all_questions-215}

During the electrolysis of fused NaC!, which reaction occurs 
at anode? 
(a) Chloride ions are oxidised 
(b). Sodium ions are oxidised 
(c) Chloride ions are reduced 
(d) Sodium ions are re/textasciitilde()uced
\end{problembox}

\begin{problembox}
\textbf{Problem 216} \hfill \textbf{[GRB (O.P. Tandon) p.70 Q103]}
\label{q:ext-grb_all_questions-216}

 0 (Ni2+ INi) == - 0.25 volt, EO(Au 3+ I Au) = 1.50 volt. The 
emf of the voltaic cell,
\end{problembox}

\begin{problembox}
\textbf{Problem 217} \hfill \textbf{[GRB (O.P. Tandon) p.71 Q104]}
\label{q:ext-grb_all_questions-217}

Which of the following does not occur at cathode? 
(a) Ag+ --7Ag e-
(b) Fez+ --7 Fe3+ + e-
(c) Cu 2+ + 2e- --7 Cu 
(d) 2W --7 Hz - 2e-
\end{problembox}

\begin{problembox}
\textbf{Problem 218} \hfill \textbf{[GRB (O.P. Tandon) p.71 Q105]}
\label{q:ext-grb_all_questions-218}

The strong oxidising agent has: 
(a) high value of reduction potential 
(b) high value of oxidation potential 
(c) low value of reduction potential 
to lose electrons
\end{problembox}

\begin{problembox}
\textbf{Problem 219} \hfill \textbf{[GRB (O.P. Tandon) p.71 Q106]}
\label{q:ext-grb_all_questions-219}

The passage of electricity in the Daniell cell when Zn and 
electrodes are connected is from: 
(a) Cu to Zn in the cell 
(b) Cu to Zn outside the cell 
. (cLZn tQ C:t[ol1.1s.i(;/(;!Jhs: ceJI 
(d) Zn to Cu in the cell
\end{problembox}

\begin{problembox}
\textbf{Problem 220} \hfill \textbf{[GRB (O.P. Tandon) p.71 Q107]}
\label{q:ext-grb_all_questions-220}

Which of the following can be used as an electrode? 
(a) A glass rod 
(b) Awooden stick 
(c) A nail 
(d) Asodastraw
\end{problembox}

\begin{problembox}
\textbf{Problem 221} \hfill \textbf{[GRB (O.P. Tandon) p.71 Q108]}
\label{q:ext-grb_all_questions-221}

W ions are reduced at platinum electrode prior to: 
(a) Zn 2+ 
(b) Cu 2+ 
(c) Ag+ 
(d) 12
\end{problembox}

\begin{problembox}
\textbf{Problem 222} \hfill \textbf{[GRB (O.P. Tandon) p.71 Q109]}
\label{q:ext-grb_all_questions-222}

Which of the following statements is wrong? 
(a) F2 is the strongest oxidising agent as its reduction potential 
is high$/cdot$ 
(b) Li. is the weakest reducing agent as its redtlction potential 
is low 
(c) Li is the strongest reducing agent as its oxidation potential 
is high 
(d) F 
ion does not show reducing property
\end{problembox}

\begin{problembox}
\textbf{Problem 223} \hfill \textbf{[GRB (O.P. Tandon) p.71 Q110]}
\label{q:ext-grb_all_questions-223}

Out of Cu, Ag, Fe and Zn, the metal which can displace all 
others from their salt solution is: 
(a) Ag 
(b) Cu 
(c) Fe : 
(d) Zn 
IH. When the electric current is passed through a cell having an 
electrolyte, the positive ions move towards cathode and 
negative ions towards the anode. If the cathode is pulled out of 
the solution: 
(a) the positive and negative ions will move towards anode$/cdot$ 
(b) the positive ions will start moving towards the. anode 
while negative ions will stop moving 
(c) the negative ions will continue to move towards anode 
while positive ions will stop moving 
. (d) the positive and negative ions will start moving randomly'
\end{problembox}

\begin{problembox}
\textbf{Problem 224} \hfill \textbf{[GRB (O.P. Tandon) p.71 Q112]}
\label{q:ext-grb_all_questions-224}

The oxidation potentials ofZn, Cu, Ag, H2 and Ni are 0.76, 
-c 0.34, - 0.80, 0 and 0.25 volt respectively. Which of the 
following reactions will provide maximum voltage') 
(a) Zn + Cu 2+ --7 Cu + Zn 2+ . 
(b) Zn + 2Ag+ --7 2Ag + Zn 2+ 
(c) Hz + Cu 2+ --7 2H+ + Cu 
(d) Hz + Ni2+ 
2H + + Ni 
[PET (KeraJa) 20071 
(e) Zn(s) + 2H' (aq.) /textasciitilde() 
Zn 2+ (aq.) + H2 (g)
\end{problembox}

\begin{problembox}
\textbf{Problem 225} \hfill \textbf{[GRB (O.P. Tandon) p.71 Q113]}
\label{q:ext-grb_all_questions-225}

Which one of the following will increase the voltage of the 
cell? 
Sn + 
---7 Sn 2+ + 2Ag 
(a) Increase in the size of silver rod 
(b) Increase in the concentration of Sn 2+ ions \underline{\hspace{1cm}} 
(c) Increase in the concentration of Ag+ ions 
(d) None of the above
\end{problembox}

\begin{problembox}
\textbf{Problem 226} \hfill \textbf{[GRB (O.P. Tandon) p.71 Q114]}
\label{q:ext-grb_all_questions-226}

A chemist wants to produce CI2 (g ) from molten NaC!. How 
many grams could be produced if he uses a steady current of 2 
ampere for 2.5 minutes? 
(a) 3.55 g 
(b) 1.775 g 
(c) 0.110 g 
(d) 0.1775 g
\end{problembox}

\begin{problembox}
\textbf{Problem 227} \hfill \textbf{[GRB (O.P. Tandon) p.71 Q115]}
\label{q:ext-grb_all_questions-227}

In the electrolysis of CuCl l solution, the mass of cathode 
------lncreaseabyoAg.Wllat occurreaa:f copperl!.llode? 
(a) 0.224 litre ofCl2 was liberated 
(b) 1.12 litre of oxygen was liberated 
(c) 0.05 mole Cu 2+ passed into the solution 
.' (dlQ.)/textasciitilde()01/textasciitilde()<=1l2+p/textasciitilde() se<i.il1\underline{\hspace{1cm}}to\underline{\hspace{1cm}}/textasciitilde()/textasciitilde()/textasciitilde() sollition
\end{problembox}

\begin{problembox}
\textbf{Problem 228} \hfill \textbf{[GRB (O.P. Tandon) p.71 Q116]}
\label{q:ext-grb_all_questions-228}

Consider the reaction, 
C12(g) + 2Br-(aq.) ---7 2Cnaq.) + Br2 
The emf of the cell, when [Cn = [Br2] 
[Br-] = om M and 
CI 2 gas is at I atm pressure, will be: (E 0 for the above reaction 
is 0.29 volt) 
. 
(a) 0.54 volt (b) 0.35 volt (c) 0.24 volt (d) - 0.29 volt 
117 .. If A:, A: and A $/circ$ refer to equivalent conductance of a cation, an 
anion and equivalent conductance of the salt at infinite 
dilution, then according to Kohlrausch's law: 
(a) AO = A: + A: 
(c) AO = A: - A: 
(b) 'A$/circ$=Ao -Ao. 
e 
a 
. (d) A $/circ$ = A: fA:
\end{problembox}

\begin{problembox}
\textbf{Problem 229} \hfill \textbf{[GRB (O.P. Tandon) p.71 Q118]}
\label{q:ext-grb_all_questions-229}

How much silver will be obtained by that quantity of current 
which displaces 5.6 litre of H)? 
[PMT (Pb.) 19931 
(a) 54 g 
(b) l3.5 g 
(c) 20 g 
(d) 108 g
\end{problembox}

\begin{problembox}
\textbf{Problem 230} \hfill \textbf{[GRB (O.P. Tandon) p.71 Q119]}
\label{q:ext-grb_all_questions-230}

The specific conductance of a salt of 0.0 I M concentration is 
l.061 x 10-4 . Molar conductance of the same solution will be: 
(a) 1.061 x 10-4 (b) 1.061 
(c) 10.61 
IPMT (Pb.) 1993) 
(d) 106,1
\end{problembox}

\begin{problembox}
\textbf{Problem 231} \hfill \textbf{[GRB (O.P. Tandon) p.71 Q120]}
\label{q:ext-grb_all_questions-231}

What is the number of coulombs required for the conversion 
of one mole of MnO; to one mole of Mn 2+? 
[MBBS (Orissa) 1993) 
(a) 5 x 96500 (b) 3 x 96500 (c) 96500 
(d) 9650
\end{problembox}

\begin{problembox}
\textbf{Problem 232} \hfill \textbf{[GRB (O.P. Tandon) p.71 Q121]}
\label{q:ext-grb_all_questions-232}

Which of the following solutions ofNaCI /vill have the highest 
specific conductance? 
Ca) 0.001 N 
(b) 0.1 N 
(c) 0.01 N 
(d) l.ON
\end{problembox}

\begin{problembox}
\textbf{Problem 233} \hfill \textbf{[GRB (O.P. Tandon) p.71 Q122]}
\label{q:ext-grb_all_questions-233}

The equivalent conductance of a I N solution of an electrolyte 
is nearly: 
Ca) 103 times its specific conductance 
(b) 10-3 times its specific conductance 
(c) 100 times its specific conductance 
(d) the same as its specific conductance
\end{problembox}

\begin{problembox}
\textbf{Problem 234} \hfill \textbf{[GRB (O.P. Tandon) p.71 Q123]}
\label{q:ext-grb_all_questions-234}

Zinc is coated over iron to prevent rusting of iron because: 
(a) it is cheaper than iron
\end{problembox}

\begin{problembox}
\textbf{Problem 235} \hfill \textbf{[GRB (O.P. Tandon) p.72 Q124]}
\label{q:ext-grb_all_questions-235}

Free energy chang/textasciitilde() (ilG) isrelated to the emf of the cell (E) 
(a)E=-nFilG 
nFE 
(c) ilG 
RT . 
-F logE 
n 
(b) ilG == 
(d) AG 
nF 
-logE 
RT 
125 .. The standard reduction potentials at 25$/circ$C of Li+ I Li, 
Ba 2+1 Ba, Na+/Na and Mg2+ IMg are - 3.05, 
2.73, 
2.71 
and - 2.37 volt respectively. Which one of the following is the 
strongest oxidising agent? 
(a) Na+ 
(c) Ba2+ 
(b) Li+ 
(d) Mg2+
\end{problembox}

\begin{problembox}
\textbf{Problem 236} \hfill \textbf{[GRB (O.P. Tandon) p.72 Q126]}
\label{q:ext-grb_all_questions-236}

Three faradays of electricity was passed through an aqueous 
solution of iron(Il) bromide. The mass of iron metal (At. mass 
56) deposited at the cathode is: 
(EAMCET 1991) 
(a) 56 g 
(b) 84 g 
(c) 112 g 
(d) 168 g 
127 $/bullet$. The standardeIectrod.epottmtialsofZn, Agand.-Cuare-0;76,$/cdot$$/cdot$ 
0.80 and 0.34 volt respectively; then: 
(a) Ag can oxidise Zn and Cu 
(b) Ag can reduce Zn 2+ and Cu 2+ 
(c) Zn can reduceAg+ andCu 2+ 
(d) Cu can oxidise Znand Ag
\end{problembox}

\begin{problembox}
\textbf{Problem 237} \hfill \textbf{[GRB (O.P. Tandon) p.72 Q128]}
\label{q:ext-grb_all_questions-237}

The standard emf forthe cell reaction, 
Zn + Cu 2+ /textasciitilde() 
Zn 2+ + Cu 
is 1.10 volt at 25$/circ$C. The emf for the cell reaCtion when 0.1 M 
Cu 2+ and 0.1 M Zn 2+ solutions are used at 25$/circ$C is: 
(a) 1.10volt 
(c) 
1.10 volt 
(b) 0.110 volt 
(d) 
0.110 volt 
(MLNR 1994)
\end{problembox}

\begin{problembox}
\textbf{Problem 238} \hfill \textbf{[GRB (O.P. Tandon) p.72 Q129]}
\label{q:ext-grb_all_questions-238}

Three mole of electrons are passed through three solutions in 
succession containingAgNO). CuS04 and AuCl3 respectively. 
The ratio ctf amounts of cations reduced at cathode will be: 
(a) 1 .. : 3 
(b) 2: 1 : 3 
(c) 3, '2 : 1 
(d) 6. 3 : 2
\end{problembox}

\begin{problembox}
\textbf{Problem 239} \hfill \textbf{[GRB (O.P. Tandon) p.72 Q130]}
\label{q:ext-grb_all_questions-239}

In the electrolysis of an aqueous solution ofNaOH, 2.8 litre of 
oxygen gas at NTP was liberated at anode. How much of 
hydrogen gas was liberated at cathode? 
(a) 2.8 litre 
(b) 5.6 litre 
(c) l1.1-litre 
(d) 22.4 litre
\end{problembox}

\begin{problembox}
\textbf{Problem 240} \hfill \textbf{[GRB (O.P. Tandon) p.72 Q131]}
\label{q:ext-grb_all_questions-240}

Two half-cells have potentials 
0.44 and 0.799 volt 
respectively. These two are coupled to make a galvanic cell. 
Which of the following will be true? 
(a) Electrode of half-cell potential- 0.44 V will act as anode 
(b) Electrode of half-cell potential - 0.44 V will act as 
cathode 
(c) Electrode of half-cell potential 0.799 V will act as.anode 
(d) Electrode of half-cell potential - 0.44 V will act as a$/cdot$ 
positive terminal
\end{problembox}

\begin{problembox}
\textbf{Problem 241} \hfill \textbf{[GRB (O.P. Tandon) p.72 Q132]}
\label{q:ext-grb_all_questions-241}

When a Itad storage battery is charged: 
(a) Pb02 dISSOl(;eS 
(b) the lead electrode becomes coated with lead sulphate$/cdot$ 
(c) sulphuric acid is regenerated 
(d) the amount of acid decreases
\end{problembox}

\begin{problembox}
\textbf{Problem 242} \hfill \textbf{[GRB (O.P. Tandon) p.72 Q137]}
\label{q:ext-grb_all_questions-242}

An example of a simple fuel cell is: 
(a) lead storage battery 
(b) H2 -
O2 cell 
(c) Daniell cell 
(d) Lechlanche cell 
For the cell reaction, 
Mg(s)+ 2Ag+(aq.)/textasciitilde() Mg 2+ (aq.)+ 2Ag(s). 
E/textasciitilde()ell is +3.17 V at 298 K. The value of E II' ilGo and Q at 
+ 
2+ 
/textasciitilde() 
Ag 
and Mg 
concentrations of 0.001 M and 0.02 M 
respectively are:' 
(ISAT 2010) 
(a) 3.04 V, -605.8 kJ mol-I, 20000 
(b) 3.04 V, 611.8 kJ mol-I, 20000 
(c) 3.13 V, -604 kJ mol-I, 20 
3.04 
-611.8 
20000 
eo=+ 3.17 V, n '= 2 
Q "" [Mi+] =/textasciitilde()=20000 
[Ag+ r [o.ooli 
AGO -nFeo 
= =2 x 96500x3.17 
= -611.8kJ 
0.059 
E = eo---loglo Q 
n 
0.059 
3.17---10g(2000G) 
2 
+ 3.04 V] 
Which of the following statements is correct? 
(a) The temperature coefficient of electrolytic conductance is 
positive 
(b) The temperature coefficient of electrolytic resistance is 
negative 
(c) The resis/textasciitilde()ance of an electrolyte decreases with decreasing 
temperature 
(d) The resistance of electrolytic conductors is independent of 
temperature 
Kohlrausch's law states that at: 
ICB1',E (PMT) 2008] 
(a) infinite dilution, each ion makes definite contribution to 
conductance of an electrolyte, whatever be the nature of 
the other ion of the electrolyte 
(b) infmite dilution, each ion makes definite contribution to 
the equivalent conductance of an electrolyte, whatever be 
the nature of the other ion of the electrolyte 
(c) finite dilution, each ion makes definite contribution to the 
equivalent conductance of an electrolyte, whatever be the 
nature ofthe other ion ofthe electrolyte 
(d) infinite dilution, each ion makes definite contribution to 
equivalent conductance of an electrolyte depending on the 
nature of the other ion of the electrolyte 
Which one of the following conditions will increase the 
voltage of the cell represented by the equation? 
Cu(s) + 
(aq.) --'>. Cu2+(aq.) + 2Ag(s)(KCET 2006) 
(a) Increase in the dimension ofCu electrode 
(b) Increase in the dimension of Ag electrode 
(c) Increase in the concentration of Cll 2+ io/textasciitilde() 
(d) Increase in the.concentration of Ag + ion 
[Hint: 
Cell voitage will increase, eitiler by increasing the 
concentration of Cu2+ ion or by decreasing the concentration of 
Ag+"
\end{problembox}

\begin{problembox}
\textbf{Problem 243} \hfill \textbf{[GRB (O.P. Tandon) p.73 Q138]}
\label{q:ext-grb_all_questions-243}

Which of the following reactions occurs at cathode during 
potential of Z > Y > X, then: 
(lIT 1999) 
charging of storage battery? 
(MGIMS 2007) 
(a) Y will oxidise X and not Z 
(a) Pb
2
+ + 2e-/textasciitilde() Pb 
(b) Y will oxidise Z and not X 
(b) Pb ----7 Pb2+ + 2e-
(c) Y will oxidise both Z and X 
(c) Pb2+ + SO/textasciitilde()-
PbS04 
(d) Y willreouce both X and Z
\end{problembox}

\begin{problembox}
\textbf{Problem 244} \hfill \textbf{[GRB (O.P. Tandon) p.73 Q146]}
\label{q:ext-grb_all_questions-244}

The number of coulombs required for the deposition ofl 07 .87 
(d) PbS04 + 2H20 
Pb02 + 4H + + SO/textasciitilde()- +2e-
g of silver is: 
[PMT (MP)1998]
\end{problembox}

\begin{problembox}
\textbf{Problem 245} \hfill \textbf{[GRB (O.P. Tandon) p.73 Q139]}
\label{q:ext-grb_all_questions-245}

The amount of silver deposited on passing 2F of electricity 
(a) 48250 
(b) 10000 
(c) 96500 
(d) 19300 
\underline{\hspace{1cm}}-':,m,\underline{\hspace{1cm}}.m\underline{\hspace{1cm}}through-<l..q1!-.e.llUSmsnbltion of A/textasciitilde()03/textasciitilde()---'MHT-CET 20filLm .. -14'7 .$/bullet$ -At-25$/circ$C,-thestandar. lemf.of.z.-cell.havingcreactionin.YOhdn/textasciitilde()c\underline{\hspace{1cm}}. 
, 
(a) 54 g 
(b) 108 g 
(c) 216 g 
(d) 324 g 
two electrons change is found to be 0.295 V. The equilibrium
\end{problembox}

\begin{problembox}
\textbf{Problem 246} \hfill \textbf{[GRB (O.P. Tandon) p.73 Q140]}
\label{q:ext-grb_all_questions-246}

Statements:' 
' 
constant ofthe reaction is: 
(lIT 1999) 
(i) Unit of specific conductivity is /Omega-I em-I. 
(a) 29.5 x 10-2 
(b) 10 
(ii) /pecific conductivity of strong electrolytes decreases on 
/cdotdilutIon./cdot/cdot/cdot 
(iii) The amount of an ion discharged during electrolysis does 
not depend upon resistance. 
(iv) The unit of electrochemical equivalence is g/coulomb. 
, (a) All are correct 
(b) All are wrong 
(c) Only (i), (ii) and (iv) are correct 
(d) Only (ii), (iii) and (iv) are correct
\end{problembox}

\begin{problembox}
\textbf{Problem 247} \hfill \textbf{[GRB (O.P. Tandon) p.73 Q141]}
\label{q:ext-grb_all_questions-247}

Which among the following expressions is/are not correct? 
(a) /textasciitilde()= = 'Y + J..";. + 'Y \underline{\hspace{1cm}} A": 
(b)' A= = \underline{\hspace{1cm}}1 J..='+ I 
+ 
+ 
-
n 
n. 
(c) J../textasciitilde()alion 
Il:tion x faraday 
(d) A:;',.ion =Il:tion x faraday
\end{problembox}

\begin{problembox}
\textbf{Problem 248} \hfill \textbf{[GRB (O.P. Tandon) p.73 Q142]}
\label{q:ext-grb_all_questions-248}

For 
the 
/textasciitilde()lectrochemical 
cell, 
M I M'" II X I X, 
EO + 
:= 0.44 V and E.
0 
\underline{\hspace{1cm}}:= 0.33 V. From this data we can 
M 1M 
XIX 
deduce that: 
lilT (Screening) 2000] 
(a) M + X 
M+ + X- is the spontaneous reaction 
(b) M+ + X-
M + X is the spontaneous reaction 
(c) 
O.77V 
(d) EcelJ 
0.77 V
\end{problembox}

\begin{problembox}
\textbf{Problem 249} \hfill \textbf{[GRB (O.P. Tandon) p.73 Q143]}
\label{q:ext-grb_all_questions-249}

For the cell reaction,
\end{problembox}

\begin{problembox}
\textbf{Problem 250} \hfill \textbf{[GRB (O.P. Tandon) p.73 Q144]}
\label{q:ext-grb_all_questions-250}

Cu2+(aq,) + Zn(s) /textasciitilde() 
Zn2+ (aq.) + Cu(s) 
(CI ) 
(C2 ) 
the change in free energy (/textasciitilde()G) at a given temperature is a 
function of: 
(eBSE 1998) 
(a) In C1 
(b) In (/textasciitilde(): J 
(c) In (Cl + C2 ) 
(d) In C2 
The standard reduction potential values of three metallic 
cations X, Y and Z are 0.52, -3.03 and -1.18 V respectively. 
The order of reducing powers of the corresponding metals is: 
(HT 1998) 
(c) 1 X 1010 
(d) 29.5 x 1010 
148:' the'-emi'of the 'ceil In wllIchthe foHoWlng -reaction;" " 
Zn(s) + Ni2+ (0.1 M) /textasciitilde() 
Zn2+ (1.0 M) + Ni(s) 
, 
occurs; is found to 0.5105 V at 298 K The standardemfofthe 
cell is: 
(lIT 1998) 
(a) 0.4810 V 
'(b) 05696V 
(c) - 0.5105 V 
(d) 0.5400 V 
[ 
E 
EO 
0.059 10 J..] 
,Hint: 
cell = 
cell 
2 
g 0.1
\end{problembox}

\begin{problembox}
\textbf{Problem 251} \hfill \textbf{[GRB (O.P. Tandon) p.73 Q149]}
\label{q:ext-grb_all_questions-251}

' The molar conductances of NaCl, HCI and CH3COONa at 
infinite dilution are 126.45, 426.16 and 91 $/Omega$-I cm2 mol-I 
respectively. The molar conductance ofCH3COOH at infinite 
dihition is: . 
" 
leBSE 1997; DeE 2009] 
(a) 201.28 $/Omega$-I cm2 mol-1 (b) 390.71 $/Omega$-I cm2 morl 
(c) 698.28 $/Omega$-l cmz morl (d) 540.48 $/Omega$-l cmz mol-l
\end{problembox}

\begin{problembox}
\textbf{Problem 252} \hfill \textbf{[GRB (O.P. Tandon) p.73 Q150]}
\label{q:ext-grb_all_questions-252}

The specific conductance of a 0.1 N KCI solution at 23$/circ$C is 
0.0112 $/Omega$-I cm-I.The resistance of the ceIl containing the 
solution atthe same temperature was found to be 55 $/Omega$. The 
cell constant will be: 
(CBSE 1999) 
(a) 0'.142 em-I 
(b) 0.918 cm-1 
(c) 1.12 cm-I 
Cd) 0.616 cm-l
\end{problembox}

\begin{problembox}
\textbf{Problem 253} \hfill \textbf{[GRB (O.P. Tandon) p.73 Q151]}
\label{q:ext-grb_all_questions-253}

Which of the following plots r/textasciitilde()presents correctly variation of 
equivalent conductance (A) with dilution for a strong 
electrolyte? 
A 
Dilution 
(a) 
Dilution 
(c) 
A 
Dilution 
(b) 
Dilution 
(d)
\end{problembox}

\begin{problembox}
\textbf{Problem 254} \hfill \textbf{[GRB (O.P. Tandon) p.74 Q152]}
\label{q:ext-grb_all_questions-254}

ELECTROCHEMISTRY 
847 
The standard reduction potential for Cu 2+ I Cu and 
Cu2+ I Cu + are 0.337 V and 0.1.53 V respectively. The 
standard electrode potential of Cu + I Cu half cell is: 
[AMU (PM.T) 2009; CBSE (PMT) 2009] 
(b) 0.827 V
\end{problembox}

\begin{problembox}
\textbf{Problem 255} \hfill \textbf{[GRB (O.P. Tandon) p.74 Q160]}
\label{q:ext-grb_all_questions-255}

What is the amount of chlorine evolved when 2 ampere of 
current is passed for 30 minutes in an aqueous solution of 
NaCl? 
(SHU 1998; AIIMS 1999) 
(a) 66 g 
(b) 1.32 g 
(c) 33 g 
(d) 99 g
\end{problembox}

\begin{problembox}
\textbf{Problem 256} \hfill \textbf{[GRB (O.P. Tandon) p.74 Q161]}
\label{q:ext-grb_all_questions-256}

When 9.65 coulomb of electricity is passed through a solution
\end{problembox}

\begin{problembox}
\textbf{Problem 257} \hfill \textbf{[GRB (O.P. Tandon) p.74 Q153]}
\label{q:ext-grb_all_questions-257}

Solubility of a sparingly soluble salt S ,specific conductance K 
of silver nitrate (Atomic mass of Ag 
108 g mor
l
), the 
and the equivalent conductance Ao are related as: 
amount of silver deposited is: 
(KCET 2000) 
(d) 0.490 V 
(a) 0.184 V 
(c) 0.521 V 
[lAS (Pre.) 1997] 
(a) 16.2 mg 
(b) 21.2 mg 
(a) S 
1000 Ao 
(c) 10.8 mg 
(d) 6.4 mg 
K 
(b) S = KAo
\end{problembox}

\begin{problembox}
\textbf{Problem 258} \hfill \textbf{[GRB (O.P. Tandon) p.74 Q162]}
\label{q:ext-grb_all_questions-258}

The quantity of electricity needed to deposit 127.08 g of 
K 
1000 
copper is: 
ICET (Pb.) 2000; PET (MP) 2004] 
(c) S =: 
(d) S /textasciitilde() 
----- --lOOOAo---$/cdot$$/cdot$$/cdot$$/cdot$----- -Ao ---- - \underline{\hspace{1cm}} .-(a)J faI:ad..a)' \underline{\hspace{1cm}} ---------Cb)A-co.ulo/textasciitilde() \underline{\hspace{1cm}} ..... \underline{\hspace{1cm}} . 
(c) 4 faraday 
(d) 1 ampere
\end{problembox}

\begin{problembox}
\textbf{Problem 259} \hfill \textbf{[GRB (O.P. Tandon) p.74 Q154]}
\label{q:ext-grb_all_questions-259}

2Ag + (aq.) + Cu(s) 
C/textasciitilde() 2+ (aq.) + 2Ag(s)
\end{problembox}

\begin{problembox}
\textbf{Problem 260} \hfill \textbf{[GRB (O.P. Tandon) p.74 Q163]}
\label{q:ext-grb_all_questions-260}

The charge required to deposit 9 g of Al fromAl3+ solution is 
(E/textasciitilde()ell 
+ 0.46 V) 
(At. wt. of /textasciitilde()l 27.0): 
rMEE (Kerala) 2000]
\end{problembox}

\begin{problembox}
\textbf{Problem 261} \hfill \textbf{[GRB (O.P. Tandon) p.74 Q159]}
\label{q:ext-grb_all_questions-261}

Whichchangewillincreasepotentialthel;nost?' 
(a) 3116.3'C (b) 96500C (c) 9650C 
(d) 32163C 
(a) DoubliJlgthe [Ag+]
\end{problembox}

\begin{problembox}
\textbf{Problem 262} \hfill \textbf{[GRB (O.P. Tandon) p.74 Q164]}
\label{q:ext-grb_all_questions-262}

The-emfofthe cellNiINi/textasciitilde()+(-1.0M)1I Au/textasciitilde()(-LOM)IAu.is --
(b) Halving the [Cu 2+] 
for-Ni 2+ INi = - 0.25 V; EO for Au 3+ IAu =: 1.5 V): 
(c) Doubling the size ofCu electrode 
(d) Halving the size of Ag electrode 
As a lead storage battery is charged: IPMT (M.P) 1993, 2000] 
(a) lead dioxide dissolves 
(b) sulphuric acid is regenerated 
(c) lead electrode becomes coated with lead sulphate 
(d) the concentration of sulphuric acid decreases 
In the electrochemical reaction, 
2Fe3+ + Zn /textasciitilde() 
Zn 2+ + 2Fe2+ 
increasing the concentration ofFe2+: 
ICEE (Kenlia) 20001 
(a) increases the cell emf 
(b) increases the current flow 
(c) decreases the cell emf 
(d) alters the pH of the solution 
In the cell rea<;tion, 
Cu(s) + 2Ag+(aq.) /textasciitilde() 
Cu 2+(aq.) + 2Ag(s). 
E/textasciitilde()ell =: 0.46 V. By doubling the concentration ofCu 2+, E/textasciitilde()ell 
is: 
[MEE (Kerala) 20001 
(a) doubled 
(b) halved 
. (c) increased but less than double 
(d) decreased by a small fraction 
The conductivity of 0.01 moIJdm3 aqueous acetic acid at 300 
K 
is 19.5 X 10-5 $/Omega$-l em-I 
and the limiting molar 
conductIvity of acetic acid at the same temperature is 390 
$/Omega$-l cm2 morl. The degree of dissociation of acetic acid is: 
(a) 0.5 
(b) 0.05 
[lAS (Pre.) 1995] 
ec) 5 x 10-3 
(d) 5 x 10-7 
The ionization constant of a weak electrolyte is 25 x 10-6 
while the eq,uivalfnt conducta/textasciitilde()ce of its 0.01 M solution is 
19.6 s cm- eq-. The eqUlvalent conductance of the 
electrolyte at infinite dilution (in s em 2 eq -I) will be: 
(a) 250 
(c) 392 
(b) 196 
(d) 384 
[lAS (Pre.) 1998] 
ICET (Pb.) 2000; PMT (MP) l,,,:'''' 
(a) 
.25 V 
(b) +1.75 V (c) -1.25 V . (d) -1.75 V
\end{problembox}

\begin{problembox}
\textbf{Problem 263} \hfill \textbf{[GRB (O.P. Tandon) p.74 Q165]}
\label{q:ext-grb_all_questions-263}

. F Of the reduction of silver ions with copper metal; the standard 
cell potential was found to be +0.46 V at 25$/circ$C. The value of
\end{problembox}

\begin{problembox}
\textbf{Problem 264} \hfill \textbf{[GRB (O.P. Tandon) p.74 Q166]}
\label{q:ext-grb_all_questions-264}

167. 
standard Gibbs energy, AGO will be : 
(F ,;" 96500C mol-I) 
(a) -89 kJ 
(b)-89 J 
(c)-44.5 kJ 
(d) -98 kJ 
(AIPMT 2010) 
In electrolysis of NaCI when Pt electrode is' taken then H2 is 
liberated at cathode while with Hg cathode it forms sodium 
amalgam: 
. 
[eBSE (PMT) 20021 
(a) Hg is more inert than Pt 
(b) more voltage is required to reduce 
at Hg than at Pt 
(c) Na is dissolved in Hg while it does not dissolve in Pt 
(d) concentration ofF ions is larger when Pt electrode is taken 
Zn gives H2 gas with H2S04 and HCl but not with HN03 
because: 
[CBSE (PMT) 20021 
(a) Zn acts as oxidising agent when reacts with HN03 
(b) HN03 is weaker acid than H2S04 and HCl 
(c) in electrochemical series, Zn is above hydrogen 
;/ 
(d) NO:l is reduced in preference to hydronium ion ,'(:"-
\end{problembox}

\begin{problembox}
\textbf{Problem 265} \hfill \textbf{[GRB (O.P. Tandon) p.74 Q168]}
\label{q:ext-grb_all_questions-265}

A current is passed through two voltameters connected in 
series. The first voltameter contains XS04 (aq.) while the 
second voltameter contains YiS04 (aq.). The relative atomic 
masses of X andY are in the ratio 2: 1. The nitio of the mass of 
X liberated to the mass of Y liberated is: 
(a) I: 1 
(b) 1: 2 
(c) 2: 1 
(d) none of these
\end{problembox}

\begin{problembox}
\textbf{Problem 266} \hfill \textbf{[GRB (O.P. Tandon) p.74 Q169]}
\label{q:ext-grb_all_questions-266}

Which of the following reactions is possible at anode? 
(AIEEE 2002) 
(a) 2Cr3+ + 7H20 /textasciitilde() 
Cr20/textasciitilde()- + 14H+ 
(b) F2 /textasciitilde() 
2F-
(c) .:. 0, + 2H+ /textasciitilde() 
H20 
2 
-
(d) None of the above
\end{problembox}

\begin{problembox}
\textbf{Problem 267} \hfill \textbf{[GRB (O.P. Tandon) p.75 Q171]}
\label{q:ext-grb_all_questions-267}

Conductivity (Unit Siemen's 'S ') is directly proportional to 
area of the vessel and the concentration of the solution in it and 
is inversely proportional to the length of the vessel, then the 
-unitoLconstant/textasciitilde()oLprQPortionalil /textasciitilde()s:/textasciitilde() ........ /textasciitilde()/textasciitilde()\underline{\hspace{1cm}}1AIEEE 
... 20\underline{\hspace{1cm}}02) .. . 
(a) S m mol-I 
(b) S m 2 mor1 
(c) S-2 
mol 
(d)S2 m:2 /textasciitilde()101-2 .
\end{problembox}

\begin{problembox}
\textbf{Problem 268} \hfill \textbf{[GRB (O.P. Tandon) p.75 Q172]}
\label{q:ext-grb_all_questions-268}

Zn(s) + Cu2+(aq.)/textasciitilde() Cu(s) + Znz+(aq.) 
-/textasciitilde()/textasciitilde()'-J(eaaionquofieirt Q;:/textasciitilde()/textasciitilde()/textasciitilde()::: /textasciitilde()-/textasciitilde()vanatiOilE celtwitli:QisOf tlie 
. "C //$/cdot$ithOA = 1.10 volt. Eceu 
1.1591 volt when: 
t/textasciitilde() 
EC/textasciitilde()IIL...I \underline{\hspace{1cm}} /textasciitilde() 
\underline{\hspace{1cm}} \underline{\hspace{1cm}} 
10gO---
(a) [Cu 2+] / [Zn2+] om 
(b) [Zn2+] / [Cu 2+] = 0.01 
[Zn2+] / [Cu 2+] = 0.1 
(d) [Zn2+] / [Cu 2+]
\end{problembox}

\begin{problembox}
\textbf{Problem 269} \hfill \textbf{[GRB (O.P. Tandon) p.75 Q173]}
\label{q:ext-grb_all_questions-269}

In which of the following cells will the emfbe independent of 
the activity of the chloride ions?
\end{problembox}

\begin{problembox}
\textbf{Problem 270} \hfill \textbf{[GRB (O.P. Tandon) p.75 Q174]}
\label{q:ext-grb_all_questions-270}

. (a) Zn I ZnC12(aq.) I Pt(Cl2) 
(b) Zn IZnCI2(aq.) II KCl(aq.) IAgCl(s),Ag(s) 
(c) Ag,AgCI(s) I KCl(aq.) I Pt(CI2) 
(d) Hg, Hg2Cl2(s) I KCl(aq.) IIAgN03(aq.) IAg(s) 
Standard electrode potential data are useful for understanding 
the suitability of an oxidant in a redox titration. Some half-cell 
reactions and their standard potentials are given below: 
MnO:; (aq.) + 8H+(aq.) + 5e- ----? Mn 2+(aq.) + 4H20(l); 
E" 
l.51 volt 
CrzOhaq.) + 14H+(aq.) + 6e- ----? 2Cr3+(aq.) + 7H20(l); 
it 
" 
Cl2(g) + 2e- ----? 2Cnaq.); . 
E" = 1.38 volt 
E" = 0.77 volt 
E" 
lAO volt 
Identify the only incorrect statement regarding the quantitative 
estimation of aqueous Fe(N03h: 
(liT 2002) 
(a) MnO; can be used in aqueous HCI 
(b) Cr20/textasciitilde()- can be used in aqueous HCI 
(c) MnO; can be used in aqueous H2S04 
(d) Cr20/textasciitilde()- can be used in aqueous H2S04
\end{problembox}

\begin{problembox}
\textbf{Problem 271} \hfill \textbf{[GRB (O.P. Tandon) p.75 Q175]}
\label{q:ext-grb_all_questions-271}

The standard reduction potential EO, for the half reaction are: 
Zn /textasciitilde() 
Zn 2+ + 2e-; 
$/cdot$Bo 
0.76V 
Cu /textasciitilde() 
Cu 2+ + 2e-; 
EO = 0.34 V 
The emffor the cell reaction, Zn(s) + Cu 2+ -7 Zn2+ + Cu(s) 
is: 
[Bihar CECE (Pre.) 20041 
(a) 0.42 V 
(b) - 0.42 V 
(c) - 1.1 V 
(d) + Ll V
\end{problembox}

\begin{problembox}
\textbf{Problem 272} \hfill \textbf{[GRB (O.P. Tandon) p.75 Q176]}
\label{q:ext-grb_all_questions-272}

In a galvanic cell, the electrons flow from: 
(a) anode to cathode through the solution 
(b) cathode to anode through the solution 
(KCET 2004) 
(c) . anode to cathode through the external circuit 
..\underline{\hspace{1cm}}--\underline{\hspace{1cm}}. 
(d) cathode to anode through ilieexternal circuit------$/cdot$-$/cdot$$/cdot$
\end{problembox}

\begin{problembox}
\textbf{Problem 273} \hfill \textbf{[GRB (O.P. Tandon) p.75 Q177]}
\label{q:ext-grb_all_questions-273}

The standard emf of a galvanic cell involving cell reaction 
with n = 2 is found to be 0.295 Vat 25$/circ$C. The equilibrium 
constant of the reaction would be: 
Given: F = 96500 C mol-I; R = 8.314JK-1 mor1 
... iCBSE /textasciitilde()(PMT)2004;AIEEE 20()4r=---
(a) 2 x 1011 
(b) 4 X 1012 
(c) 1 X 102 
(d) 1 x 101$/circ$
\end{problembox}

\begin{problembox}
\textbf{Problem 274} \hfill \textbf{[GRB (O.P. Tandon) p.75 Q178]}
\label{q:ext-grb_all_questions-274}

The emf of the cell, 
Zn 1 Zn 2+ (0.01 M) II Fe2+ (0.001 M) IFe 
at 298 K is 0.2905 V then the value of equilibrium constant for 
the cell reaction is: 
11IT (S) 2004] 
(a) eO.3210.0295 
(b) 10$/circ$.3210.0295 
(c) 10$/circ$$/cdot$261 0.0295 
(d) 10$/circ$.3210.0591
\end{problembox}

\begin{problembox}
\textbf{Problem 275} \hfill \textbf{[GRB (O.P. Tandon) p.75 Q179]}
\label{q:ext-grb_all_questions-275}

The standard emf of the following electrodes are;
\end{problembox}

\begin{problembox}
\textbf{Problem 276} \hfill \textbf{[GRB (O.P. Tandon) p.75 Q180]}
\label{q:ext-grb_all_questions-276}

181. 
E;e3+ fFe2+ =+ 0.77 V; E;n 2+ fSn = - 0.14 V 
under standard conditions, the potential for the reaction, 
Sn(s) + 2F/textasciitilde()/textasciitilde()+(aq.) ----? 2Fe2+(aq.) + Sn2+(aq.) is: 
(AIEEE 2004) 
(a) 1.68 V 
(b) lAO V 
(c) 0.91 V 
(d) 0.63 V 
The highest electrical conductivity of the following aqueous 
solutions is of: 
(AlEEE 2()05) 
(a) 0.1 M acetic acid 
(b) 0.1 M chloroacetic apid 
(c) 0.1 M fluoroacetic acid 
(d) 0.1 M difluoroacetic acid 
Aluminium oxide may be electro lysed at 1000$/circ$C to furnish 
aluminium metal (Atomic mass = 27 amu; 1 faraday = 96500 
coulomb). The cathode reaction is: 
A13+ + 3e- /textasciitilde() 
AI 
To prepare 5.12 kg of aluminium metal by this method we 
would require: 
(AIEEE 200S) . 
(a) 5049 x 107 C of electricity (b) 1.83 x 107 C of electricity 
(c) 5A9 x 104 C of electricity (d) 5.49 x 1010 C of electricity
\end{problembox}

\begin{problembox}
\textbf{Problem 277} \hfill \textbf{[GRB (O.P. Tandon) p.75 Q182]}
\label{q:ext-grb_all_questions-277}

During the process of electrolytic refining of copper, some 
metals present as impurity settle as 'anode mud'. These are: 
(AIEEE 2005) 
(a) Sn and Ag 
(b) Pb and Zn 
(c) Ag and Au 
(d)Fe and Ni
\end{problembox}

\begin{problembox}
\textbf{Problem 278} \hfill \textbf{[GRB (O.P. Tandon) p.75 Q183]}
\label{q:ext-grb_all_questions-278}

When an acid cell is charged, then: 
(AFMC 2(05) 
(a) voltage of the cell increases 
(b) electrolyte of the cell dilutes
\end{problembox}

\begin{problembox}
\textbf{Problem 279} \hfill \textbf{[GRB (O.P. Tandon) p.76 Q184]}
\label{q:ext-grb_all_questions-279}

How$/cdot$many coulombs of electricity are required for the reduc-
tion of I mole of Mn04 to Mn2+? 
IPMT (KeraJa) 2005] 
(a) 96500C 
(b) L93 X 105 C 
(c) 4.83 xI05 C 
(d) 9.65 x 106 C 
(e) 5.62 x 105 C, .
\end{problembox}

\begin{problembox}
\textbf{Problem 280} \hfill \textbf{[GRB (O.P. Tandon) p.76 Q185]}
\label{q:ext-grb_all_questions-280}

The standard electrode pot/textasciitilde()ntial of Ag+ lAg is + O.SO$/cdot$i/1:and of 
Cu2+ ICu is + 0.34 V. These electrodes are connected through a 
,.galt bridge and if: 
. (PET (Kerala) 2005] 
(a) copper electrode acts as cathode, then E;ell is + 0.46 volt 
(6)slfverclectrodeaCtsasallode;tnen E:ell/textasciitilde()ls':':/textasciitilde() 0 .34 volt ... 
, '(c) copper electrode acts as anode, then E:ell is + 0.46 volt 
(d) silver electrode acts as cathode, then E/textasciitilde()I is - 0.34 volt 
$/cdot$$/cdot$$/cdot$(e)silverelectrode acts as $/cdot$anode;$/cdot$then E:1f IS + );l4 volt$/cdot$
\end{problembox}

\begin{problembox}
\textbf{Problem 281} \hfill \textbf{[GRB (O.P. Tandon) p.76 Q186]}
\label{q:ext-grb_all_questions-281}

The half-cell reaction for the corrosion, 
2ft +.!. O2 + 2e- /textasciitilde() 
H20; 
EO = 1.23 V 
. 2 . 
+ 2e- /textasciitilde() 
Fe(s); 
-0.44 V 
Find the I1G $/circ$ (in kJ) for the overall reaction:$/cdot$ (lIT (S) 2005) 
(a) 
76 kJ 
(b) 
322 kJ . (c) - 161 kJ 
(d) 
152 kJ 
[Hint: 
Applying 
Fe(s) /textasciitilde() 
Fe2++ 2e-; I1G/textasciitilde() 
2H+ + 2e- + /textasciitilde() 02 /textasciitilde() 
H20(l);I1G/textasciitilde() 
I1G/textasciitilde() = (-2F x 0.44) + (-2F x 1.23) 
= (-2 x 96500 x 0.44) + (-2 x 96500 x 1.23) 
-322310J = -322kJ]
\end{problembox}

\begin{problembox}
\textbf{Problem 282} \hfill \textbf{[GRB (O.P. Tandon) p.76 Q187]}
\label{q:ext-grb_all_questions-282}

What is the cell reaction occurring in Daniell cell (galvanic . 
cell)? 
[CET (J and K) 20061 
(a) Cu(s) + ZnS04(aq.) 
CuS04(aq.) + Zn(s) 
. (b) Zn(s) + CuS04(aq.) /textasciitilde() 
Cu(s) + ZnS04(aq.) 
(c) Ni(s) + ZnS04(aq.) /textasciitilde() 
NiS04(aq.) + Zn(s) 
(d) 2Na(s) + CdS04(aq.) /textasciitilde() 
Na2S04(aq.) + Cd(s) 
188.$/cdot$ What are the units of equivalent conductivity of a solution? 
.
\end{problembox}

\begin{problembox}
\textbf{Problem 283} \hfill \textbf{[GRB (O.P. Tandon) p.76 Q189]}
\label{q:ext-grb_all_questions-283}

[CET (J and K) 2006) 
(a) mhocm/textasciitilde()l 
(b) $/Omega$cm-1 gequiv-1 
(c) Inho cm-2 g equiv-1 
(d) mho cm2 g equiv-1 
The amount of copper deposited by one faraday current will be 
maximum in an acidic solution of one litre of: 
(a) 1M Cu 2C12 
(c) 5M CuS04 
(e) 10M CuPz 
[PMT (KeraJa) 2006) 
(b) 2M Cu(N03h 
(d) 5M CU3(P04h 
[Hint: 
Greater is the equivalent mass, more is the amount 
deposited by 1 F charge.]
\end{problembox}

\begin{problembox}
\textbf{Problem 284} \hfill \textbf{[GRB (O.P. Tandon) p.76 Q190]}
\label{q:ext-grb_all_questions-284}

The reduction potential values of M , Nand 0 are + 2.46 V, 
- 1.13 V, 
3.13 V respectively. Which of the following 
orders $/cdot$is correct regarding their reducing property? 
. 
(JEE (Orissa) 2006) 
(a) 0> N > M'. 
(b) 0> M > N 
(c) M > N > 0 
(d) M > 0 > N
\end{problembox}

\begin{problembox}
\textbf{Problem 285} \hfill \textbf{[GRB (O.P. Tandon) p.76 Q191]}
\label{q:ext-grb_all_questions-285}

The molar conductivities of A/textasciitilde()aoAc and A/textasciitilde()CI at infinite 
dilution in water at 25$/circ$C are 91 and 426.2 S cm2 mol-I 
respectively. To calculate, A /textasciitilde()OAC' the additional value 
required is: 
(AIEEE 2006) 
(a) AOH20 
(b). A/textasciitilde()C1 
(c) A';.aOH 
(d) A/textasciitilde()acl 
[lIillt:/textasciitilde()OAc=fu.aOAc$/pm$/textasciitilde()CL/textasciitilde()A/textasciitilde()actL .
\end{problembox}

\begin{problembox}
\textbf{Problem 286} \hfill \textbf{[GRB (O.P. Tandon) p.76 Q192]}
\label{q:ext-grb_all_questions-286}

The equivalent conductances at infinite dilution of HCI and 
NaCi are 426.15 and 126.15 mho C:r;n2 g eq-I respectively. It 
can be said that the mobility of: 
. (CET (Gujarat) 2006) 
(a) W:ions is much more than that ofCl- ions 
. 
(b) CFions is much more than thatofW /textasciitilde()ions$/cdot$ 
(c) IV ions is much more than that ofNa+ ions 
(d) Na+ ions is much more than that ofW ions
\end{problembox}

\begin{problembox}
\textbf{Problem 287} \hfill \textbf{[GRB (O.P. Tandon) p.76 Q193]}
\label{q:ext-grb_all_questions-287}

The tendencies of the electrodes made up ofCu, Zn and Ag to 
release electrons when dipped in their respective salt solutions 
decrease in the order: 
(VITEEE 2006) 
(a) Zn>Ag>Cu 
(b) Cu>Zn.>Ag 
(c) Zti > Cu > Ag 
(d)' Ag > eu > Zn
\end{problembox}

\begin{problembox}
\textbf{Problem 288} \hfill \textbf{[GRB (O.P. Tandon) p.76 Q194]}
\label{q:ext-grb_all_questions-288}

The electrode reaction that takes place at the anode of 
CH4 -
02 $/cdot$fuel cell is: 
(VITEEE 2006) 
. (a) 202 + 8H+ + 8e- /textasciitilde() 
4HzO 
(b) CH4 + 2HzO 
CO2 + 8H+ + 8e-
(c) CH4 + 202 /textasciitilde() 
CO2 + 2H20 
(d) 2H+ + 
/textasciitilde()H2
\end{problembox}

\begin{problembox}
\textbf{Problem 289} \hfill \textbf{[GRB (O.P. Tandon) p.76 Q196]}
\label{q:ext-grb_all_questions-289}

Zh(s)IZn 2+(IMIICu 2+(lM)ICu(s) (E/textasciitilde()ell 
+1.10 V) 
was allowed to be completely discharged at 298K. The relative 
concentration of Zn 2+ to Cu 2+ [[zn:+] J is: 
[Cu +] 
(AIEEE 2007) 
(a) 9.65 x 104 
(b) antilog (24.08) 
(e) 37.3 
(d) l(f73 
. [Hint: 
E =  0 \underline{\hspace{1cm}} 0.059 10 [Zn2+] 
2 
g [Cu 2+] 
0= 1.10 - 0.059 log 
. 
'.2 
[Zn 2+ ] 
log --2 
-
37.3 
[eu + 1 
[Zn 2+] = 1037.3] 
[Cu2+] 
The efficiency ofa cell is given by: 
(CaSE (Med.) 2007) 
(a) I1G 
(b) AG 
/ 
I1S 
!,M! 
f,. 
(c) I1G 
(d)'/textasciitilde() /textasciitilde():-
I
\end{problembox}

\begin{problembox}
\textbf{Problem 290} \hfill \textbf{[GRB (O.P. Tandon) p.77 Q197]}
\label{q:ext-grb_all_questions-290}

An electric current of I amp is passed through acidulated 
water for 160 minutes and 50 seconds. What is the volume'of 
the hydrogen liberated at the anode (as reduced to NTP)? 
(SCRA2007) 
(a) 1.12 litre 
(b) 2.24 litre (c) 11.2 litre 
(d) 22.4 litre 
[Hint: 
V = I + Ve 
96500 
1x9650x11.2 
1121' 
] 
= 
=. 
Itre 
96500
\end{problembox}

\begin{problembox}
\textbf{Problem 291} \hfill \textbf{[GRB (O.P. Tandon) p.77 Q198]}
\label{q:ext-grb_all_questions-291}

'The resistance of Nil 0 solution is found to be /textasciitilde()2.5 x 103 $/Omega$. 
The equivalent conductance of the solution is (cell constant 
= 1.25 cm-I ): 
[PMT (Kerala) 20071 
(a) 2.5 $/Omega$-I cm2 equiv-1 
(b) 5 $/Omega$-I cm2 equiv-1
\end{problembox}

\begin{problembox}
\textbf{Problem 292} \hfill \textbf{[GRB (O.P. Tandon) p.77 Q205]}
\label{q:ext-grb_all_questions-292}

Electrolysis of dilute aqueous NaCI solution was carried out 
by passing 10 milli ampere current. The time required to 
liberate 0.01 mole of H2 gas at the cathode is: (I faraday = 
96500 C mol- 1) 
(lIT 2008) 
(a) 9.65 x 104 sec 
(b) 19.3 x 104 sec 
. 
(c) 28.95 X 104 sec 
(d) 38.6 X 104 sec 
[Hint: Mass of 0.01 mol H2 = 0.02 g 
. 
ItE 
W=--
96500 
0.02=10XlO-
3 xtX1 
96500 
t = 19.3 X 104 sec] 
(c) 2.5 $/Omega$-I cm-2 equiv-I 
(e) 1.25 $/Omega(-1)$ cm2 equiv-1 
(d) 5 $/Omega(-1)$ cm-2 equiv-1 - -- --- -TO()/textasciitilde()lhe emf of a cell containing sodium/copper electrodes is 
3.05 V, if the electrode potential of copper electrode is
\end{problembox}

\begin{problembox}
\textbf{Problem 293} \hfill \textbf{[GRB (O.P. Tandon) p.77 Q199]}
\label{q:ext-grb_all_questions-293}

For strong electrolytes, the plot of molar conductance versus 
.JC is: 
(VITEEE 2007) 
- - - "( iiHiaraooliC--- --
----- - ---(0) Iineai-------
(c) sinusoidal 
(d) circular
\end{problembox}

\begin{problembox}
\textbf{Problem 294} \hfill \textbf{[GRB (O.P. Tandon) p.77 Q200]}
\label{q:ext-grb_all_questions-294}

How long (in hours) must a current of 5 ampere be maintained 
to electroplate 60 g of calcium from molten CaCI2 ? 
(VITEEE 2007) 
(a) 27 hours 
(b) 8.3 hours (c) II hours 
(d) 16 hours
\end{problembox}

\begin{problembox}
\textbf{Problem 295} \hfill \textbf{[GRB (O.P. Tandon) p.77 Q201]}
\label{q:ext-grb_all_questions-295}

Emfofhydrogen electrode in term of pH is (at I atm pressure):
\end{problembox}

\begin{problembox}
\textbf{Problem 296} \hfill \textbf{[GRB (O.P. Tandon) p.77 Q204]}
\label{q:ext-grb_all_questions-296}

(MHT/textasciitilde()CET 2007) 
RT 
(a)EH = /textasciitilde() 
x pH 
2 
F 
(c) EH = 2.303RT pH 
2 
F 
(b)E 
=RT\underline{\hspace{1cm}}I 
H2 . 
F pH 
(d) EH2 = - 0.0591 pH 
The rusting of iron is catalysed by which of the following? 
(a) Fe 
(c) Zn 
(MGIMS 2007) 
(d) H+ 
On the basis of EO values, the strongest oxidising agent is: 
[Fe(CN)6t- /textasciitilde() 
[Fe(CN)6]3- + e-
EO = - 0.35 V 
Fe2+ /textasciitilde() 
Fe3+ + e-
EO = - 0.77 V 
(a) Fe3+ 
[eBSE (PMT) 2008) 
(b) [Fe(CN)6]3-
(c) [Fe(CN)6t-
(d) Fe2+ 
Given EO 1/textasciitilde() 
= - 0.72 V,EO 2+ 
= - 0.42 V 
Q-/textasciitilde() 
h 
Ih 
. 
The potential for the cell, Cr1Cr3+ (O.IM) I! Fe2+ (O.OIM)I Fe 
is: 
(AIEEE 2008) 
(a) - 0.26 V 
(b) 0.26 V 
(c) 0.339 V 
(d) - 0.339 V 
[Hint: E/textasciitilde()ell = - 0.42 - (- 0.72)=+ 0.30 V 
2Cr(s) + 3Fe2+ (0.0 1M) /textasciitilde() 
2Cr3+ (0. 1M) -t /textasciitilde()Fe(s) 
Q = [Cr
3+ ]2 = [0.1]2 = 104 
[Fe2+]3 
[0.01]2 
According to Nemst equation, 
0.059 
E =  0 - ----j 10giO Q 
n . 
= 0.30 \underline{\hspace{1cm}} 0.059 10g104 
6 
= 0.261 V] 
(": n = 6) 
+ 0.34 V, the electrode potential of sodium is: 
reomed (Karnataka) 20081 
(a)-2.7IV 
(b)+2.7IV (c)-3.7IV 
(d)+3.7IV
\end{problembox}

\begin{problembox}
\textbf{Problem 297} \hfill \textbf{[GRB (O.P. Tandon) p.77 Q207]}
\label{q:ext-grb_all_questions-297}

-Whlii-lsthe number -ormoles-Qt'oxygen 'gas-ev01ve'/textasciitilde()d=D=y/textasciitilde()
\end{problembox}

\begin{problembox}
\textbf{Problem 298} \hfill \textbf{[GRB (O.P. Tandon) p.77 Q208]}
\label{q:ext-grb_all_questions-298}

electrolysis of 180 g of water? 
(SeRA 2009) 
(a) 2.5 
(b) 5.0 
(c) 7.5 
(d) 10.0 
The Gibbs energy for the decomposition of Ai20 3 at 500$/circ$C is 
as follows: 
. 
2 
4 
/textasciitilde() 
-Ai20 3 --7 -AI + 02; !::J.rG = + 966 kJ mol 
3 
3 
The potential difference needed for the electrolytic reduction 
of Al 20 3 at 500$/circ$C is at least: 
(AIEEE 2010) 
(a) 2.5 V 
(b) 5.0 V 
(c) 4.5 V 
(d) 3.0.V 
[Hint: In the given reaction, 
3.[2 AI 3+ ]+ 4e ---7.$/pm$ AI 
3 
3 
(n = 4) 
!::J.G =-nFE 
E = \underline{\hspace{1cm}} !::J.G = 966 x 1000 
nF 
4x96500 
=-2.5 V] 
Set-2: The questions given below may have more 
than one correct answers
\end{problembox}

\begin{problembox}
\textbf{Problem 299} \hfill \textbf{[GRB (O.P. Tandon) p.77 Q1]}
\label{q:ext-grb_all_questions-299}

What is the difference between galvanic cell and electrolytic 
cell? 
(a) In galvanic cell, electrical energy is produced while in 
electrolytic cell electrical energy is consumed 
(b) In galvanic cell, anode is (-)ve while in electrolytic cell 
anode is (+ )ve 
(c) In galvanic cell, cathode is( + )ve while in electrolytic cell 
anode is (-)ve 
(d) All are correct
\end{problembox}

\begin{problembox}
\textbf{Problem 300} \hfill \textbf{[GRB (O.P. Tandon) p.77 Q2]}
\label{q:ext-grb_all_questions-300}

Ag I Ag +, KI II AgI I Ag emf is E, then K sp of AgI is given as: 
nF 
[ 8E
O 
J 
(a) K sp = 2.303RT log EO 
(b) In K sp = nF 
8T - EO .. 
nF 
(c) In Ksp =-
EO 
(d) 10 K 
= 
nFEo . 
g 
sp 
2.303RT
\end{problembox}

\begin{problembox}
\textbf{Problem 301} \hfill \textbf{[GRB (O.P. Tandon) p.77 Q3]}
\label{q:ext-grb_all_questions-301}

A hydrogen electrode is placed in a buffer solution of acetic 
acid and sodium acetate in the ratio y: x and x: y has
\end{problembox}

\begin{problembox}
\textbf{Problem 302} \hfill \textbf{[GRB (O.P. Tandon) p.78 Q5]}
\label{q:ext-grb_all_questions-302}

. The main factors whlch affect corrosion are: 
(a) position of rile tal in electrochemical series 
(b) presence of CO2 in water 
( c) presence of impurities in metal 
. (d) presence of protective coating
\end{problembox}

\begin{problembox}
\textbf{Problem 303} \hfill \textbf{[GRB (O.P. Tandon) p.78 Q6]}
\label{q:ext-grb_all_questions-303}

Which is correct about silver platin/textasciitilde()? 
(a) Anode-pure Ag 
(b) Cathode---ubject to be electroplated 
(c) Electrolyte-Na[Ag(CN)z] 
\underline{\hspace{1cm}} . \underline{\hspace{1cm}}ldJJ=:lec1:!s>lyte.:-::-.Agl'!"Q3.\underline{\hspace{1cm}}. \underline{\hspace{1cm}} /textasciitilde() \underline{\hspace{1cm}} . \underline{\hspace{1cm}} . \underline{\hspace{1cm}}. \underline{\hspace{1cm}} \underline{\hspace{1cm}}
\end{problembox}

\begin{problembox}
\textbf{Problem 304} \hfill \textbf{[GRB (O.P. Tandon) p.78 Q7]}
\label{q:ext-grb_all_questions-304}

Lead storage battery contains: 
(a) Pb rod as anode 
(b) Pb rod as cathode 
(c) Pb plates coated with Pb02 act as cathode 
(d) electrolyte is, H2S04 
.
\end{problembox}

\begin{problembox}
\textbf{Problem 305} \hfill \textbf{[GRB (O.P. Tandon) p.78 Q8]}
\label{q:ext-grb_all_questions-305}

During the electrolysis of AgN03 (/lsing Pt electrodes) 
concentration around cathode as well as anode falls from 4M 
to 3M. What will happen if this happened with Ag electrodes? 
. (a) Result /yill remain same 
(b) Concentration around cathode will, fall from 4M to 3M 
but around anode will increase from 4M to 5M 
( /textasciitilde()) Reverse of statement (b) 
(d) Concentration increases from 4M to 5M on both the 
electrodes
\end{problembox}

\begin{problembox}
\textbf{Problem 306} \hfill \textbf{[GRB (O.P. Tandon) p.78 Q9]}
\label{q:ext-grb_all_questions-306}

Emf of the cell Pt: H2 (l atm) I H+(aq.) II AgCIIAg is 0.27 V 
and 0.26 V at 25$/circ$C and 35$/circ$C .. Heat of reaction occurring' 
inside the cell at 25$/circ$C is: 
(a) - 54.S kJ 
(c) -$/cdot$26.05 kJ
\end{problembox}

\begin{problembox}
\textbf{Problem 307} \hfill \textbf{[GRB (O.P. Tandon) p.78 Q10]}
\label{q:ext-grb_all_questions-307}

Given that, 
(b) 26.05 kJ 
(d) + 54.S kJ 
Ni2+ /Ni = 0.25 V, Cu2+ /Cu = 0.34 V, 
Ag+/Ag = O.SO V and Zn2+ /Zn = - 0.76 V 
which of the following reactions under standard condition will 
not take place in the specified direction? 
(a) Ni2+(aq.) + Cu(s) -----7 Ni(s) + Cu 2+(aq.) 
(b) Cu(s) + 2Ag+(aq.) -----7 Cu 2+(aq.) + 2Ag(s) 
(c) Cu(s) + 2W(aq.) -----7 Cu 2+(aq.) + H2(g) . 
(d) Zn(s) + 2W(aq.) -----7 Zn 2+(aq.) + 3H2(g) 
11.' Which of the following statements is/are correct? 
(a) One faraday is the charge carried by one mole of electrons 
(b) If same quantity of electricity flows through the sOlutions 
of 0.1 M AgN03 and 0.1 MCuS04 solutions, same weight 
of silver and copper will be deposited 
(c) Electrochemical equivalent has the units of grams per 
coulomb 
(d) Passage of one faraday of electricity produces one gn$/cdot$.l 
equivalent of the substance at the electrode
\end{problembox}

\begin{problembox}
\textbf{Problem 308} \hfill \textbf{[GRB (O.P. Tandon) p.78 Q12]}
\label{q:ext-grb_all_questions-308}

Which of the following statements is/are not con'ect? 
(a)Zn-Cu cell is called Daniell cell 
(b) Rust is Fe20 3 . 
(c) Saline water slows down rusting 
(d) Pure metals undergo corrosion faster than impure metals
\end{problembox}

\begin{problembox}
\textbf{Problem 309} \hfill \textbf{[GRB (O.P. Tandon) p.78 Q13]}
\label{q:ext-grb_all_questions-309}

In electrolysi/textasciitilde() of very dilute NaOH solution using platinum 
electrodes: 
(a) H2 is evolved at cathode 
(b) H2 is produced at anode 
(c) Na is obtained at cathode 
(d) O2 is produced at anode 
14.' We observe blue colour if: 
(a) Cuelectrode is placed in the AgN03 solution 
(b) Cu electrode is placed in the ZnS04 solution 
J/textasciitilde()LC::/textasciitilde()\underline{\hspace{1cm}}/textasciitilde()/textasciitilde().c.!T().o/textasciitilde()\underline{\hspace{1cm}}isplace.dJn)/textasciitilde()/textasciitilde()/textasciitilde()il. Jfr.!.QL \underline{\hspace{1cm}} \underline{\hspace{1cm}} 
(d) Cu electrode is placed in dil. H2S04
\end{problembox}

\begin{problembox}
\textbf{Problem 310} \hfill \textbf{[GRB (O.P. Tandon) p.78 Q15]}
\label{q:ext-grb_all_questions-310}

In which of the following cell (s); Ecea == E;ell? 
(a) Cu(s) I Cu 2+(0.01 M) II Ag+(O.1 M) I Ag(s) 
(b) Pt(H2) I pH = III Zn 2+(0.01 M) I Zn(s) 
(c) Pt(H2) I pH = 111 Zn 2+(1 M) I Zn(s) 
(d) Pt(H2) I W = 0.01 M II Zn2+(0.01 M) I Zn(s)
\end{problembox}

\begin{problembox}
\textbf{Problem 311} \hfill \textbf{[GRB (O.P. Tandon) p.78 Q16]}
\label{q:ext-grb_all_questions-311}

Rusting on$/cdot$ the surface of iron involves: 
(a) Fe(s) -----7 Fe2+(aq.) + 2e- (at anodic site) 
(b) 02(g) + 4W(aq:) + 4e- -----7 2H20(l) (at cathodic 
(c) 4Fe2+(aq.) + 02(g) + 4B20(l) -----7 2F/textasciitilde()203(S) + sW' 
(d) Fe20 3(s) + xH20(l) -----7 Fe203 $/cdot$xH20
\end{problembox}

\begin{problembox}
\textbf{Problem 312} \hfill \textbf{[GRB (O.P. Tandon) p.78 Q17]}
\label{q:ext-grb_all_questions-312}

Fuel cell involves following reaction(s): 
(a) 02(g) + 2H2.O(l) + 4e- -----7 4mr(aq.) (at cathode) 
(b) 02(g) + 2H20(l) + 4e- -----7 40W(aq.) (at a'lode) 
(c) 2H2(g) + 40W(aq.) -----7 4H20(l) +4e- (at anode) 
(d) 2H2(g) + 4011(aq.) -----7 4H20(l) + 4e- (at cathode) 
.
\end{problembox}

\begin{problembox}
\textbf{Problem 313} \hfill \textbf{[GRB (O.P. Tandon) p.78 Q18]}
\label{q:ext-grb_all_questions-313}

In the following question, more than one of the answers given 
may be correct. Select the correct answers and mark it 
according to the code: 
Codes'; 
[BHU (Mains) 2008] 
(a) 1,2 ancl'3are correct 
(b) 1 and 2 are correct 
(c) 2 and 4 are correct 
(d) I and 3. are correct 
lHint: In a cell Zn(s) I Zn2+ II H+ IH2(Pt); the addition ofH2SO, 
to the cathode compartment, will:
\end{problembox}

\begin{problembox}
\textbf{Problem 314} \hfill \textbf{[GRB (O.P. Tandon) p.78 Q3]}
\label{q:ext-grb_all_questions-314}

shift equilibrium to left
\end{problembox}

\begin{problembox}
\textbf{Problem 315} \hfill \textbf{[GRB (O.P. Tandon) p.78 Q4]}
\label{q:ext-grb_all_questions-315}

shift equilibrium to right]
\end{problembox}

\begin{problembox}
\textbf{Problem 316} \hfill \textbf{[GRB (O.P. Tandon) p.79 Q1]}
\label{q:ext-grb_all_questions-316}

(A) When acidified zinc sulphate solution is electrolysed 
between zinc electrodes, it is zinc that is deposited at the 
cathodeilnd-hydrogen evolution doesnoHakeplace; . 
(R) The electrode potential of zinc is more negative than 
hydrogen as the Dver voltage for the hydrogen evolution 
on zinc is quite large.
\end{problembox}

\begin{problembox}
\textbf{Problem 317} \hfill \textbf{[GRB (O.P. Tandon) p.79 Q2]}
\label{q:ext-grb_all_questions-317}

(A) In electrolysis, the quantity of electricity' needed for 
depositing 1 mole of silver is different from that required 
for 1 mole of copper. 
(R) The molecular weights of silver and copper are different. 
(AlIMS 1996)
\end{problembox}

\begin{problembox}
\textbf{Problem 318} \hfill \textbf{[GRB (O.P. Tandon) p.79 Q3]}
\label{q:ext-grb_all_questions-318}

(A) Equivalent conductance of all electrolytes decreases with 
increasing concentration. 
(R) Lesser munber of ions are aVai,l.able per gram equivalent at 
higher concentration. 
" . 
(AIIMS .1999)
\end{problembox}

\begin{problembox}
\textbf{Problem 319} \hfill \textbf{[GRB (O.P. Tandon) p.79 Q4]}
\label{q:ext-grb_all_questions-319}

(A) Zinc displaces copper from copper sulphate solution. 
(R) The EO is Zn of -0.76 volt and that of copper is +0.34 
volt. 
(AlIMS 1999)
\end{problembox}

\begin{problembox}
\textbf{Problem 320} \hfill \textbf{[GRB (O.P. Tandon) p.79 Q5]}
\label{q:ext-grb_all_questions-320}

(A) An electrochemical cell can be set-up only if the redox 
reaction is ,spontaneous. 
(R) A reaction is spontaneous iffree energy change is negative.
\end{problembox}

\begin{problembox}
\textbf{Problem 321} \hfill \textbf{[GRB (O.P. Tandon) p.79 Q6]}
\label{q:ext-grb_all_questions-321}

(A) If an aqueous solution ofNaCI is electrolysed, the product 
ob/textasciitilde()ained at the cathode is H2 gas and not Na. 
(R) Gases are liberated faster than the metals.
\end{problembox}

\begin{problembox}
\textbf{Problem 322} \hfill \textbf{[GRB (O.P. Tandon) p.79 Q7]}
\label{q:ext-grb_all_questions-322}

(A) Specific conductance decreases with dilution whereas 
equivalent conductance increases:' 
(R) On dilution, number of ions per cc decreases but total 
number of ions increases considerablr 
8 .. (A) The cell constant of a cell depends upon the nature of the 
material of the electrodes. 
' 
(R) The observed conductance of a solution depends upon the 
nature of the material of the electrodes.
\end{problembox}

\begin{problembox}
\textbf{Problem 323} \hfill \textbf{[GRB (O.P. Tandon) p.79 Q9]}
\label{q:ext-grb_all_questions-323}

(A) The ratio of specific conductivity to the observed 
conductance does not depend upon the concentration of 
the solution taken in the conductivity cell. 
(R) Specific conductivity decreases with dilution whereas 
observed conductance increases with dilutie!1.
\end{problembox}

\begin{problembox}
\textbf{Problem 324} \hfill \textbf{[GRB (O.P. Tandon) p.79 Q10]}
\label{q:ext-grb_all_questions-324}

(A) Molar conductivity of a weak electrolyte at infinite 
dilution cannot be determined experimentally. 
(R) Kohlrausch's law helps to find the molar conductivity of a 
weak electrolyte at infinite dilution.
\end{problembox}

\begin{problembox}
\textbf{Problem 325} \hfill \textbf{[GRB (O.P. Tandon) p.79 Q11]}
\label{q:ext-grb_all_questions-325}

(A) One coulomb of electric charge deposits weight equal to 
the electrochemical equivalent of the substance. 
(RTOne faraday depOSIts one mole of the substance. 
rZ:-(A) If 
standard reductlOn potentIal 
for 
the reactlOn, 
Ag+ + e- /textasciitilde() 
Ag is 0.80 volt, then for the reaction, 
2Ag + + 2e - /textasciitilde() 
2Ag, it will be 1.60 volt. 
(R) If concentration of Ag + ions is doubled, the electrode 
potential is also doubled.
\end{problembox}

\begin{problembox}
\textbf{Problem 326} \hfill \textbf{[GRB (O.P. Tandon) p.79 Q13]}
\label{q:ext-grb_all_questions-326}

(A) Gold \underline{\hspace{1cm}}cJ,lJQri<k (AyQ;JJsQi1!tion \underline{\hspace{1cm}}.QfUlIlQLbtu,;jQIed in a 
vessel made of copper, iron, nickel, chromium, zinc or tin. 
(R) Gold is very precious metal.
\end{problembox}

\begin{problembox}
\textbf{Problem 327} \hfill \textbf{[GRB (O.P. Tandon) p.79 Q14]}
\label{q:ext-grb_all_questions-327}

(A)Jn the Daniell cell, if concentrations of Cu 2+ and Zn 2+ 
. ions are doubled, the emf of the cell will be doubled. 
(R) If the concentration of ions in contact with the metals is 
doubled, the electrode potential is doubled.
\end{problembox}

\begin{problembox}
\textbf{Problem 328} \hfill \textbf{[GRB (O.P. Tandon) p.79 Q15]}
\label{q:ext-grb_all_questions-328}

(A) H2 + 02 fuel. cell gives a constant voltage throughout its 
life. 
(R) In this fuel cell, H2 reacts with OH- ions, yet the overall 
. concentration ofOH- ions does not change.
\end{problembox}

\begin{problembox}
\textbf{Problem 329} \hfill \textbf{[GRB (O.P. Tandon) p.79 Q16]}
\label{q:ext-grb_all_questions-329}

(A) Presence of CO2 in the air accelerates cormsion. 
(R) CO2 is a poisonous gas.
\end{problembox}

\begin{problembox}
\textbf{Problem 330} \hfill \textbf{[GRB (O.P. Tandon) p.79 Q17]}
\label{q:ext-grb_all_questions-330}

(A) For the Daniell cell, Zn I Zn2+ II Cu 2+ I Cu with Ecell = 1.1 
volt, the application of opposite potential greater than 1.1 V 
results into flow of electrons from cathode to anode. 
(R) Zn is deposited at anode and Cu is dissolved at cathode. 
(AlIMS 2006)
\end{problembox}

\begin{problembox}
\textbf{Problem 331} \hfill \textbf{[GRB (O.P. Tandon) p.79 Q18]}
\label{q:ext-grb_all_questions-331}

(A) A current of 96.5 ampere is passed into aqueous AgN03 
solution for 100 second. The weight of silver deposited is 
10.8 g. (Atomic weight of Ag = 108) 
(R) The mass of a substance deposited during the electrolysis 
of an electrolyte is inversely proportional to the quantity 
o!:i,electricity passing through the electrolyte . 
. ',". 
[EAMCET (Engg.) 20061
\end{problembox}

\begin{problembox}
\textbf{Problem 332} \hfill \textbf{[GRB (O.P. Tandon) p.79 Q19]}
\label{q:ext-grb_all_questions-332}

(A) According to Kohlrausch'sJaw, the molar conductance of 
a strong electrolyte at infmite dilution is sum of molar 
conductivities of its ions. 
(R) The current carried by cation and anion is always equal. 
(AIIMS 2007)
\end{problembox}

\begin{problembox}
\textbf{Problem 333} \hfill \textbf{[GRB (O.P. Tandon) p.79 Q20]}
\label{q:ext-grb_all_questions-333}

(A) The cell potential of mercury cell is 1.35 V, which remains 
constant. 
(R) In mercury cell, the electrolyte is a paste of KOH and 
ZnO. 
(AIIMS 2008) .
\end{problembox}

\begin{problembox}
\textbf{Problem 334} \hfill \textbf{[GRB (O.P. Tandon) p.80 Q71]}
\label{q:ext-grb_all_questions-334}

(e) 
. 72. (d) 
. 73. (a)
\end{problembox}

\begin{problembox}
\textbf{Problem 335} \hfill \textbf{[GRB (O.P. Tandon) p.80 Q112]}
\label{q:ext-grb_all_questions-335}

(b) 
\underline{\hspace{1cm}} JJJ.\underline{\hspace{1cm}} J<:L \underline{\hspace{1cm}} 1 14 \underline{\hspace{1cm}} (<:L \underline{\hspace{1cm}} Jl/textasciitilde().$/cdot$\underline{\hspace{1cm}}7J9) \underline{\hspace{1cm}} 11/textasciitilde() \underline{\hspace{1cm}} (Q)\underline{\hspace{1cm}} \underline{\hspace{1cm}} 117 $/bullet$. \underline{\hspace{1cm}}(liL \underline{\hspace{1cm}} 118.\underline{\hspace{1cm}}/textasciitilde()/textasciitilde()) \underline{\hspace{1cm}} \underline{\hspace{1cm}},\underline{\hspace{1cm}}112!..\underline{\hspace{1cm}}Jf) \underline{\hspace{1cm}} 
12fi.\underline{\hspace{1cm}}jliL \underline{\hspace{1cm}} . 
$/bullet$
\end{problembox}

\begin{problembox}
\textbf{Problem 336} \hfill \textbf{[GRB (O.P. Tandon) p.80 Q14]}
\label{q:ext-grb_all_questions-336}

(a, e) 
.', 
6 $/bullet$. (c).
\end{problembox}

\begin{problembox}
\textbf{Problem 337} \hfill \textbf{[GRB (O.P. Tandon) p.81 Q1]}
\label{q:ext-grb_all_questions-337}

When aluminium oxide (AI20 3) is electrolysed for the 
production of aluminium metaL. For a given quantity of 
electricity, the number of moles of aluminium obtained if the 
volume of Oz gas obtained is 201.6 litre measured at NTP, is: 
(a:) 9 
(b) 6 
(c) 12 
(d) 4.5 
201.6 
[Hint: 
Number of equivalents of oxygen 
36 
5.6 
(Equivalent volume of oxygen == 5.6 litre at NTP) 
:. Number of equivalents of Al = 36
\end{problembox}

\begin{problembox}
\textbf{Problem 338} \hfill \textbf{[GRB (O.P. Tandon) p.81 Q2]}
\label{q:ext-grb_all_questions-338}

A smuggler could not carry gold by depositing iron on the gold 
surface since: 
(a) gold is denser 
(b) iron rusts 
(ctgoldha:s-higherreducti<mpotentialthaniron . 
. d) gold has lower reduction potential than iron
\end{problembox}

\begin{problembox}
\textbf{Problem 339} \hfill \textbf{[GRB (O.P. Tandon) p.81 Q3]}
\label{q:ext-grb_all_questions-339}

Oil electrolysis, which of the following does not give out 
oxygen? 
(a) Acidic water using Pt electrode 
(b) Fused NaOH using Pt electrode 
(c) Dilute H2S04 using Pt electrode 
(d) Diiute H2S04 using Cu electrode
\end{problembox}

\begin{problembox}
\textbf{Problem 340} \hfill \textbf{[GRB (O.P. Tandon) p.81 Q4]}
\label{q:ext-grb_all_questions-340}

During electrolysis of a solution of AgN03, 9650 coulomb of 
charge pass through the . electro lytic cell; the mass of silver 
deposited on the cathode will be: 
(a) 21.6 g 
(b) 108 g 
(c) 10.8 g 
(d) 1.08 g
\end{problembox}

\begin{problembox}
\textbf{Problem 341} \hfill \textbf{[GRB (O.P. Tandon) p.81 Q5]}
\label{q:ext-grb_all_questions-341}

An electrolytic ceIL contains a solution of AgzS04 and has 
platinum electrodes. A current is passed until 1.6 g of O2 has 
been liberated at anode. The amount of silver deposited at 
cathode will be: 
(a) 108 g 
(b) 1.6 g 
(c)0.8g 
[Hint: 
NUliJber of equivalents of oxygen 
.. 
(d) 21.60 g 
1.6 
0.2 
8 
:. Number of equivalents of Ag deposited 
0.2 
Mass of Ag deposited 
0.2 x 108 
21.6 g]
\end{problembox}

\begin{problembox}
\textbf{Problem 342} \hfill \textbf{[GRB (O.P. Tandon) p.81 Q6]}
\label{q:ext-grb_all_questions-342}

In the process of electroplating, m g of silver is deposited when 
4 ampere of current flows for 2 minutes. The amount (in g) of 
silver deposited by 6 ampere of current flowing for 40 seconds 
will be: 
(a) 4m 
(b)/textasciitilde() 
4m 
(d) 3m 
(c)-
2 
3 
[Hint: 
WI = 
W2 
Q2 
X 
W2 
12 X t2 
m 
4 x 2 x 60 
Wz 
6x 40 
m 
W, =:- g] 
-
2
\end{problembox}

\begin{problembox}
\textbf{Problem 343} \hfill \textbf{[GRB (O.P. Tandon) p.81 Q7]}
\label{q:ext-grb_all_questions-343}

If four moles of electrons are transferred from anode to 
cathode in ail experiment on electrolysis of water, then total 
volume of the two gases produced at STP will be: 
(a) 224 L 
(b) 72.6 L 
(c) 67.2 L 
/textasciitilde()d) 89.4 L 
[Hint: 
Equivalent volume of Hz = 11.2 L 
Equivalent volume of 02 = 5.6 L 
electrolytic cells contammg Ag+, 
ions 
respectively, the deposited Ag (At mass = 108), Ni (At. mass 
= 59) and Cr (At. mass = 52) are: 
Ag 
Ni 
Cr 
-(a )--1-08-g-- --- /textasciitilde()/textasciitilde()/textasciitilde()2-9-,S-g---'---- . ... . $/cdot$$/cdot$$/cdot$$/cdot$$/cdot$'$/cdot$-1$/cdot$$/cdot$$/cdot$$/cdot$ .1;(.\underline{\hspace{1cm}}--
(b) 
108 g 
59 g 
(c) 
108 g 
108 g 
(d) 
108 g 
29.5 g 
[Hint: 
. I 
fA 
108 
EqUlva ent mass 0 
g =: -
= 108 
. 
Equivalent mass ofNi '" 59 =: 29.5 
2 
EquivalentmassofCr = 52 = 17.3 
3 
I 
52 g 
17.3 g 
166 g 
:. Amount of these metals deposited by 1 faraday charge will be: 
108 gAg, 29.5 g Ni and 17.3 g Cr respectively.]
\end{problembox}

\begin{problembox}
\textbf{Problem 344} \hfill \textbf{[GRB (O.P. Tandon) p.81 Q9]}
\label{q:ext-grb_all_questions-344}

Which of the following reactions occur at the cathode during 
the charging of lead storage battery? 
(a) Pb2+ + 2e- -t Pb 
(b) Pb2+ + SO/textasciitilde()- -t PbS04 
(c) Pb -t Fl)+ + 2e-
(d) PbS04 + 2H20-t Pb02 + 4W + SO/textasciitilde()- + 2e-
\end{problembox}

\begin{problembox}
\textbf{Problem 345} \hfill \textbf{[GRB (O.P. Tandon) p.81 Q10]}
\label{q:ext-grb_all_questions-345}

A current of2.6 amp was passed through CUS04 solution for 6 
minutes and 20 seconds. The amount of copper deposited is: 
(a)0.3175g 
(b)0.003Ig (c)635g 
(d)3.l75g 
[II. 
W = ItE = 2.6 x 380 x 31.75"" 0.32 ]. 
lilt: 
96500 
96500 
g
\end{problembox}

\begin{problembox}
\textbf{Problem 346} \hfill \textbf{[GRB (O.P. Tandon) p.81 Q11]}
\label{q:ext-grb_all_questions-346}

In the electrolysis of fused salt, the weight of the substance 
deposited on an electrode will not depend on: 
(a) temperature ofthe bath 
(b) current intensity 
. (c) time of electrolysis 
(d) electrochemical equivalent of the ions
\end{problembox}

\begin{problembox}
\textbf{Problem 347} \hfill \textbf{[GRB (O.P. Tandon) p.81 Q12]}
\label{q:ext-grb_all_questions-347}

When an aqueous solution of sodium chloride is electrolysed 
using platinum electrodes, the ions discharged at the 
electrodes are:
\end{problembox}

\begin{problembox}
\textbf{Problem 348} \hfill \textbf{[GRB (O.P. Tandon) p.82 Q13]}
\label{q:ext-grb_all_questions-348}

How many coulombs are required for the oxidation of I mole
\end{problembox}

\begin{problembox}
\textbf{Problem 349} \hfill \textbf{[GRB (O.P. Tandon) p.82 Q14]}
\label{q:ext-grb_all_questions-349}

ofH20 2 to02? 
(a) 9.65 x 104 C 
(c) 1.93 X 105 C 
(b) 93000 C 
(d) 19.3 x 102 C 
[Hint: 
H 20 2 /textasciitilde() 
02 + 2H" + 2e-
[Hint: 
1 mol H20 2 
2 mol e-
== 2x 96500 C 
== 193000 C 
== 1.93 X 105 C] 
. W '" ItE 
96500 
E 
W x 965JO == 22.2 x 96500 = 59.5 
It 
2 x 5 x 3600 
. 
. 
Atomic. mass 
EqUivalent mass = --/textasciitilde()--- 
Oxidation state 
59.5 == 177 
n 
.;
\end{problembox}

\begin{problembox}
\textbf{Problem 350} \hfill \textbf{[GRB (O.P. Tandon) p.82 Q15]}
\label{q:ext-grb_all_questions-350}

When water is electrolysed,hydrogen and oxygen gases are 
produced. It 1.008 g ofH2 is liberated at cathode, what mass 
of 02 is formed at the anode? 
(a)32g 
(b)16g 
(c)8g 
(d)4g 
[Hint: 
WI 
Wz 
1.008 
1.008 
Wz 
8 
W2 =8g 
where, EI and E2 are equivalent mass of hydrogen and oxygen 
respectively.]
\end{problembox}

\begin{problembox}
\textbf{Problem 351} \hfill \textbf{[GRB (O.P. Tandon) p.82 Q16]}
\label{q:ext-grb_all_questions-351}

The cell potential (E) and the free energy change (/textasciitilde()G) 
accompanying an electrochemical reaction are related by: 
(a) AG == nF log E 
(b) /textasciitilde()G == nFE 
(c) -AG = nFE 
(d) -/textasciitilde()G =: nE. log E
\end{problembox}

\begin{problembox}
\textbf{Problem 352} \hfill \textbf{[GRB (O.P. Tandon) p.82 Q11]}
\label{q:ext-grb_all_questions-352}

The units of conductivity are: 
(a) siemens-1 cm-I 
(b) siemens cm 
(c) siemens cm-1 
(d) siemens cm-2 mol-l 
HI-
The calomel electrode used as a reference electrode contains: 
(a) PbOz-PbS04 mixture 
(b) HgClz 
(c) Hg2Cl2 
(d) ZnCl 2 . 
19./textasciitilde() KCI is used in salt bridge because: 
(a) it forms a good jelly with agar-agar 
(b) it is a strong electrolyte 
(c) it is a good conductor of electricity 
(d) the transference number ofK+ and Cl- ions are almost 
equal 
[Hint: 
Only those electrolytes are used in salt bridge whose 
ions have same transference number in agar-agar gel.]
\end{problembox}

\begin{problembox}
\textbf{Problem 353} \hfill \textbf{[GRB (O.P. Tandon) p.82 Q20]}
\label{q:ext-grb_all_questions-353}

The increase in the equivalent conductance of a salt solution 
on dilution is due to increase in the: 
(a) attraction between the ions 
(b) degree of ionization of the salt 
(c) molecular attraction 
(d) association of the salt
\end{problembox}

\begin{problembox}
\textbf{Problem 354} \hfill \textbf{[GRB (O.P. Tandon) p.82 Q21]}
\label{q:ext-grb_all_questions-354}

When 96500 coulombs of electricity are passed through nickel 
sulphate solution, the amount of nickel deposited will be: 
(a) I mol 
(b) 0.5 mol 
(c) 0.1 mol 
(d).2 mol 
[Hint: 
I faraday deposits 1 equivalent of nickel. 
1 equivalent ofNi = 1/2 mole of nickel.]
\end{problembox}

\begin{problembox}
\textbf{Problem 355} \hfill \textbf{[GRB (O.P. Tandon) p.82 Q22]}
\label{q:ext-grb_all_questions-355}

AI3+(aq.)+3e-/textasciitilde()Al(s); EO==-1.66V 
Cu2+ (aq.) + 2e- /textasciitilde() 
Cu(s);E <> == + 0.34 V 
potentials? 
(a) 1.32 V 
(b) 2.00 V 
[Hint; 
The eell will be: 
(c) 2.30 V 
(d) 4.34 V 
Al(s) I A13+(aq.) II Cuz+(/textasciitilde()q.) I Cu(s) 
E/textasciitilde()ell == E/textasciitilde()athode 
<> 
= Ecu 2+ ICu 
fAI 
= + 0.34 - (-1.66)= + 2.00 V]
\end{problembox}

\begin{problembox}
\textbf{Problem 356} \hfill \textbf{[GRB (O.P. Tandon) p.82 Q23]}
\label{q:ext-grb_all_questions-356}

For which of these oxidation/reduction pairs will the reduction 
potential vary with pH? 
I. AmO/textasciitilde()+ IAmO! 
II. AmO/textasciitilde()+ I Am4+ 
Ill. Am4+ I Am2+ 
(a) I only 
(b) II only 
(c) I and II only 
(d) I, II and III 
(Hi 11 t: 
4H+ + AmO/textasciitilde()+ + 2e- /textasciitilde() 
Am4+ + 2H20 
It includes H+ ions; hence the electrode potential depends on pH.]
\end{problembox}

\begin{problembox}
\textbf{Problem 357} \hfill \textbf{[GRB (O.P. Tandon) p.82 Q24]}
\label{q:ext-grb_all_questions-357}

2Ag + (aq.) + Cu(s) /textasciitilde() 
Cu 2+ (aq.) + 2Ag(s) 
The-standard potential EO for this reaction is 0.46 V. Which-
change will increase the potential the most? 
(a) Doubling the [Ag+] 
(b) Halving the [Cu 2+] 
(c) Doubling the size of the Cu(s) electrode 
(d) Decreasing the size of the Ag electrode by one-half
\end{problembox}

\begin{problembox}
\textbf{Problem 358} \hfill \textbf{[GRB (O.P. Tandon) p.82 Q25]}
\label{q:ext-grb_all_questions-358}

10CI-(aq.)+ 2 Mn0"4 (aq.) + 16W(aq.) 
5Cli(g) 
+ 2Mn 2+(aq.) + 8H20(l) 
. The value of EO for this reaction is 0.15 V. What is the value of 
the equilibrium constant (K ) for this reaction? 
(a) 2.4 x 1025 (b)4.9 X loI2 (c) 1.2 X 105 
(d)3.4 X 102 
Hint: 
K 
antilog [0/textasciitilde():;9J-
n 
[IOXO.l5] 
an 1 o/textasciitilde() 
0.059 
2.4 x 1OZ5 ] 
U. What takes place when zinc metal is added to an aqueous 
solution containing magnesium nitrate and silver nitrate? 
$/bullet$
\end{problembox}

\begin{problembox}
\textbf{Problem 359} \hfill \textbf{[GRB (O.P. Tandon) p.83 Q4]}
\label{q:ext-grb_all_questions-359}

No reaction takes place 
(a) I and 2 only 
(b) I and 3 only 
(c) 1,2 and 3 only 
(d) 4 only
\end{problembox}

\begin{problembox}
\textbf{Problem 360} \hfill \textbf{[GRB (O.P. Tandon) p.83 Q27]}
\label{q:ext-grb_all_questions-360}

In the galvanizing process, iron is coated with zinc. The 
resulting chemical protection is most similar to that provided 
when: 
. 
(a) a magnesium bar is connected to an iron pipe 
. (b) an iron can is plated with tin 
(c) copper pipes are connected using lead s.uIder 
(d) a copper pipe is covered with epoxy paint
\end{problembox}

\begin{problembox}
\textbf{Problem 361} \hfill \textbf{[GRB (O.P. Tandon) p.83 Q28]}
\label{q:ext-grb_all_questions-361}

What is the sign of 110$/circ$ and the value of K for an
\end{problembox}

\begin{problembox}
\textbf{Problem 362} \hfill \textbf{[GRB (O.P. Tandon) p.83 Q33]}
\label{q:ext-grb_all_questions-362}

electrochemical cell for which 
0.80 V? 
(c) 
(d) 
+ 
K 
>1 
<1 
<1 
'Consider a voltaic cell based on these half-cells: 
Ag+(aq.)+e-
Ag(s); EO=tO.80V 
Cd(sj;-eo-;;-(f46/textasciitilde()v' 
Identify the anode and give the voltage of this cell under 
standard conditions. 
(a) Ag; 
= 0.40V 
(c) Cd; 
= 1.20V 
(b) Ag; Eceu = 2.00V 
(d) Cd; Ecell = 2.00V 
[Hint: 
Anode has lower standard reduction potential; thus Cd 
will be considered as anode. 
ECalhode - E Anode 
=E 
-
Ag+ lAg 
led 
, 
= 0.80 
(/textasciitilde()0.40) 1.20 V] 
If the E/textasciitilde()1l for a given reaction has a negative value, which 
gives the correct relationships for the values of 110 0 and K eq ? 
(a) 110$/circ$ > 0; Keq < I 
(b) 110$/circ$ > 0; Keq > 1 
(c) 110$/circ$ < 0; Keq > 1 
(d) 110$/circ$ < 0; Keq < I 
Which of the following solutions is used as an -anti-rusting 
solution? 
(a) Na2S04 
(b) Na3P04 
(c) Na3B03 
(d) Na2S 
The \underline{\hspace{1cm}}pressure of hydrogen gas .is increased from I atm to 100 
atm. Keeping the H+ (I M) constant, the voltage of the 
hydrogen half-cell at 25$/circ$C will be: . 
(a) 0.059 V 
(b) 0.59 V 
(c) 0.0295 V (d) 0.118 V 
Efficiency ofthe following cell is 84/%. 
A(s)+ 
(aq.)/textasciitilde() A 2+(aq.)+B(s); Mi=-285kJ 
Then the standard electrod/textasciitilde() potential of the cell will be: 
(a) 1.20 V 
(b) 2.40 V 
(c) 1.10 V 
(d) 1.24 V 
110$/circ$ 
-nFEo 
[Hint: 
Efficiency 
MfO 
MfD 
0.84 = \underline{\hspace{1cm}} 2 x EO x 96500 
-285 x 1000 
eo =+ 1.24 V]
\end{problembox}

\begin{problembox}
\textbf{Problem 363} \hfill \textbf{[GRB (O.P. Tandon) p.83 Q34]}
\label{q:ext-grb_all_questions-363}

A chemical reaction will be spontaneous if: 
(a) 
is positive 
(b) 110$/circ$ is negative 
(c) reaction quotient, Q < K ( d) E/textasciitilde()xid is negative
\end{problembox}

\begin{problembox}
\textbf{Problem 364} \hfill \textbf{[GRB (O.P. Tandon) p.83 Q35]}
\label{q:ext-grb_all_questions-364}

Match the List-I with List-II: 
List-I (Electrode) 
List-II (Type)
\end{problembox}

\begin{problembox}
\textbf{Problem 365} \hfill \textbf{[GRB (O.P. Tandon) p.83 Q1]}
\label{q:ext-grb_all_questions-365}

Calomel 
(A) Reference
\end{problembox}

\begin{problembox}
\textbf{Problem 366} \hfill \textbf{[GRB (O.P. Tandon) p.83 Q3]}
\label{q:ext-grb_all_questions-366}

Hydrogen 
(C) Membrane
\end{problembox}

\begin{problembox}
\textbf{Problem 367} \hfill \textbf{[GRB (O.P. Tandon) p.83 Q4]}
\label{q:ext-grb_all_questions-367}

Quinhydrone 
(D) Gas 
Codes: 
(a) I-A 
2-;C 
3-D 
4-B 
(b) I-B 
2-A 
3-D 
4-C 
(c) I-C 
2-B 
3-A 
4-D 
(d) I-D 
2-A 
3-C 
4-B 
mho 
and its equivalent conductance is 1.53 
mho cm2 eq-I. The Ksp for BaS04 will be: 
(a) 4 X 10-12 M 
(b) 4 X 106 M -
4 X 10-6 M2
\end{problembox}

\begin{problembox}
\textbf{Problem 368} \hfill \textbf{[GRB (O.P. Tandon) p.83 Q37]}
\label{q:ext-grb_all_questions-368}

The standard reduction potential of hydrogen is zero because: 
(a) it is assumed 
(b) hydrogen is easiest to oxidise 
(c) hydrogen has single electron 
(d) hydrogen is electronegative 
In the following three questions, three statements I, II and III are 
given. Mark: 
(a) if all the statements are correct 
(b) if II and III are correct-
(c) ifIarid III are correct 
( d) if only II is correct 
IBHU (Mains) 200'll
\end{problembox}

\begin{problembox}
\textbf{Problem 369} \hfill \textbf{[GRB (O.P. Tandon) p.83 Q38]}
\label{q:ext-grb_all_questions-369}

I. Conductance of electrolyte solution increases with 
temperature. 
II. Resistivity is reciprocal of molar conductivity of 
electrolyte. 
III. Cell constant has unit cm -1. 
(a) 
(b) 
(c) 
(d)
\end{problembox}

\begin{problembox}
\textbf{Problem 370} \hfill \textbf{[GRB (O.P. Tandon) p.83 Q39]}
\label{q:ext-grb_all_questions-370}

1. The conductivity of molten NaCI is due to movement of 
Na+ and CI - ions. 
II. Solid NaCI is also conductor of electricity. 
III. Molten sodium is a good conductor because of mobile 
electrons. 
(a) 
(b) 
(c) 
(d)
\end{problembox}

\begin{problembox}
\textbf{Problem 371} \hfill \textbf{[GRB (O.P. Tandon) p.83 Q40]}
\label{q:ext-grb_all_questions-371}

1. Cathode is -ve terminal both III electrochemical and 
electrolytic cells. 
. 
II. Reduction occurs at cathode both in galvanic as well as 
electrolytic cells. 
III. Chemical charge in electrolytic cell is$/cdot$non-spontaneous. 
(a) 
(b) 
(c) 
(d)
\end{problembox}

\begin{problembox}
\textbf{Problem 372} \hfill \textbf{[GRB (O.P. Tandon) p.84 Q41]}
\label{q:ext-grb_all_questions-372}

In an experiment, 0.04 F was passed through 400 mL of a 1M 
solution of NaCI. What would be the pH of the solution after 
the electrolysis? 
[PMT (Kerala) 2007) 
(a) 8 
$/bullet$ (b) 10 
(c) 13 
(d) 6 
(ej 9 
[Hint: 
Electrolysis of aq. NaCI gives hydrogen gas at 
cathode $/cdot$.and oxygen gas at anode, the electrolyte solution 
contains NaOH after electrolysis. 
Number of equivalents of NaOH formed' = 0.04 
Normality, 
N 
0.04 x 1000 = 0.1 
. :' 400 
10a-] O'.lM 
.' pOH.= I 
:. pH = 13]
\end{problembox}

\begin{problembox}
\textbf{Problem 373} \hfill \textbf{[GRB (O.P. Tandon) p.84 Q42]}
\label{q:ext-grb_all_questions-373}

An alloy of Pb/textasciitilde()Ag weighing 1.08 g was dissolved in dilute 
and the volume made to 100 mL. A silver electrode was 
is 0.80 V, what is the percentage of Ag in the alloy? 
(At 25$/circ$C, RT / F 
0.06) 
(PET (Kerala) 2001) 
(a) 25 
(b) 2.50 
(c) 10 
(d) I 
(e) 50 
lfIint: Overall cell reaction is: 
H2(g) + 2Ag+ 
2Ag(s) + 2H+(aq.) 
0.06 X 2.303 I 
[H+ ]2 
og 
2 
[Ag+]2 pH2 
0.62 = 0.80 + 2 x 0.06 x 2.303 log [Ag+] 
2 
[Ag+]= 0.05 M 
Number of moles of Ag + in 100 mL 
MV 
1000 
--/textasciitilde()=0.005 
Mass of silver = 0.005 x 108 g 
. 
0.005 x 108 x 100 
Percentage of Ag m 1.08 g of alloy = 
= 50/%] 
1.08
\end{problembox}

\begin{problembox}
\textbf{Problem 374} \hfill \textbf{[GRB (O.P. Tandon) p.84 Q43]}
\label{q:ext-grb_all_questions-374}

Select the correct statements about dry cell: 
(a) It is also called Lec1anche cell 
. (b) It is also called Daniell cell 
(c) Electrolyte used is moist paste ofNH4Cl and ZnC12 
(d) Cathodic process is: 
2Mn02(s) + 2NH;t(aq.) + 2e- /textasciitilde() Mn20is) + 2NH3(g)
\end{problembox}

\begin{problembox}
\textbf{Problem 375} \hfill \textbf{[GRB (O.P. Tandon) p.84 Q44]}
\label{q:ext-grb_all_questions-375}

Given the standard oxidation potentials, 
. 
+0.4 V 
2+ 
-0.8 V 
3+ 
Fe --7 Fe 
(aq.) --7 Fe 
(aq.) 
+0.9V 
0.6 V 
Fe --7 Fe(OH)2 --7 Fe(OH)3 
It is easier to oxidise Fe2+ to Fe3+ in: 
(a) acid medium 
(b) alkaline medium 
(c) neutral medium 
(d) both in acidic and alkaline mediums 
+ H2$/circ$(l) 
[Hint: 
E'p 2+ p 3+ is positive in alkaline solution, therefore, it is 
e Ie 
.. 
easier to oxidise 
to Fe3+ in alkaline medium.]
\end{problembox}

\begin{problembox}
\textbf{Problem 376} \hfill \textbf{[GRB (O.P. Tandon) p.84 Q45]}
\label{q:ext-grb_all_questions-376}

Dipping iron article into a strongly alkaline solution of sodium 
phosphate: 
. 
(VIJ'EEE 2008) 
(a) does not affect the article. 
(b) forms Fe203 $/cdot$J;H20 on the surface 
(c) forn:s iron phosphate film 
(d) forms ferric hydroxide
\end{problembox}

\begin{problembox}
\textbf{Problem 377} \hfill \textbf{[GRB (O.P. Tandon) p.84 Q46]}
\label{q:ext-grb_all_questions-377}

For the redox process, 
'47. 
Zn(s)+ Cu2
rl-
Zn 2+ + Cu(S)E/textasciitilde()ell /textasciitilde()+ 1.l0V 
which graph correctly represents Ecell (Y-axis) as a f.Unction of 
[Zn 2+] 
log -- (X-axis)? 
[Cu 2+] 
1.10V 
-1.0 
o 
+ 1.0 
-1.0 
o 
+ 1.0 
---.-- . \underline{\hspace{1cm}} .\underline{\hspace{1cm}}---
-/textasciitilde()------ ......... ---------/textasciitilde().--.-- ............. ---/textasciitilde() 
(c) 
(d) 
-1.0 . 
1.1OV 
" 
o 
.+ 1.0 
-1.0 
o 
+ 1.0 
A fuel cell involves combustion of the butane at i atm and 
298K 
13 
C4HIO (g)+ 
02(g)----> 4C02(g) + 5H20(l) 
. 2 , 
fj.Go = - 2746 kJ / mol 
what is EO of a cell ? 
(a) + 4.74 V 
(b) + 0.547 V (c) + $/cdot$1.09 V 
(d) + 4.37 V 
[Hint: In the reaction 
.. 10+ 10 
13 
+ 4 - 4 
. 
C4 H IO (g) + 2 .02(g)/textasciitilde()4CQ2(g)+ 5H20(l) 
Change in oxidation number of carbon = + 16 (- 10) 
+ 26 
:. Number of electrons involved in cell process will be 26. 
EO = \underline{\hspace{1cm}}fj.Go 
(- 2746) x 1000 
nF 
26x 96500 
=+ L09V]
\end{problembox}

\begin{problembox}
\textbf{Problem 378} \hfill \textbf{[GRB (O.P. Tandon) p.84 Q48]}
\label{q:ext-grb_all_questions-378}

Molar conductance Am is plotted agains.JC (mol litre-I) for 
three electrolytes (NaCl, Hel, NH40H) 
----(3)
\end{problembox}

\begin{problembox}
\textbf{Problem 379} \hfill \textbf{[GRB (O.P. Tandon) p.85 Q49]}
\label{q:ext-grb_all_questions-379}

In the concentration cell 
HCI 
NaCI 
NaCl 
HC! 
3 
NH/textasciitilde()OH 
NH40H. 
HC! 
NaC! 
Pt(H )iHA (O.lM) iiHA (1M) I 
(H ) Pt 
.2 
NaA (1M) 
NaA (lM) 
2 
(PKa ofHA = 4) 
Cell potential will be: 
(a) 0.03 V 
(b) 0.06 V 
(c) - 0,06 V 
(d)' 0.03 V 
pH Anode 
3 
'pHCathode 
4 (fromeq. I) 
Eccll 
- 0.06 V] 
50;--In the following-proeess'ofdisproportionation:' 
E/textasciitilde()104/CI03 + 0.36 V 
2CIO)' 
Chlorate 
CIO; + CIO:;. 
+ 0.33 V 
Per chlorate 
ion 
ion 
[,4' ......... 
: ..... /textasciitilde() 
.... 
. 
-
.-
1- (c)
\end{problembox}

\begin{problembox}
\textbf{Problem 380} \hfill \textbf{[GRB (O.P. Tandon) p.85 Q50]}
\label{q:ext-grb_all_questions-380}

(d) , 
(c) 
(c) 
(b) 
(a) 
(d) 
(b) 
Initial concentration of chlorate ion was 0.1 M, The 
equilibrium concentration of per chlorate ion will be : 
(a) 0,19 V 
(b) OJ M 
(c) 0,024 M 
(d) 0.019 M 
[Hint: 
CIO;(aq.) + H20(1) -----:-l- CIO:;(aq.) + 2H+ (aq.) + 2e-
2H+ + CIO; (aq;) + 2e- -----:-l- ClO; (aq.) + H20(/) 
ZCIO; (aq,) /textasciitilde() 
CIO; (aq.) + CIO:; (aq.) 
= 0.33 - 0.36 = - 0.03 V 
E - EO 
0.059 I' Q 
-
--- og 
n 
At equilibrium, 
E 
0, 
n = 2, Q = K 
o 
0.03 - 0,059 log K 
2 
2CIO/textasciitilde() (Jq.)'/textasciitilde() ClO; (aq.) + C10:; (aq.) 
0.1 
o 
o 
0.1 
2x 
x 
x 
K= 
xXx 
=/textasciitilde() 
. 
-(OX:':'-2X)2/textasciitilde()10'--
... 
x=0.019]
\end{problembox}

\begin{problembox}
\textbf{Problem 381} \hfill \textbf{[GRB (O.P. Tandon) p.86 Q1]}
\label{q:ext-grb_all_questions-381}

How many grams of water will be electrolysed by 96500 
coulomb charge?
\end{problembox}

\begin{problembox}
\textbf{Problem 382} \hfill \textbf{[GRB (O.P. Tandon) p.86 Q2]}
\label{q:ext-grb_all_questions-382}

Number offaradays required to convert 1 mole of Cr20/textasciitilde()- into
\end{problembox}

\begin{problembox}
\textbf{Problem 383} \hfill \textbf{[GRB (O.P. Tandon) p.86 Q4]}
\label{q:ext-grb_all_questions-383}

The ratio Of( /textasciitilde(): J for Ca3 (P04 h will be equal to :
\end{problembox}

\begin{problembox}
\textbf{Problem 384} \hfill \textbf{[GRB (O.P. Tandon) p.86 Q6]}
\label{q:ext-grb_all_questions-384}

If an aqueous solution of NiCI is electrolysed using platinum 
electrode by a current of 5 amp, then what volume of Cl2 gas 
in litres at STP will be produced? 
7, Charge of 6.24 x Hi8 electrons will be (in coulomb):
\end{problembox}

\begin{problembox}
\textbf{Problem 385} \hfill \textbf{[GRB (O.P. Tandon) p.86 Q8]}
\label{q:ext-grb_all_questions-385}

A current of 2 amp-when passed for 5 hour through a molten 
salt deposits 22.2 g of metal of atomic mass 177. The positive 
oxidation state of the metal in the metal salt is :
\end{problembox}

\begin{problembox}
\textbf{Problem 386} \hfill \textbf{[GRB (O.P. Tandon) p.86 Q9]}
\label{q:ext-grb_all_questions-386}

. 
Cr(s) I Cr3+ II Fe2+ r Fe(s) 
In above cell, the value of n in theNemst equation: 
i.e., E = E"- o.05910gloQ will be: 
n
\end{problembox}

\begin{problembox}
\textbf{Problem 387} \hfill \textbf{[GRB (O.P. Tandon) p.86 Q10]}
\label{q:ext-grb_all_questions-387}

In the Nernst Equation, 
Q will be equal to the equilibrium constant Kc, when the cell 
potential E is equal to :
\end{problembox}

\begin{problembox}
\textbf{Problem 388} \hfill \textbf{[GRB (O.P. Tandon) p.86 Q5]}
\label{q:ext-grb_all_questions-388}

12 (s) / r (0.1 M) half-cell is connected to W (aq) /(H2 1 atm) 
lla/textasciitilde()If-/textasciitilde()c(:lLaI!(U:tsc/textasciitilde()R1!9tenti/textasciitilde()lJ  found toJ,!e 0.7714VJ( 
0,535 V, the pH of W / Hi half-cell will be : 
.. -.-.... -----.. -/textasciitilde()/textasciitilde()- -/textasciitilde()------- .--.. ------.----,-.--/textasciitilde()., .. -.--.-.. /textasciitilde(), .. --.----/textasciitilde() ----
\end{problembox}

\begin{problembox}
\textbf{Problem 389} \hfill \textbf{[GRB (O.P. Tandon) p.87 Q1]}
\label{q:ext-grb_all_questions-389}

Equivalent mass of sulphuric acid in lead storage battery is: 
(a) 49 
(b) 98 
(c) 24.5 
(d) none of these
\end{problembox}

\begin{problembox}
\textbf{Problem 390} \hfill \textbf{[GRB (O.P. Tandon) p.87 Q2]}
\label{q:ext-grb_all_questions-390}

Normalities of sulphuric acid before and after discharge are: 
(a) 5.15, 2.32 (b) 2.32; 5.15 (c) 5.15, 5.15 (d) 2.32, 2.32 
(c) 265.04 amp-hrs 
(d) 26.504 amp-hrs
\end{problembox}

\begin{problembox}
\textbf{Problem 391} \hfill \textbf{[GRB (O.P. Tandon) p.87 Q4]}
\label{q:ext-grb_all_questions-391}

The amount of charge which the battery must have been used 
is: 
(a) 9.88 F 
(b) 8.98 F 
(c) 8.89 F 
(d) 7.88 F 
50$/cdot$$/cdot$ Which oHhe$/cdot$foU(')wing-takes$/cdot$placeindischarge-proeess at$/cdot$ 
anode? 
(a) PbOz + 4H+ + SO/textasciitilde()- + 2e-
PbS04 + 2H20. 
(b) PbS04 + 2H20/textasciitilde() Pb02 + 4W + SO/textasciitilde()- + 2e-
(c) Pb + SO/textasciitilde()- /textasciitilde() 
PbS04 + 2e-
(d) PbS04 + 2e/textasciitilde() /textasciitilde() 
Pb + SO!-
$/bullet$ Passage2 
Electrolysis is the process in which electrical energy is converted 
to chemical energy. In electrolytic cell, oxidation. takes place at 
anode and reduction at cathode. Electrode process depends on the 
electrode taken for electrolysis. Amount of substance liberated at an 
electrode is directly proportional to the amount of charge passed 
through it. The mass oj substance liberated at electrode is calculated 
using the follOWing relation: 
ItE 
m=--
96500 
Here, E represents the equivalent mass and 96500 C 'is called the 
Faraday constant. Faraday (96500C ) is the charge of I mole 
electron, i.e., 6.023 X 1023 electrons; it is used to liberate one gram 
equivalent of the substance. 
Answer the following questions:
\end{problembox}

\begin{problembox}
\textbf{Problem 392} \hfill \textbf{[GRB (O.P. Tandon) p.87 Q1]}
\label{q:ext-grb_all_questions-392}

Th/textasciitilde() platinum electrodes were immersed in a'solution of cupric 
sulphate (CUS04) and electric current is passed through the 
solution. After sometime, it was observed that the colour of 
copper sulphate disappeared with evolution of a gas at the 
electrode. The colourless solution contains: 
. (a) platinum sulphate 
(b) copper nitrate 
(c) copper sulphate 
(d) sulphuric acid
\end{problembox}

\begin{problembox}
\textbf{Problem 393} \hfill \textbf{[GRB (O.P. Tandon) p.87 Q2]}
\label{q:ext-grb_all_questions-393}

The passage of current liberates Hz at cathode and Cl 2 at 
anode. The solution is: 
(a) copper chloride in water (b) NaCl in water 
(c) mercuric chloride in water(d) AuCl 3 in water
\end{problembox}

\begin{problembox}
\textbf{Problem 394} \hfill \textbf{[GRB (O.P. Tandon) p.87 Q3]}
\label{q:ext-grb_all_questions-394}

On -electrolysis of dilute sulphuric acid using platinum 
electrodes, the product obtained at the anode will be: 
(a) hydrogen 
(b) oxygen 
(c) hydrogen sulphide 
(d) sulphur oxide
\end{problembox}

\begin{problembox}
\textbf{Problem 395} \hfill \textbf{[GRB (O.P. Tandon) p.87 Q4]}
\label{q:ext-grb_all_questions-395}

How many faradays are required to reduce 1 mol BrO) to 
Br-? 
003 
005 
/textasciitilde()6 
004
\end{problembox}

\begin{problembox}
\textbf{Problem 396} \hfill \textbf{[GRB (O.P. Tandon) p.87 Q5]}
\label{q:ext-grb_all_questions-396}

$/cdot$Calculate the volume of gas liberated at the an9de at STP 
during the electrolysis of a CuSO 4 solution by a cqrrent of I A 
passed for 16 minutes arid 5 seconds:)' 
(a) 224/mL 
(b) 56 mL 
(c) 112 mL 
(d) 448'mL 
. 
I 
,-
+ -
+ 2e-
v 
:= 1 x 965 x 5.6 x 10
3 
:= 56 mL] 
96500 
96500
\end{problembox}

\begin{problembox}
\textbf{Problem 397} \hfill \textbf{[GRB (O.P. Tandon) p.87 Q6]}
\label{q:ext-grb_all_questions-397}

The quantity of electricity required to liberate 112 cc hydrogen 
. .!l.LS,T.:efro.macjdifi/textasciitilde()d\underline{\hspace{1cm}}.w.l!/textasciitilde()rj/textasciitilde(); ... \underline{\hspace{1cm}} .......... \underline{\hspace{1cm}} .. \underline{\hspace{1cm}} .... \underline{\hspace{1cm}} .... \underline{\hspace{1cm}} .. \underline{\hspace{1cm}} \underline{\hspace{1cm}} 
(Comed (Karnataka) /%008] 
(a) 965 C 
(b) 9650 C 
(c) 96500 C 
(d) 4825,C 
$/bullet$ Passage 3 
The concentration of potassium ions inside a biological cell is at 
least twenty times higher than the outside. The resulting potential 
difference across the cell'is important in several processes such as 
transmission of nerve impulses and maintaining the ion balance. A 
simple model for such a concentration cell involving a metal Mis: 
M(s) IM+ (aq. 0.05 molar) IIM+ (aq; I molar) I M(s) 
For the above electrolytic cell, the magnitude of the cell potential 
IEeell 1=70mV, 
Answer the following questions:
\end{problembox}

\begin{problembox}
\textbf{Problem 398} \hfill \textbf{[GRB (O.P. Tandon) p.87 Q1]}
\label{q:ext-grb_all_questions-398}

For the above cell 
(a) Eeell <: 0; IJ.G > 0 
(b) EeelJ > 0; IJ.G < 0 
(c) Eeell < 0; IJ.Go> 0 
(d) Eeell > 0; M?< 0
\end{problembox}

\begin{problembox}
\textbf{Problem 399} \hfill \textbf{[GRB (O.P. Tandon) p.87 Q2]}
\label{q:ext-grb_all_questions-399}

If the 0.05 molar solution of M+ is replaced by a 0.0025 molar 
M+ solution, then the magnitude of the cell potential would 
be: 
(a) 35 mV 
(c) 140 mV 
(b) 70mV 
(d) 700 mV 
(lIT 20101 
!Hint: 
1 (b) Electrolyte concentration cell will be spontaneous 
when the concentration in cathodic half cell is greater than that of 
anodic half cell. Thus, the given cell is spontaneous hence 
ll.G < 0, Ecell > 0
\end{problembox}

\begin{problembox}
\textbf{Problem 400} \hfill \textbf{[GRB (O.P. Tandon) p.87 Q2]}
\label{q:ext-grb_all_questions-400}

(c) E;ell = 0, for every concentration cell 
E = 0 0.059 log [M+ J.node 
n 
[M+ lathode 
= \underline{\hspace{1cm}} 0.059 Iog[0.0025] 
1 
= +153 mV 
It is close to 140 mV.] .
\end{problembox}

\begin{problembox}
\textbf{Problem 401} \hfill \textbf{[GRB (O.P. Tandon) p.88 Q1]}
\label{q:ext-grb_all_questions-401}

On the basis of information available for the reaction: 
4 
2 
- Al + 02 /textasciitilde() 
- Al20 3; I1G = - 827 kJ I mol of 02 
-/textasciitilde()-- - --
- -3- ---
--- ---- -
the minimum emf required to carry out an electrolysis of 
Al20 3 is: 
(Given: 1 F = 96500 C) 
(a) 2.14 V 
(b) 4.28 V 
(c) 6.42 V 
(d) 8.56 V 
[Hint: 
Al /textasciitilde()AI3+ + 3e-
i. e., 
4 
4 
\underline{\hspace{1cm}} 
- mol Al == - x 3 mol e 
3 
3 
== 4 mol e-
n=4 
t:..G =-nFE 
-827 x 1000 = - 4 x 96500 x E 
E=2.14 V]
\end{problembox}

\begin{problembox}
\textbf{Problem 402} \hfill \textbf{[GRB (O.P. Tandon) p.88 Q2]}
\label{q:ext-grb_all_questions-402}

The equilibrium .constant Kc will be equal to Q, when: 
(a)E=Eo 
(b)RTlnF=1 
(c) E = 0 
(d) EO = 1
\end{problembox}

\begin{problembox}
\textbf{Problem 403} \hfill \textbf{[GRB (O.P. Tandon) p.88 Q3]}
\label{q:ext-grb_all_questions-403}

The nature of graph of E/textasciitilde()ell against log Kc is alan: 
(a) straight line 
(b) parabola 
(c) hyperbola 
(d) elliptical curve
\end{problembox}

\begin{problembox}
\textbf{Problem 404} \hfill \textbf{[GRB (O.P. Tandon) p.88 Q4]}
\label{q:ext-grb_all_questions-404}

The equilibrium constant K c for the reaction: 
Cu(s) + 2Ag+ (aq.) /textasciitilde() 
Cu2+ (aq.) + 2Ag(s) (/textasciitilde()ell = 0.46 V) 
will be: 
(a) antilog 15.6 
(c) antilog 1.5 
(b) antilog 2.5 
(d) antilog 12.2
\end{problembox}

\begin{problembox}
\textbf{Problem 405} \hfill \textbf{[GRB (O.P. Tandon) p.88 Q5]}
\label{q:ext-grb_all_questions-405}

n - lor" tJ.te electrochemical cell, 
Zn(s) I Z.' 2+ (1 M) aq.1I Cu 2+(1 M) aq.1 Cu(s) is 1.10 volt at 
25$/circ$C. 
The equilibrium con/textasciitilde()tant for the cell reaction: 
will be: 
(a) 10-37 
Zn(s) + Cu2+ (aq.) /textasciitilde() 
Cu(s) +. Zn2+ (aq.) 
(b) 1037 
(c) 10-39 
(d) 1039 
concentration: 
Am = A: - bFc 
where, A: = molar conductance at infinite dilution 
-:-c-=-concentration of-tile soltIfioff-------
b = constant 
The degrees of dissociation of weak electrolytes are calculated 
as: 
Answer the following questions:
\end{problembox}

\begin{problembox}
\textbf{Problem 406} \hfill \textbf{[GRB (O.P. Tandon) p.88 Q1]}
\label{q:ext-grb_all_questions-406}

Which of the following decreases on dilution of electrolyte 
solution? . 
(a) Equivalent conductance . (b) Molar conductance 
(c) Specific conductance 
. (d) Conductance
\end{problembox}

\begin{problembox}
\textbf{Problem 407} \hfill \textbf{[GRB (O.P. Tandon) p.88 Q2]}
\label{q:ext-grb_all_questions-407}

The correct order of equivalent conductances at infinite 
dilution ofLiCl, NaCi and KCl is: 
(a) LiCI > NaCI > KCl 
(b) KCl > NaCI > LiCI 
(c) NaCI > KCl > LiCI 
(d) LiCl > KCl > NaCI
\end{problembox}

\begin{problembox}
\textbf{Problem 408} \hfill \textbf{[GRB (O.P. Tandon) p.88 Q3]}
\label{q:ext-grb_all_questions-408}

For which of the following electrolytic solutions Ani and Ae 
are equal? 
. 
(a) BaCl2 
(b) KCl 
(c) A12(S04h (d) CaCl2
\end{problembox}

\begin{problembox}
\textbf{Problem 409} \hfill \textbf{[GRB (O.P. Tandon) p.88 Q4]}
\label{q:ext-grb_all_questions-409}

The conductance of a solution of an electrolyte is equal to that 
of its specific conductance. The cell constant of the 
conductivity cell is equal to: 
. (a) resistance 
(b) faraday 
(c) zero 
(d) unity
\end{problembox}

\begin{problembox}
\textbf{Problem 410} \hfill \textbf{[GRB (O.P. Tandon) p.88 Q5]}
\label{q:ext-grb_all_questions-410}

Which of the following equality holds good for the strong 
electrolytes? 
(a) A = AD asc-71 
(c) A = AD as c -700 
$/bullet$ Passage 6 
(b) A = AD as c -70 
(d)A=Ao asc-7Jb 
At . infinite dilution, when the dissociation of electrolyte is 
complete, each ion makes a definite contribution towards the molar 
conductance of electrolyte, irrespective of the nature of the other ion 
with which it is associated.
\end{problembox}

\begin{problembox}
\textbf{Problem 411} \hfill \textbf{[GRB (O.P. Tandon) p.89 Q2]}
\label{q:ext-grb_all_questions-411}

Whi/textasciitilde()h of the following oxides will be thermally most stable? 
(a)Zno. 
(b) MgO 
(c) CU20 
(d) Ag20
\end{problembox}

\begin{problembox}
\textbf{Problem 412} \hfill \textbf{[GRB (O.P. Tandon) p.89 Q3]}
\label{q:ext-grb_all_questions-412}

Which of the following reactions is not correct? 
(a) 2;n + H2S04 /textasciitilde() 
ZnS04 + H2 
(b) Fe + H2S04 /textasciitilde() 
FeS04 + H2 
Answer th$/cdot$e following questions: 
(c) Mg + H2S04 /textasciitilde() 
MgS04 + H2
\end{problembox}

\begin{problembox}
\textbf{Problem 413} \hfill \textbf{[GRB (O.P. Tandon) p.89 Q1]}
\label{q:ext-grb_all_questions-413}

The ionic conductances of Al3+ and SO/textasciitilde()- ions at infmite 
(d) Cu + H2S04 /textasciitilde() 
C/textasciitilde()S04 + H2 
dilution are x and ; $/Omega$ -I cm2 mol-I respectively. If
\end{problembox}

\begin{problembox}
\textbf{Problem 414} \hfill \textbf{[GRB (O.P. Tandon) p.89 Q4]}
\label{q:ext-grb_all_questions-414}

Which of the following couples will have highest value of 
Kohlrausch's iaw is valid, then molar conductance of 
emf? 
. 
. 
aluminium sulphate at infinite dilution will be: 
(a) Mg I Mg2+ II Ag+ / Ag 
(b) Zn I Zn 2+ II Cu 2+ I Cu 
(a) 3x + 2y(b)3y+ 2X 
(c).a + 2y 
(d)3x + 3y 
. (c) Zn I Zn2+11 Ag+ / Ag 
(d) Cu I Cu 2+ II Ag+ lAg 
------/textasciitilde()--------------/textasciitilde()------------------------/textasciitilde()----/textasciitilde()----------------/textasciitilde()------------------------/textasciitilde()/textasciitilde()--/textasciitilde()------/textasciitilde()-----+2.--Thelllotar/textasciitilde()onUrr/textasciitilde()c/textasciitilde()rorinfinrt/textasciitilde()onfmelectrm/textasciitilde()A/textasciitilde()--------/textasciitilde()/textasciitilde()/textasciitilde()/textasciitilde()--/textasciitilde()./textasciitilde()/textasciitilde()/textasciitilde()--/textasciitilde()/textasciitilde()/textasciitilde()/textasciitilde()/textasciitilde()/textasciitilde()/textasciitilde()/textasciitilde()/textasciitilde()---------
/textasciitilde() 
/textasciitilde() 
and .CA are 140 and 120 $/Omega$-I cm2 mol-I. If the molarS. Whichof.the following metals will not displace hydrogen 
conductance at infinite dilution ofBX is 198 $/Omega(-1)$ cm2 mol-I, 
from water? 
. 
then at infinite dilution, the molar conductance ofCX is: 
(a) 178 
(b) 198 
(c) 218 
(d) 130 
= 120 + 198 -140 = 178]
\end{problembox}

\begin{problembox}
\textbf{Problem 415} \hfill \textbf{[GRB (O.P. Tandon) p.89 Q3]}
\label{q:ext-grb_all_questions-415}

The molar conductance of 0.001 M acetic acid is 50 $/Omega$-I 
cm2 mol-I. The maximum value of molar conductance of 
acetic acid is 250 $/Omega$-I cm2 mol-I. What is the degree of 
dissociation (a ) of acetic acid?$/cdot$ 
(a)0.5 
(b) 0.2 
(c) 0.3 
(d) 0.4
\end{problembox}

\begin{problembox}
\textbf{Problem 416} \hfill \textbf{[GRB (O.P. Tandon) p.89 Q4]}
\label{q:ext-grb_all_questions-416}

Which of the follow.ing solutions will have highest value of 
the molar conductance of CH3COOH? 
(a) 1 MCH 3COOH 
(b) 0.5 MCH 3COOH 
(c) 0.3 MCH 3COOH 
(d) 0.1 MCH 3COOH
\end{problembox}

\begin{problembox}
\textbf{Problem 417} \hfill \textbf{[GRB (O.P. Tandon) p.89 Q5]}
\label{q:ext-grb_all_questions-417}

The Unit of molar conductance of an electrolyte solution will 
. be.: 
(a) $/Omega$-I cm2 mol/textasciitilde()1 
(c) S cm2 mol-I 
$/bullet$ Passage 7 
(b) mho cm2 mo1-1 
(d) $/Omega(-1)$ cm-1 mol-I 
. The potential associaied with each electrode is known as 
el(!Ctrode potential. If the concentration of each speCies taking part 
in the electrode reaction is unity (if any appears in the electrode 
reaction, it is confined to 1 atmospheric pr.essure) and further the 
reaction is carried out at 298 K, then the potential of each electrode 
is said to be the standard electrode poteritial.By convention, the 
standard electrode potential of hydrogen electrode is 0.0 volt. The 
electrode potential value for each electrode process is a measure of 
relative tendency of the active species in the process to remain in the 
oxidisedlreducedform. A negative go means that the redox couple.is 
a stronger redUcing agent than the H+ / H 2 couple. A positive E 0 
means that the redox couple isa weaker reducing agent than the 
H+ / H 2 couple. The meial with greater positive value of standard 
reduc.tion potential forms the oxide of greater thermal stability. 
Answer the following questions: 
. 1. Given the standard reduction potentials, 
E/textasciitilde()+ IK = - 2.93 V, E:g+ lAg = + 0.80 V, E/textasciitilde()g2+ IHg = 0.79 V 
(a) Mg$/cdot$ 
(b) Zn 
(c) Sn 
(d) Ag 
$/bullet$ Passage 8 
Chemical reactions involve interaction of atoms and molecules, A 
large number of atoms/molecules (approximately 6.023 x 1023) are 
present in a few grams of any chemical compound varying with their 
atomic/molecular masses. 
To 
handle such 
large 
numbers 
conveniently, ,the mole concept was $/cdot$introduced. This concept has 
implications in diverse areas such as analytical chemistry,$/cdot$ 
biochemistry, electrochemistry and radiochemistry. The following 
example 
illustrates 
a 
typical 
case, 
involving 
chemical/electrochemical 
reaction, 
which 
requires 
a 
clear 
understanding of the mole concept.$/cdot$ 
A 4.0 molar aqueous solution of NaCI is prepared and 500 mL of 
this solution is electro lysed. This leadS to the evdlution of chlorine 
gas at one of the electrodes (Atomic mass: Na = 23; Hg = 200; 
1 Faraday = 96500 coulombs). . 
". (lIT 2007) 
Answer the following $/cdot$questions:
\end{problembox}

\begin{problembox}
\textbf{Problem 418} \hfill \textbf{[GRB (O.P. Tandon) p.89 Q1]}
\label{q:ext-grb_all_questions-418}

The total$/cdot$number of moles of chlorine gas evolved is: 
(a) O.S 
(b) 1.0 
(c) 2.0 
(d) 3.0 
[Hint: 
Number of moles of NaCl = MV = 4 x 500 = 2 
. 
1000 
1000 
2Cl-:- /textasciitilde() 
C12 + 2e-
2 mol Cl- ions give 1 mol C12]
\end{problembox}

\begin{problembox}
\textbf{Problem 419} \hfill \textbf{[GRB (O.P. Tandon) p.89 Q2]}
\label{q:ext-grb_all_questions-419}

If the cathode is a Hg electrode, the maximum weight (in g) of 
amalgam formed from this solution is: 
(a) 200 
(b) 225 
(c) 400 
(d) 446 
[Hint: 
Electrolysis gives 2 mol sodium at cathode. Thus, $/bullet$ 
amalgam (Na / Hg) will contain 2 mol of each sodium and Hg. 
Mass of amalgam = 2 x 200 + 2 x 23 = 446]
\end{problembox}

\begin{problembox}
\textbf{Problem 420} \hfill \textbf{[GRB (O.P. Tandon) p.89 Q3]}
\label{q:ext-grb_all_questions-420}

The total charge (in coulombs) required for complete 
electrolysis is: 
(a) 24125 
(b) 48250 
(c) 96500 
(d) 193000 
fRint: 
2 mol electrons will be required, therefore, required 
charge will be 2 faraday or 193000 coulombs.]
\end{problembox}

\begin{problembox}
\textbf{Problem 421} \hfill \textbf{[GRB (O.P. Tandon) p.90 Q5]}
\label{q:ext-grb_all_questions-421}

(a, b, c) 
Passage 7.
\end{problembox}

\begin{problembox}
\textbf{Problem 422} \hfill \textbf{[GRB (O.P. Tandon) p.90 Q3]}
\label{q:ext-grb_all_questions-422}

(d) 
ASSIGNMENT NO. 12 
'SEcnON/textasciitilde()1 
StraighfObjective TYpeQuestiOns 
This section contains 10' multiple choice questions. Each 
question has 4 choices (a), (b), (c) and (d), out of/textasciitilde()hich only 
one is correct.
\end{problembox}

\begin{problembox}
\textbf{Problem 423} \hfill \textbf{[GRB (O.P. Tandon) p.90 Q1]}
\label{q:ext-grb_all_questions-423}

The specific conductance (K) of an electrolyte of O.lN 
concentration is related to equivalent conductance (Ae ) by the 
following formula:' 
ICET (J and K) 2007]' 
, (a) Ae == K 
(b) Ae ;:::; 10K 
(c)Ae=100K 
(d)Ae=10000K
\end{problembox}

\begin{problembox}
\textbf{Problem 424} \hfill \textbf{[GRB (O.P. Tandon) p.90 Q2]}
\label{q:ext-grb_all_questions-424}

The standard E/textasciitilde()ed' values' of A, Band Care + 0.68 V, 
2.54 V, 
0.50 V respectively. The order of their reducing 
power is: 
(MHT-CET 2007) 
(a) A> B > C (b) A > C:> B(c) C > B > A (d) B > C > A
\end{problembox}

\begin{problembox}
\textbf{Problem 425} \hfill \textbf{[GRB (O.P. Tandon) p.90 Q3]}
\label{q:ext-grb_all_questions-425}

In the electrochemical reaction, 
+ Zn /textasciitilde() 
Zn 2+ + 2Fe2+, 
increasing the concentration ofFe2+: 
IJEE (W8) 2007) 
(a) increasing the cell emf 
(b) increasing the current flow 
(c) decrease the cell emf 
(d) alter the pH of the solution
\end{problembox}

\begin{problembox}
\textbf{Problem 426} \hfill \textbf{[GRB (O.P. Tandon) p.90 Q4]}
\label{q:ext-grb_all_questions-426}

Fully charged lead storage battery contains 1.5 L of 
5 M H2S04 , If 2.5 amp of current is taken from the cell for 
965 minutes, then what will be the molarity of remaining 
H2S04? Assume that volume of battery fluid to be constant: 
(a) 4 M 
(b) 3.5 M 
(c) 2M 
(d) 4.25 M
\end{problembox}

\begin{problembox}
\textbf{Problem 427} \hfill \textbf{[GRB (O.P. Tandon) p.90 Q5]}
\label{q:ext-grb_all_questions-427}

In a hydrogen-oxygen, 67.2 litre of H2 at STP is used in 
15 ..... ;",ltes. What is the average current produced? 
(a) 549.4 a'11p 
(b) 643.33 amp 
(c) 965 amp 
(d) 129.8 amp 
.' /textasciitilde() 
.
\end{problembox}

\begin{problembox}
\textbf{Problem 428} \hfill \textbf{[GRB (O.P. Tandon) p.90 Q6]}
\label{q:ext-grb_all_questions-428}

The cell reaction involving quiIihydr<;me electrode is: 
OH 
0 
    
+ 2lI' /textasciitilde():/textasciitilde()/textasciitilde()E/textasciitilde() = L30volt 
OR 
o 
What will be the electrode potenticilat pH: 3? , 
(a) 1.48 V 
fj1"J.20 V 
(c) 1.l0 V 
(d) 1.30 V 
llIint: 
E ",/Eo - 0.0591 10g [H+]2 
2 
=1.30+ 0.0591 x3 =1.48 V]
\end{problembox}

\begin{problembox}
\textbf{Problem 429} \hfill \textbf{[GRB (O.P. Tandon) p.90 Q7]}
\label{q:ext-grb_all_questions-429}

The standard reduction potential EO for ocr / Cr and f()r 
cr /1I2CI2 are 0.86 Nand - 1.10 volt respectively. The EO 
value of ocr /1/2Clz will be: 
(a) + 1.96 V 
(b) .,..1.96 V 
(c) + 0.24 V 
(d) 
0.24 V 
llIint: 
OC1- + H 20 + e- ---7 cr + 20W 
E/textasciitilde()d = 0.86 V 
Cl- ---7 !..C12 + e-
E';,xi = -1.10 V 
, 
2' 
=-0.24 V]
\end{problembox}

\begin{problembox}
\textbf{Problem 430} \hfill \textbf{[GRB (O.P. Tandon) p.90 Q8]}
\label{q:ext-grb_all_questions-430}

The standard reduction potential for the following two 
reactions ,are given: 
AgCI + e-
Ag(s) + a- (aq.); 
EO == 0.22 V 
... (i) 
Ag+ (aq.) + e- /textasciitilde() 
Ag(s) 
, 
EO == 0.80 V 
... (ii) 
The solubility prOduct of AgCI under standard condition will be: 
(a) 1.613 y. 10-5 M2 
(b) 1.535 x 10-8 M2 
(c) 3.213 X 10-10 M2 
(d) 1.535 x 10-10 M2
\end{problembox}

\begin{problembox}
\textbf{Problem 431} \hfill \textbf{[GRB (O.P. Tandon) p.91 Q9]}
\label{q:ext-grb_all_questions-431}

The equilibrium constant for the reaction, 
question has fQllowing 4 choices (a), (b), (c) and (d), out of 
. Cu(s) + Cu 2+ (aq.)/textasciitilde() 2Cu+ (aq.) 
which oniy one is correct.' 
. 
. 
(a) Statement-l is true; statement-2 is true; statement-2 is a 
E/textasciitilde()u2+ ICu = 0.34 V 
. E/textasciitilde()u2/textasciitilde() ICu = 0.15 V 
correct explanation for statement-I. 
(Given: log 3.72:::; 0.571) 
(b) Statement-1 is true; statement-2 is true; statement-2 is not 
(a) 3.72 x 10-6 
.. 
(b) 3.72 x 10-5 
. 
Ii correct explanation for statement-I. 
(c)3.72 x 10-7 
(d) 3.72 x 10-8 
(c) Statement-1 is true; statement-2 is false. 
. 
. .' 022 
://1,' 
(d) Statement-I is false; statement-2 is true.
\end{problembox}

\begin{problembox}
\textbf{Problem 432} \hfill \textbf{[GRB (O.P. Tandon) p.91 Q10]}
\label{q:ext-grb_all_questions-432}

Cations absorb 6.023 x 1 
electrons for the/textasciitilde()dtlction. How
\end{problembox}

\begin{problembox}
\textbf{Problem 433} \hfill \textbf{[GRB (O.P. Tandon) p.91 Q16]}
\label{q:ext-grb_all_questions-433}

StatemenM .. During the electrolysis of water, two faraffiry 6f 
SECTION/textasciitilde()II 
Multiple Answers TYpe Objective Questions
\end{problembox}

\begin{problembox}
\textbf{Problem 434} \hfill \textbf{[GRB (O.P. Tandon) p.91 Q11]}
\label{q:ext-grb_all_questions-434}

. In an erectro[Ytlc/textasciitilde()cerf: 
(a) anode is positively charged 
(b) cathode is negatively charged 
(c) oxidation takes place at anode . 
(d) reduction takes place at cathode
\end{problembox}

\begin{problembox}
\textbf{Problem 435} \hfill \textbf{[GRB (O.P. Tandon) p.91 Q12]}
\label{q:ext-grb_all_questions-435}

One gram equivalent of a substance is liberated at an electrode 
by: 
(a) 6.023 x IOz3 electrons 
(b) 96500 C$/cdot$ 
(c) I amp current for I second 
(d) I amp current for 96500 sec
\end{problembox}

\begin{problembox}
\textbf{Problem 436} \hfill \textbf{[GRB (O.P. Tandon) p.91 Q13]}
\label{q:ext-grb_all_questions-436}

If 9 gm. H20 is electrolysed completely with the current of 
50/% efficiency then: 
(a) 96500 charge is required 
(b) 2 x 96500 C charge is required 
(c) 5.6 L ofOz at STP will be formed 
(d) 11.2 L of 0z at STP will be formed
\end{problembox}

\begin{problembox}
\textbf{Problem 437} \hfill \textbf{[GRB (O.P. Tandon) p.91 Q14]}
\label{q:ext-grb_all_questions-437}

A galvanic cell involves the following reaction: 
Zn(s) + 2Ag+(aq.)/textasciitilde() Znz+ (aq.)+ 2Ag(s) 
Select the correctstatetnents among the following: 
(a) Zinc is negatively charged 
(b) The given redox process is spontaneous 
(c) Ag + + e -
Ag, takes place at anode 
(d) Zn(s) -----t Zn 2+ + 2e-, takes place at cathode
\end{problembox}

\begin{problembox}
\textbf{Problem 438} \hfill \textbf{[GRB (O.P. Tandon) p.91 Q15]}
\label{q:ext-grb_all_questions-438}

Given that, 
E/textasciitilde()i2+ INi = - 0.25 V, E/textasciitilde()u2+ ICu ::: + 0.34 V 
E:g+ lAg 
+ 0.80 V E;n2+ Izn' 
0.76 V 
Which of the following redox processes will not take place in 
specified direction? 
(a) Ni2+ (aq.) + Cu(s) -----t Ni(s) + Cu 2+ (aq.) 
(b) Cu(s) + 2Ag+ (aq.) -----t Cu2+ (aq.) + 2Ag(s) 
Because 
Statement-2: In the electrolysis of water, two faraday of 
charge will produce half mole of Hz gas and one fourth mole 
ofOz gas. 
17.-Statement:'l: Aqueous solution of'CuSO;;-wrnscoloutleslnrn---
complete electroiysis using platinum electrode. 
Because 
Statement-2: CUS04 is converted to Cu(OHh on electrolysis. 
18; Statement-I:' Sodium ions are discharged at a mercury 
cathode in preference to hydrogen ion. 
Because 
Statement-2:. Na + is stronger reducing agent than W.
\end{problembox}

\begin{problembox}
\textbf{Problem 439} \hfill \textbf{[GRB (O.P. Tandon) p.91 Q19]}
\label{q:ext-grb_all_questions-439}

Statement-I: KCl and NH4Cl cannot be used in salt bridge of 
a cell containing Ag + , Hg/textasciitilde()+ and Tl+ ions. 
$/cdot$Because 
Statement-2: Cell will be destroyed due to precipitation of 
metal chlorides.
\end{problembox}

\begin{problembox}
\textbf{Problem 440} \hfill \textbf{[GRB (O.P. Tandon) p.91 Q20]}
\label{q:ext-grb_all_questions-440}

Statement-I: The voltage of mercury cell remains constant 
for its life time. 
. 
Because 
Statement-2: Overall cell reaction does not involve any ion.
\end{problembox}

\begin{problembox}
\textbf{Problem 441} \hfill \textbf{[GRB (O.P. Tandon) p.91 Q21]}
\label{q:ext-grb_all_questions-441}

Statement-I: In alkaline version of dry cell, /textasciitilde()Cl is 
replaced by KOH. 
Because 
Statement-2: Zinc container does not undergo corrosion ill 
alkaline medium. 
$/cdot$$/cdot$$/cdot$SECTION-IV 
Matrix-Matching Type Questions 
, 
This section contains 3 questions;. Each question contains 
statement given in two columns which have to be matched. 
Statements (a, b, c and d) in Coiumn-I have to be matched with 
statements (p, q, r and s) in Column-II. The answers to these 
questions have to be appropriately bubbled as illustrated in the 
following examples: 
If the correct matches are (a-p,s); (b-q,r); (c-P,q) and (d-s); 
then correct bubbled 4 x 4 matrix should be as follows:
\end{problembox}

\begin{problembox}
\textbf{Problem 442} \hfill \textbf{[GRB (O.P. Tandon) p.92 Q22]}
\label{q:ext-grb_all_questions-442}

Match the Column-I with Column-II: 
Column-I 
Column-II 
(a) Nickel-Cadmium cell 
(p) Used in auto vehicles 
(b) Lithium battery 
(q) Secondary cell
\end{problembox}

\begin{problembox}
\textbf{Problem 443} \hfill \textbf{[GRB (O.P. Tandon) p.92 Q23]}
\label{q:ext-grb_all_questions-443}

Match the Column-I with Col limn-II: 
Column-I 
Column-TI 
(a) Specific conductance, K 
(p) A/textasciitilde(), I A/textasciitilde(), 
(b) Molar conductance, Am 
(q) Decreases with diluti:on 
(c) Resistance of electrolyte (r) Increases with dilution 
solution, R 
(d) Degree of ionization of 
(s) Increases with increase in 
weak electrolyte, a 
the distance between parallel 
plates
\end{problembox}

\begin{problembox}
\textbf{Problem 444} \hfill \textbf{[GRB (O.P. Tandon) p.92 Q24]}
\label{q:ext-grb_all_questions-444}

Match the Column-I with Column-II: 
Column-I 
(a) Concentration cell 
(b) Edison cell 
(c) Mercury cell 
(d) Dry cell
\end{problembox}

\begin{problembox}
\textbf{Problem 445} \hfill \textbf{[GRB (O.P. Tandon) p.92 Q23]}
\label{q:ext-grb_all_questions-445}

(a-q) (b-r) (c-q,s) (d-p,r) 
Column-II 
(P) Fe is oxidised by Ni20 3 
(q) Zinc anode 
(r) HgO cathode 
(s) EO 
0
\end{problembox}

\begin{problembox}
\textbf{Problem 446} \hfill \textbf{[GRB (O.P. Tandon) p.92 Q24]}
\label{q:ext-grb_all_questions-446}

(a-s) (b-p) (c-q,r) (d-q) 
SEcnON ... V 
Linked Comprehension Type Questions 
Redox reactions playa pivotal role in chemistry and biology. 
The values of standard redox potential (EO) of two half-cells 
reactions decide which way the reaction is expected to 
proceed. A simple example is a Daniell cell in which zinc goes 
into solution and copper gets deposited. Given below are a set 
of half-cell reactions (acidic medium) along with their EO (V 
with respect to normal hydrogen electrode) values. 
12 + 2e- /textasciitilde() 
2r 
EO 
0.54 
Cl2 + 2e- /textasciitilde() 
2cr 
EO 
1.36 
Mn3++e-/textasciitilde()Mn2+ 
EO 
1.50 
. Fe3+ + e- /textasciitilde()Fe2+ 
EO 
0.77 
02 + 4W + 4e- /textasciitilde() 
2H20 
E () 
1.23 
(lIT 2007) 
Answer the following questions: 
(b) Fe2+ is oxidised by iodine 
. (c) Iodide ion is oxidised by chlorine 
(d) Mn 2+ is oxidised by chlorine 
[Hint: Species with greater reduction potential, oxidises other 
with lower reduction potential.]
\end{problembox}

\begin{problembox}
\textbf{Problem 447} \hfill \textbf{[GRB (O.P. Tandon) p.92 Q26]}
\label{q:ext-grb_all_questions-447}

While Fe3+ is stable, Mn 3+ is not stable in acid solution 
because: 
(a) 02 oxidises Mn 1+ to Mn 3+ 
(b) O2 oxidises both Mn2+ and Fe2+ 
(c) Fe3+ oxidises H20 to O2 
(d) Mn 3+ oxidises H20to O2 
[Hint: Mn3+ oxidises H20 to 02 because the standard reduction 
potential of (Mn 3+ /textasciitilde() 
Mn 2+) is greater than that of 
(02'/textasciitilde() H 20).]
\end{problembox}

\begin{problembox}
\textbf{Problem 448} \hfill \textbf{[GRB (O.P. Tandon) p.92 Q22]}
\label{q:ext-grb_all_questions-448}

(a-q) (b-q) (c-r,s) (d-p,q)
\end{problembox}